\documentclass[11pt,oneside]{book}
\usepackage[
    paperwidth=6in,paperheight=9in,
    top=0.75in,bottom=0.75in,left=0.6in,right=0.6in,twoside=false
]{geometry}
\input{../pri-end/T0_preamble_De.tex}
% T0_preamble_patches.tex — LOKALER PLATZHALTER
% Im Repo durch die offizielle Version ersetzen.

\usepackage{graphicx}
\makeatletter\renewcommand{\thechapter}{\arabic{chapter}}\makeatother
\raggedbottom\allowdisplaybreaks
\hfuzz=5pt\vfuzz=5pt\emergencystretch=2em
\newcommand{\status}[1]{\penalty10000\,{\scriptsize\textcolor{gray!60}{[#1]}}}
\newcommand{\quelle}[1]{\par\vspace{1pt}\noindent{\footnotesize\color{gray!60}\itshape Quellen: #1}\par\vspace{6pt}}
\newenvironment{laierspur}{\par\vspace{4pt}\noindent\small\itshape\color{blue!60!black}\textbf{Für alle Leser:}\enspace\ignorespaces}{\par\vspace{4pt}\normalsize\normalfont\normalcolor}
\setcounter{tocdepth}{1}
\begin{document}
\frontmatter
\pagestyle{fancy}
\makeatletter\let\ps@plain\ps@fancy\makeatother
\begin{titlepage}
  \centering\vspace*{1.5cm}
  {\Huge\bfseries\color{blue!70!black}Quantenmechanik ist deterministisch}\\[0.6cm]
  {\Large\bfseries\color{blue!40!black}Die geometrische Lesart der Quantenmechanik}\\[1cm]
  {\LARGE\itshape Fundamentale Fraktale Geometrische\\[0.2cm]Feldtheorie (FFGFT)}\\[2cm]
  {\large Johann Pascher}\\[0.3cm]
  {\normalsize ORCID: 0009-0000-6518-4064}\\[0.3cm]
  {\normalsize\url{https://github.com/jpascher/T0-Time-Mass-Duality}}\\[1.5cm]
  {\normalsize September 2026}\vfill
  {\small Version 1.0}
\end{titlepage}
\newpage\thispagestyle{empty}%
{\small\noindent\textbf{Stichworte:} Bellsche Ungleichung $\cdot$ Verschränkung $\cdot$ Hilbertraum $\cdot$ Qubit $\cdot$ Quantencomputer $\cdot$ $T^4/\mathbb{Z}_3$ $\cdot$ $\tilde T\cdot m=1$\\[0.8em]
\noindent Lizenz: Creative Commons BY 4.0}\par\vspace{1em}
{\small\noindent\textbf{Hinweis:} Dieses Buch stellt den Stand des FFGFT-Korpus (September 2026) dar. Die zugrundeliegenden Dokumente (Dok.~020--334) sind im Repositorium frei zugänglich. Statusmarker: [K] konfirmiert $\cdot$ [B] bewiesen $\cdot$ [S] skizziert $\cdot$ [E] etablierte Physik $\cdot$ [X] geprüft negativ.}
\tableofcontents
\mainmatter

% --- Teil I ---
\part{Grundlagen}
% QMB_n01_grundlage.tex — Kapitel 1: Die geometrische Grundlage
% Quellen: Dok. 230, 307, 314, 365
\chapter{Die geometrische Grundlage}
\label{ch:grundlage}

Die Fundamentalrelation der FFGFT lautet
\begin{equation}
  \tilde T \cdot m = 1,
  \label{eq:tm}
\end{equation}
wobei $\tilde T$ eine kompakte Zeitkoordinate und $m$ die Masse einer
Mode ist.\status{SETZUNG} Diese Relation ist keine abgeleitete
Gleichung, sondern die geometrische Setzung, aus der alles Weitere
folgt.

\section{Der Träger: $T^4/\mathbb{Z}_3$}

Die FFGFT rechnet auf einem kompakten Torus $T^4 = (S^1)^4$ mit einer
$\mathbb{Z}_3$-Gittersymmetrie. Dieser Träger hat zwei wesentliche
Eigenschaften.

\textbf{Kompaktheit.} Jede Koordinate ist periodisch. Das hat eine
unmittelbare Konsequenz: Die zulässigen Moden sind diskret — nicht weil
die Natur diskret postuliert wird, sondern weil Periodizität
diskrete Fouriermoden erzwingt. Die Fouriermoden des Trägers liefern
genau die Teilchenspektren des Standardmodells.\status{K}

\textbf{$\mathbb{Z}_3$-Symmetrie.} Die Quotientierung durch
$\mathbb{Z}_3$ erzeugt neun Fixpunkte und eine Trialitätssymmetrie mit
$|\mathrm{Aut}(D_4)| = 1152 = 2^7\cdot3^2$. Diese Symmetrie wird im
Spektraloperator sichtbar und steuert, welche Moden stabil sind.\status{B}

\section{Der Strukturparameter $\xi$}

Der einzige freie Parameter der Theorie ist
\begin{equation}
  \xi = \frac{4}{30000} = \frac{1}{7500} \approx 1{,}333\times10^{-4}.
\end{equation}
Er ist keine Kalibrierungsgröße, sondern folgt aus den
Galois-Gruppenordnungen des Trägers: Speziell gilt
$\xi = C_2(SU(3)_{\mathrm{fund}})/N_{\mathrm{Fourier}}= (4/3)/10^4$.\status{K}
(Dok.~338, Dok.~353)

Alle Massen, Kopplungen und Korrekturen der FFGFT lassen sich aus $\xi$
ableiten; SI-Ankerwerte (Elektronenmasse, $H_0$) dienen nur der
Umrechnung in SI-Einheiten, nicht der Ableitung.

\section{Die Relation $\tilde T\cdot m = 1$ und ihre Konsequenzen}

Gleichung~\eqref{eq:tm} hat drei unmittelbare Folgen.

\textbf{Modengebundene Zeitskala.} Jede Mode mit Masse $m_i$ trägt
eine eigene kompakte Zeitskala $\tilde T_i = 1/m_i$.\status{B} Die Zeit
ist kein äußerer Parameter, sondern eine modespezifische Eigenschaft des
Trägers. Pauli-Theorem und der fehlende Zeitoperator der Standard-QM
sind in dieser Lesart kein Problem: Auf einem kompakten Träger gilt das
Pauli-Argument nicht (Dok.~334).

\textbf{Massenarithmetik aus Geometrie.} Das Verhältnis
$(m_\mu/m_e)^2 = 43200$ folgt sowohl geometrisch aus den
Wicklungszahlen als auch aus den Galois-Gruppenordnungen
$|GF(9)^*|^2\cdot 5^2\cdot|GF(27)|$.\status{K} (Dok.~338)

\textbf{Fraktale Korrektur.} Der Übergang vom kompakten Träger zur
SI-Skala bringt eine rekursive Korrektur
$K_{\mathrm{frak}} = 1 - 100\xi \approx 0{,}9867$, die aus der
Typ-III-Pullback-Struktur des Trägers folgt.\status{B} (Dok.~333)

\begin{laierspur}
Stell dir vor, Physik passiert nicht auf einer unendlich langen
geraden Linie, sondern auf einer Trommel — einer Oberfläche, die
in sich geschlossen ist. Auf einer Trommel sind nur bestimmte
Schwingungen möglich: die, die nach einem Umlauf wieder dort
ankommen, wo sie angefangen haben. Genau so funktioniert $T^4$:
Nur Moden, die auf dem Torus „passen", existieren — und das sind
exakt die Teilchen, die wir messen.
\end{laierspur}
\quelle{Dok.~230, 307, 314, 338, 353, 365}

% QMB_n02_zustand_messung.tex — Kapitel 2: Zustände, Messung, Zeit
% Quellen: Dok. 230, 307, 322, 334
\chapter{Zustände, Messung und Zeit}
\label{ch:zustand}

\section{Zustände als Modenkonfigurationen}

In der Standard-QM ist ein Zustand ein Vektor $|\psi\rangle$ in einem
Hilbertraum $\mathcal{H}$. In der FFGFT ist ein Zustand eine
Modenkonfiguration $(z, r, \theta)$ auf dem Träger $T^4/\mathbb{Z}_3$.
Die Verbindung ist bijektiv: Der Spektraloperator
\begin{equation}
  \hat F = \sum_k f_k\, |k\rangle\langle k|,
  \label{eq:spektral}
\end{equation}
mit Moden $k$ aus dem Gitter $D_4$ ist selbstadjungiert auf
$\mathcal{H} = L^2(T^4)\otimes\mathbb{C}^3$.\status{B} (Dok.~327)
Jede QM-Rechnung ist in FFGFT-Sprache abbildbar und umgekehrt;
die Formalismen liefern dieselben statistischen Vorhersagen.\status{B}

\subsection{Die Übersetzungstabelle}

\begin{center}
\small\renewcommand{\arraystretch}{1.25}
\begin{tabular}{p{3.4cm}p{4.2cm}}
\toprule
Hilbertraum-QM & FFGFT \\
\midrule
Qubit $|\psi\rangle = \alpha|0\rangle+\beta|1\rangle$ & Mode $(z,r,\theta)\in S^1\times\mathbb{R}_{>0}\times S^1$ \\
Bloch-Sphäre $S^2$ & Zylindrische Faser $(r,\theta)$ \\
Superposition & Intermediäre Torsion einer Mode \\
Tensorprodukt $|\psi\rangle\otimes|\phi\rangle$ & Unabhängige Modenkonfigurationen \\
Verschränkung & Topologische Verbindung (nicht reduzierbar) \\
Unitäre Evolution & Phasendrehung auf $T^4$ \\
Messung / Kollaps & Polübergang zu $z=\pm1$ \\
Born-Regel $P=|\alpha|^2$ & Statistisches Mittel über unbekannte Phase $\theta$ \\
\bottomrule
\end{tabular}
\end{center}

\section{Messung ohne Kollaps-Postulat}

Der Kollaps der Wellenfunktion ist in der Standard-QM ein Postulat
ohne mechanische Erklärung. In der FFGFT ist Messung eine
\emph{geometrische Wechselwirkung}: Das Messgerät treibt die Mode
auf das nächstgelegene stabile Extrem ($z=+1$ oder $z=-1$). Das ist
eine deterministische Bewegung im Modenraum — aber wegen der
unbekannten Anfangsphase $\theta$ statistisch ununterscheidbar
von Born-Regel-Verhalten:\status{B}
\begin{equation}
  P_{\mathrm{FFGFT}}(z=+1) = \frac{1+z}{2} = |\alpha|^2 = P_{\mathrm{QM}}(0).
\end{equation}
Die beiden Formalismen liefern identische statistische Vorhersagen
(Dok.~230). Interpretativ unterscheiden sie sich: Was die QM als Zufall
beschreibt, ist in der FFGFT Unkenntnis der genauen Trägerkonfiguration.

\section{Die Abwesenheit des Zeitoperators}

Paulis Einwand zeigt, dass es zu einem nach unten beschränkten Hamilton-
operator keinen selbstadjungierten, kanonisch konjugierten Zeitoperator
geben kann.\status{E} Die Zeit fehlt im Hilbertraum — nicht als Versehen,
sondern strukturell.

In der FFGFT ist das kein Mangel: $\tilde T$ ist kein Operator, sondern
eine modengebundene Eigenschaft des Trägers. Für Mode $i$ gilt
$\tilde T_i = 1/m_i$; verschiedene Moden tragen verschiedene Zeitskalen.
Superposition von Moden verschiedener Masse ergibt damit
\emph{inkompatible Zeitskalen} — eine strukturelle Aussage über die
Grenzen modenübergreifender Superposition.\status{B} (Dok.~334)

Die Schrödinger-Gleichung bleibt als effektive Beschreibung vollständig
gültig; die Zeit $t$ darin ist die globale Koordinate, die aus dem
Träger ausgerollt wird (Typ-I-Pullback, Dok.~333).

\begin{laierspur}
Dass die QM keinen Zeitoperator hat, klingt technisch — bedeutet aber:
Die Theorie kann über die Dauer von Prozessen nichts sagen, nur über
Zustände zu festen Zeitpunkten. Die FFGFT macht das explizit: Jedes
Teilchen hat eine eigene innere Uhr ($\tilde T = 1/m$). Zwei Teilchen
verschiedener Masse schwingen buchstäblich in verschiedenen Rhythmen.
\end{laierspur}
\quelle{Dok.~230, 307, 322, 327, 333, 334}

% QMB_n03_verschraenkung.tex — Kapitel 3: Verschränkung und scheinbare Instantanität
% Inhaltliche Grundlage: Dok. 131, 230, 334, 364; eigene Darstellung
\chapter{Verschränkung ohne Fernwirkung und scheinbare Instantanität}
\label{ch:verschraenkung}

\section{Was Verschränkung in der FFGFT ist}

Zwei Qubits heißen verschränkt, wenn ihr gemeinsamer Zustand nicht als
Produkt $|\psi_1\rangle\otimes|\psi_2\rangle$ geschrieben werden kann.
Der maximal verschränkte Bell-Zustand
\begin{equation}
  |\Phi^+\rangle = \tfrac{1}{\sqrt2}(|00\rangle+|11\rangle)
\end{equation}
hat in der FFGFT eine geometrische Charakterisierung: Die beiden Moden
teilen einen gemeinsamen Träger $T^4/\mathbb{Z}_3$. Ihre
$(z_i,r_i,\theta_i)$-Koordinaten sind nicht unabhängig --- sie sind
Knoten derselben topologischen Konfiguration, die nicht auf Produktform
reduziert werden kann.\status{B}

\textbf{Verschränkung ist in der FFGFT eine topologische Verbindung,
keine nichtlokale Korrelation.} Die Teilchen hatten nie zwei getrennte
Träger; sie teilen von Anfang an denselben. Was die QM als Fernwirkung
beschreibt, ist eine Eigenschaft der gemeinsamen Geometrie.

\section{EPR und Lokalität}

Das EPR-Paradoxon (1935): Wenn QM vollständig ist, sind verschränkte
Teilchen nichtlokal verbunden; Einstein hielt das für inakzeptabel und
vermutete verborgene Variablen.\status{E}

Die FFGFT löst das Paradoxon geometrisch. Lokalität gilt auf dem Torus
in dem Sinn, dass kein Signal schneller als Licht übertragen wird. Die
Korrelation ist eine Eigenschaft der gemeinsamen Geometrie, keine
Signalübertragung. Was aufgegeben wird, ist nur die Vorstellung, dass
jedes Teilchen einen \emph{unabhängigen} Träger hätte --- und die ist
nicht beobachtbar.\status{B}

\section{Scheinbare Instantanität: kein physikalisches Problem}

Der Kollaps der Wellenfunktion scheint instantan zu geschehen: Man misst
Qubit A und B weiß sofort, in welchem Zustand es ist, egal wie weit
entfernt. Das scheint der Relativitätstheorie zu widersprechen.

Bevor man die FFGFT-Auflösung erklärt, lohnt ein Schritt zurück. In
der Standardmathematik gelten Rechenoperationen als instantan: Wenn
$2 + 2 = 4$, dann gilt das gleichzeitig für alle Stellen des Ausdrucks
--- niemand verknüpft damit eine physikalische Instantanität, die
kausal problematisch wäre. Die Gleichheit ist eine strukturelle
Aussage über das Objekt, kein Ereignis im Raum. Genauso verhält es
sich mit der Zustandsbeschreibung in der QM: Die Wellenfunktion
$|\Phi^+\rangle = \frac{1}{\sqrt2}(|00\rangle+|11\rangle)$ ist eine
mathematische Beschreibung eines Zustands --- kein Prozess, der sich
irgendwo ausbreitet. Die Instantanität der Beschreibungsanpassung nach
der Messung ist mathematischer, nicht physikalischer Natur.

In der FFGFT gilt dasselbe Argument auf der Ebene des Trägers: Die
Relation $\tilde T\cdot m = 1$ ist eine \textbf{strukturelle
Eigenschaft} jeder Mode auf $T^4/\mathbb{Z}_3$ --- analog dazu, dass
eine Gleichung für alle ihre Variablen gleichzeitig gilt. Sie ist eine
\textbf{lokale Zwangsbedingung} am selben Raumpunkt, kein
Signal zwischen entfernten Punkten.

Die Auflösung kommt aus der Trennung zweier Ebenen: Eine Änderung bei $\mathbf{x}_1$ breitet sich dann über
die dynamischen Feldgleichungen aus:\status{B}
\begin{equation}
  \text{Änderung bei }\mathbf{x}_1
  \to \text{Ausbreitung mit }v\le c
  \to \text{Wirkung bei }\mathbf{x}_2.
\end{equation}
Die Zeitverzögerung ist $\Delta t = |\mathbf{x}_2-\mathbf{x}_1|/c$.

Lokale Anpassung und globale Ausbreitung sind zwei verschiedene Prozesse.
Die lokale Anpassung ($T\cdot m = 1$ am selben Punkt) erfolgt auf der
Planck-Zeitskala $t_P\sim10^{-43}$\,s --- praktisch instantan für jede
messbare Zeitskala. Die Ausbreitung zu entfernten Punkten ist durch die
Lichtgeschwindigkeit begrenzt. Diese Hierarchie erklärt, warum viele
Quantenprozesse scheinbar instantan wirken.

\section{Die Feldgleichung und ihre Kausalstruktur}

Die FFGFT-Felddynamik folgt einer Wellengleichung:
\begin{equation}
  \frac{\partial^2\tilde T}{\partial t^2} = c^2\nabla^2\tilde T + Q(\tilde T,m,\rho),
\end{equation}
mit der Quellfunktion $Q = -4\pi G\rho\cdot\tilde T$. Das ist eine
hyperbolische Differentialgleichung, charakteristisch für
Wellenausbreitung mit endlicher Geschwindigkeit.

Die Kausalstruktur ist durch die Greensche Funktion manifestiert:
\begin{equation}
  G(\mathbf{x},\mathbf{x}',t-t') = \theta(t-t')\cdot
  \frac{\delta(|\mathbf{x}-\mathbf{x}'|-c(t-t'))}{4\pi|\mathbf{x}-\mathbf{x}'|}.
\end{equation}
Die Heaviside-Funktion $\theta(t-t')$ garantiert Kausalität: Keine
Wirkung vor ihrer Ursache. Die Delta-Funktion kodiert Ausbreitung mit
Lichtgeschwindigkeit. Diese Struktur ist vollständig mit der speziellen
Relativitätstheorie kompatibel.\status{B}

Ein konkretes Beispiel: Ändert sich $m(\mathbf{x}_0)$ bei $t = 0$, dann
passt sich $\tilde T(\mathbf{x}_0)$ sofort an (lokale Zwangsbedingung);
gleichzeitig entsteht $\nabla^2\tilde T\ne0$, das eine Feldstörung
erzeugt; diese breitet sich als Welle mit $v = c$ aus und erreicht
$\mathbf{x}_1$ zur Zeit $t = r/c$.

\section{Die Hierarchie der Zeitskalen}

Die scheinbare Instantanität resultiert aus der extremen Trennung
verschiedener Zeitskalen:

\begin{center}
\small\renewcommand{\arraystretch}{1.3}
\begin{tabular}{p{3.5cm}p{2.5cm}p{3.2cm}}
\toprule
\textbf{Prozess} & \textbf{Zeitskala} & \textbf{Charakter} \\
\midrule
Lokale T-E-Anpassung & $t_P\sim10^{-43}$\,s & unmessbar schnell \\
Wellenfunktionsevolution & $\sim10^{-15}$\,s & Quantenbereich \\
Kausale Feldausbreitung & $t = r/c$ & relativistisch \\
Makroskopische Messung & $\gg10^{-9}$\,s & klassisch \\
\bottomrule
\end{tabular}
\end{center}

Prozesse auf der Planck-Skala sind so schnell, dass sie mit keiner
vorstellbaren Technologie zeitlich aufgelöst werden könnten. Was als
instantan erscheint, läuft auf dieser Skala --- und ist damit kausal
vollständig verträglich.

\section{Die vollständige Dualität: Zeit, Masse, Energie und Länge}

Die FFGFT beschreibt ein umfassendes System von Dualitäten, in dem alle
fundamentalen Größen durch $\tilde T\cdot m = 1$ miteinander verbunden sind:
\begin{align}
  \tilde T\cdot E &= 1\qquad\text{(Zeit-Energie-Dualität)} \\
  \tilde T &= 1/m\qquad\text{(Zeit-Masse-Beziehung)} \\
  E &= mc^2\qquad\text{(Masse-Energie-Äquivalenz)} \\
  \ell &= \hbar/(mc)\qquad\text{(Compton-Wellenlänge als Länge-Energie-Brücke)}
\end{align}
An der Planck-Skala konvergieren alle diese Dualitäten; dort sind Zeit,
Länge, Masse und Energie nicht mehr unabhängig:
\begin{equation}
  t_P = \ell_P/c = \sqrt{\hbar G/c^5}\approx5{,}4\times10^{-44}\,\text{s}.
\end{equation}
Die Dualitäten sind nicht nur mathematische Formulierungen, sondern
haben physikalische Konsequenz: Sie zeigen, dass die scheinbar
verschiedenen Konzepte von Zeit, Raum, Masse und Energie verschiedene
Aspekte derselben geometrischen Realität sind. Je nach Skala dominiert
ein anderer Aspekt: Bei atomaren Skalen die Zeit-Energie-Dualität, bei
makroskopischen Skalen die Masse, bei kosmologischen die Länge-Zeit-Relation.

\section{Topologische Verbindung: der heutige Stand}

Die frühen FFGFT-Dokumente (Herbst 2025) versuchten, die Bell-Korrelationen
durch eine $\xi$-Dämpfung zu erklären, die den CHSH-Wert leicht unter
die Tsirelson-Schranke drückt. Dieser Ansatz ist durch Dok.~230 (Mai 2026)
überholt. Die heutige Position:

Verschränkte Teilchen teilen einen Träger $T^4/\mathbb{Z}_3$ ohne
Fernwirkung. Die statistischen Vorhersagen sind identisch mit der
Quantenmechanik.\status{B} Es gibt keine Fernwirkung, weil es keinen
getrennten Träger gibt. Die topologische Verbindung ist keine
Wirkung über Distanz, sondern eine geometrische Tatsache des gemeinsamen
Trägers. Eine kleine $\xi$-Abweichung $\Delta\mathrm{CHSH}\approx
\xi/(2\pi)\approx10^{-5}$ pro Messung ist als Arbeitshypothese offen,
aber nicht mehr das Kernargument.\status{S}

\begin{laierspur}
Zwei Punkte auf demselben Torus sind verbunden --- nicht durch ein
Signal zwischen ihnen, sondern weil sie dieselbe Oberfläche teilen. Das
ist der Unterschied: EPR fragte nach einer Verbindung zwischen zwei
getrennten Objekten. Die FFGFT sagt: Die Prämisse ist falsch. Die
Objekte waren nie getrennt.
\end{laierspur}

\quelle{Dok.~131 (Scheinbare Instantanität, 2026), Dok.~230 (Verschränkung als topologische Verbindung, Mai 2026), Dok.~334, Dok.~364}


% --- Teil II ---
\part{Bellsche Ungleichung}
% QMB_n04_bell_theorem.tex — Kapitel 4: Bells Theorem
% Quellen: Dok. 023, 023a, 074
\chapter{Bells Theorem}
\label{ch:bell}

\section{Die Grundfrage}

Bell-Tests prüfen, ob die Korrelationen verschränkter Teilchen durch
eine lokal-realistische Theorie erklärt werden können — eine Theorie,
in der jedes Teilchen von Anfang an definierte Eigenschaften besitzt,
die unabhängig vom anderen Teilchen gemessen werden.

John Bell zeigte 1964, dass keine solche Theorie die
Quantenvorhersagen reproduzieren kann.\status{E} Das Argument ist
mathematisch scharf und lässt sich nicht durch technische Verfeinerungen
umgehen.

\section{Die CHSH-Ungleichung}

In der von Clauser, Horne, Shimony und Holt formulierten Form
(CHSH, 1969) lautet die Bell-Ungleichung:\status{E}
\begin{equation}
  |E(a,b)-E(a,c)| + |E(a',b)+E(a',c)| \le 2,
  \label{eq:chsh}
\end{equation}
wobei $E(a,b)$ die Korrelation der Messergebnisse bei den
Detektorwinkeln $a$ und $b$ bezeichnet.

Die Quantenmechanik verletzt diese Schranke bis zur
Tsirelson-Schranke $2\sqrt2\approx2{,}828$.\status{E} Alle
schlupflochfreien Experimente seit 2015 bestätigen die
QM-Vorhersage und widerlegen lokal-realistische Theorien.\status{E}

\section{Die drei Annahmen}

Bells Beweis setzt drei Annahmen voraus:
\begin{enumerate}
  \item \textbf{Lokalität:} Das Ergebnis bei Detektor A hängt nicht
        von der Einstellung bei Detektor B ab.
  \item \textbf{Realismus:} Die Messergebnisse sind vorher bestimmt
        (durch verborgene Variablen).
  \item \textbf{Messfreiheit:} Die Detektoreinstellungen werden
        unabhängig von den verborgenen Variablen gewählt.
\end{enumerate}
Mindestens eine dieser Annahmen ist falsch — das ist Bells
Leistung. Die QM gibt Lokalität und Realismus gemeinsam auf.

\section{Historische FFGFT-Ansätze und ihr heutiger Stand}

Die frühen FFGFT-Dokumente (Dok.~023, 023a, 074, Herbst 2025)
versuchten verschiedene Wege, die Verletzung mit lokaler Realität
zu versöhnen:

\begin{itemize}
  \item \textbf{Zeitfeld-Dämpfung} (Dok.~023): Die Korrelationen
        werden durch einen $\xi$-abhängigen Faktor abgeschwächt.
        Ergebnis: CHSH-Reduktion von $2{,}828$ auf $2{,}827$ —
        zu klein für heutige Experimente.\status{S}
  \item \textbf{Superdeterminismus} (Dok.~074): Die Messfreiheit
        ist verletzt, weil das Energiefeld Detektoreinstellungen
        und Teilcheneigenschaften gemeinsam bestimmt. Dieser Weg
        ist durch Dok.~230 überholt (s. Kapitel~\ref{ch:bell_ffgft}).
  \item \textbf{Topologische Verbindung} (Dok.~230, Mai 2026):
        Kein Zugeständnis an die Annahmen nötig — der gemeinsame
        Träger macht die Korrelation zu einer geometrischen
        Eigenschaft. Dieser Weg ist heute maßgeblich.\status{B}
\end{itemize}

\quelle{Dok.~023, 023a, 074, 190 (Präz.~11, 12), 230}

% QMB_n05_bell_ffgft.tex — Kapitel 5: FFGFT-Antwort auf Bell
% Quellen: Dok. 230, 190, 353, 364
\chapter{Die FFGFT-Antwort auf Bell}
\label{ch:bell_ffgft}

\section{Topologische Verbindung statt Fernwirkung}

Die heutige FFGFT-Position (Dok.~230, Mai 2026) greift keine der drei
Bell-Annahmen an. Sie verändert stattdessen den ontologischen Rahmen:
Verschränkte Teilchen haben keinen getrennten Träger, sondern teilen
denselben Torus $T^4/\mathbb{Z}_3$. Die Korrelation ist eine
topologische Eigenschaft dieses gemeinsamen Trägers.

Das bedeutet konkret:
\begin{itemize}
  \item \textbf{Lokalität} gilt im Sinne des Trägers: Kein Signal
        überträgt Information schneller als Licht.
  \item \textbf{Realismus} gilt in dem Sinn, dass die Moden-
        konfiguration vor der Messung existiert.
  \item \textbf{Messfreiheit} wird nicht eingeschränkt.
\end{itemize}
Was aufgegeben wird, ist die Vorstellung, dass jedes Teilchen
einen \emph{unabhängigen} Träger hätte. Die topologische Verbindung
ist keine Fernwirkung — sie ist eine geometrische Tatsache.\status{B}

\section{Statistische Vorhersagen}

Die FFGFT reproduziert die Born-Regel und damit alle statistischen
Vorhersagen der QM vollständig:\status{B}
\begin{equation}
  P_{\mathrm{FFGFT}}(z=+1) = \tfrac{1+z}{2} = |\alpha|^2 = P_{\mathrm{QM}}(0).
\end{equation}
Insbesondere wird die CHSH-Schranke $2\sqrt2$ erreicht — nicht
unterschritten, wie es die frühen Dämpfungsansätze vorhergesagt
hatten. Der Bell-Test-Streit über lokalen Realismus löst sich
geometrisch auf: Es gibt keine Nichtlokalität, weil es keinen
getrennten Träger gibt.

\section{Die $\xi$-Abweichung}

Eine kleine Abweichung von der Tsirelson-Schranke ist aus der
$\theta_4$-Projektionsgeometrie des Messacts ableitbar. Der
Kandidat für die elementare Abweichung pro Messung ist:\status{S}
\begin{equation}
  \Delta\mathrm{CHSH}_{\mathrm{elem}} = \frac{\xi}{2\pi}
  \approx 2{,}1\times10^{-5}.
\end{equation}
Das ist die Phasenkosten eines $\theta_4$-Zyklus in geometrischen
Einheiten (Dok.~190, Präzisierung~12). Kumulativ über $N$ Messläufe
wächst diese Abweichung als:
\begin{equation}
  \Delta\mathrm{CHSH}(N) \approx \xi\,\frac{\ln N}{D_f}\cdot2\sqrt2.
\end{equation}
Bei $N=73$ ergibt das $\approx5\times10^{-4}$ — eine andere
Observable als die elementare, beide unter dem NISQ-Rauschen
($\sim10^{-2}$).

\textbf{Experimentelle Reichweite:} Hardware-Messungen (IBM Heron~r2,
Mai 2026, Dok.~147) zeigen Bell-Verletzung $S=2{,}74$ ($96{,}9\,\%$
der Tsirelson-Schranke). Der $\xi$-Effekt selbst erfordert
sub-ppm-stabile fehlerkorrigierte Qubits — nicht mit heutiger
NISQ-Hardware auflösbar.\status{K}

\begin{laierspur}
Die FFGFT sagt: Die starken Quantenkorrelationen sind echt und
korrekt — die QM hat recht. Was sie anders sieht: Die Korrelation
kommt nicht durch eine Verbindung \emph{zwischen} den Teilchen,
sondern weil die Teilchen nie wirklich getrennt waren. Sie sitzen
auf demselben geometrischen Träger wie zwei Punkte auf derselben
Trommel.
\end{laierspur}
\quelle{Dok.~190 (Präz.~11, 12), 230, 353, 364, 365}

% QMB_n06_nogo.tex — Kapitel 6: No-Go-Theoreme
% Quellen: Dok. 074, 055
\chapter{No-Go-Theoreme und was sie zeigen}
\label{ch:nogo}

No-Go-Theoreme begrenzen, was alternative Interpretationen der QM
leisten können. Die FFGFT muss sich an ihnen messen.

\section{Bells Theorem (revisited)}

Bells Theorem (Kapitel~\ref{ch:bell}) schließt lokal-realistische
Theorien mit verborgenen Variablen aus. Die FFGFT umgeht dies
nicht durch eine Lücke, sondern durch eine andere Ontologie: Die
„verborgene Variable" ist die vollständige Trägerkonfiguration
$(z,r,\theta)_i$ aller Moden — aber der Träger ist geteilt, nicht
lokal getrennt.

\section{Kochen-Specker-Theorem}

Das Kochen-Specker-Theorem (1967) zeigt: Einem Quantensystem können
keine definiten Werte für alle Observablen gleichzeitig und
kontextunabhängig zugewiesen werden.\status{E} Messergebnisse sind
kontextabhängig.

In der FFGFT: Der Polübergang bei der Messung hängt von der
vollständigen Träger-Apparatur-Wechselwirkung ab, nicht nur von
der Mode selbst. Kontextabhängigkeit ist damit eine Eigenschaft
der Messdynamik auf dem Träger, nicht ein Beweis für Unvollständigkeit
der Beschreibung.\status{S} Die Born-Regel wird reproduziert;
ob die Kontextabhängigkeit vollständig aus der Geometrie herleitbar
ist, bleibt offen.

\section{PBR-Theorem}

Das Pusey-Barrett-Rudolph-Theorem (2012) zeigt: Wenn der Quanten-
zustand nur epistemisch wäre (Ausdruck des Wissens, nicht der
Realität), würden bestimmte Messungen Ergebnisse erlauben, die die
QM ausschließt.\status{E}

In der FFGFT ist der Zustand $(z,r,\theta)$ eine reale Träger-
konfiguration, kein Wissenszustand — das PBR-Theorem ist damit
kompatibel: Die FFGFT hat eine ontische Wellenfunktion.\status{B}

\section{Hardys Theorem und GHZ-Zustand}

Hardy (1992) und Greenberger-Horne-Zeilinger (1989) konstruieren
Szenarien, in denen jede lokal-realistische Zuweisung zu einem
Widerspruch führt — auch ohne statistische Ungleichungen.\status{E}

Beide Argumente setzen getrennte Träger voraus. In der FFGFT
entfällt diese Voraussetzung: Die Moden teilen einen Träger, und
die GHZ-Korrelationen sind topologisch kodiert.\status{S}

\section{Was die FFGFT nicht behauptet}

Die No-Go-Theoreme zeigen, was nicht geht. Die FFGFT behauptet:
\begin{itemize}
  \item Keine neuen Vorhersagen bei Standard-Bell-Tests (alle
        statistischen Ergebnisse wie in der QM).\status{B}
  \item Eine kleine $\xi$-Abweichung bei Hochpräzisions-Tests
        ($\Delta\mathrm{CHSH}_{\mathrm{elem}}\approx2\times10^{-5}$),
        heute nicht messbar.\status{S}
  \item Keine Verletzung der Messfreiheit (kein Superdeterminismus
        — dieser Ansatz aus Dok.~074 ist überholt).
  \item Keine Information schneller als Licht.
\end{itemize}

\quelle{Dok.~055, 074, 230}

% QMB_n07_chsh_experiment.tex — Kapitel 7: ΔCHSH und Experiment
% Quellen: Dok. 190, 230, 147
\chapter{CHSH-Abweichung und experimentelle Reichweite}
\label{ch:chsh_exp}

\section{Zwei verschiedene Observablen}

Das Korrekturregister (Dok.~190, Präzisierung~11 und 12, Mai 2026)
unterscheidet präzise zwischen zwei verschiedenen Größen, die frühere
Dokumente zusammenwarfen:

\textbf{Elementare Abweichung.} Pro einzelner Messung entstehen
Phasenkosten des $\theta_4$-Projektionszyklus. Der natürliche
Kandidat ist:
\begin{equation}
  \Delta\mathrm{CHSH}_{\mathrm{elem}} = \frac{\xi}{2\pi}
  \approx 2{,}1\times10^{-5}.
\end{equation}
Diese Identifikation ist eine Arbeitshypothese; eine formale Herleitung
aus der $T^4$-Projektionsgeometrie fehlt noch.\status{S}

\textbf{Kumulative Abweichung.} Über $N$ Messläufe wächst die
Abweichung nach der Formel aus Dok.~022:
\begin{equation}
  \Delta\mathrm{CHSH}(N) = 2\sqrt2 \Bigl[1 - \exp\!\Bigl(-\xi\,\frac{\ln N}{D_f}\Bigr)\Bigr]
  \approx \xi\,\frac{\ln N}{D_f}\cdot2\sqrt2.
\end{equation}
Bei $N=73$ ergibt das $\approx5\times10^{-4}$. Diese Größe ist
experimentell zugänglich, sobald Fehlerkorrektur-Schwellenwerte
unterschritten sind.

\section{Experimentelle Situation 2026}

IBM Heron~r2 (Mai 2026, Dok.~147) liefert:
\begin{itemize}
  \item Bell-Verletzung bestätigt: $S=2{,}74$ ($96{,}9\,\%$ der
        Tsirelson-Schranke).\status{K}
  \item FFGFT-Schaltkreise lauffähig auf aktueller Hardware.\status{K}
  \item Geräte-zu-Geräte-Variation: $1{,}4\,\%$ — rund
        $1400\times$ größer als der gesuchte $\xi$-Effekt.\status{K}
\end{itemize}

Der $\xi$-Effekt erfordert sub-ppm-stabile fehlerkorrigierte Qubits.
Mit NISQ-Hardware nicht auflösbar.\status{K}

\section{Offene Herleitung}

Die Identifikation $\Delta\mathrm{CHSH}_{\mathrm{elem}} = \xi/(2\pi)$
ist plausibel, aber nicht aus der $T^4$-Projektionsgeometrie hergeleitet.
Das Verhältnis zwischen kumulativer und elementarer Abweichung bei
$N$ Läufen beträgt:
\begin{equation}
  \frac{\Delta\mathrm{CHSH}(N)}{\Delta\mathrm{CHSH}_{\mathrm{elem}}}
  \approx N\cdot\frac{2\pi}{D_f},
\end{equation}
bei $N=73$ einen Faktor $\approx153$, konsistent mit dem beobachteten
Faktor $\sim50$--$100$.\status{S}

Die formale Herleitung von $\xi/(2\pi)$ aus der Trägergeometrie
bleibt eine offene Brücke.

\quelle{Dok.~022, 147, 190 (Präz.~11, 12), 230}


% --- Teil III ---
\part{Quantenrechner}
% QMB_n08_qubit_formalismus.tex — Kapitel 8: Qubit-Formalismus
% Inhaltliche Grundlage: Dok. 175 (März 2026), Dok. 034, Dok. 230; eigene Darstellung
\chapter{Mehr als zwei Zustände: das Qubit geometrisch}
\label{ch:qubit}

\section{Was steckt zwischen $|0\rangle$ und $|1\rangle$?}

Spin-up und Spin-down sind zwei Punkte. Zwei Punkte sind eindeutig
beschriftet --- ein Bit Information. Aber ein Qubit ist kein Bit. Die
eigentliche Kraft des Qubits liegt nicht darin, dass es $|0\rangle$
\emph{oder} $|1\rangle$ sein kann; das kann eine klassische Münze auch.
Sie liegt darin, dass es beides in einem kontinuierlichen Verhältnis
überlagern kann:
\begin{equation}
  |\psi\rangle = \cos\tfrac{\theta}{2}\,|0\rangle
  + e^{i\phi}\sin\tfrac{\theta}{2}\,|1\rangle.
\end{equation}
Zwei Parameter, der Mischungswinkel $\theta\in[0,\pi]$ und die relative
Phase $\phi\in[0,2\pi)$, parametrieren eine Kugeloberfläche: die
Bloch-Kugel. Auf ihr gibt es unendlich viele Punkte --- nicht nur die
Pole $|0\rangle$ und $|1\rangle$, sondern den gesamten Äquator, alle
Breitenkreise, jeden Punkt dazwischen. Jeder ist ein gültiger,
physikalisch unterscheidbarer Zustand.

\begin{keypoint}[Ein Qubit kodiert nicht 2 Zustände, sondern unendlich viele]
Die zwei klassischen Pole sind Sonderfälle einer kontinuierlichen
Familie. Was ein Qubit vom Bit unterscheidet, ist nicht die Zahl der
diskreten Zustände, sondern die kontinuierliche Geometrie des
Zustandsraums. Das Auslesen koppelt an die Konfiguration und treibt sie
in den nächstgelegenen stabilen Extremzustand. Das Ergebnis ist in
FFGFT-Sicht deterministisch --- es hängt vom genauen Ausgangszustand
ab, den wir vor der Messung nicht vollständig kennen. Die Zustände
selbst sind diskret; die Kontinuität liegt im Phasenraum \emph{vor}
der Messung.
\end{keypoint}

\section{Der zylindrische Phasenraum}

Die Bloch-Kugel ist die Standarddarstellung. Die FFGFT hat eine
physikalisch direktere Alternative: den zylindrischen Phasenraum mit
drei Koordinaten $(z,r,\theta)$. Dabei ist $z\in[-1,+1]$ die Höhe
entlang der Zylinderachse --- der Anteil der Torsionskonfiguration in
Richtung $|0\rangle$ gegenüber $|1\rangle$; $r\in[0,1]$ der Radius im
Querschnitt --- die Stärke der intermediären Torsion, also wie weit der
Zustand von beiden Polen entfernt ist; $\theta\in[0,2\pi)$ der
Azimutwinkel --- die relative Phase zwischen den Basiskomponenten. Die
Beziehung zur Bloch-Kugel ist direkt: $z = \cos\theta_{\mathrm{Bloch}}$,
$r = |\sin\theta_{\mathrm{Bloch}}|$; die formale Brücke steht in
Kapitel~\ref{ch:hilbert}.

Der Zylinder ist kein geometrisches Schönheitsmittel: Er macht
Gatteroperationen zu geometrischen Transformationen im Realraum.

\begin{center}
\small\renewcommand{\arraystretch}{1.4}
\begin{tabular}{lp{2.4cm}p{4.2cm}}
\toprule
\textbf{Gatter} & \textbf{Matrix} & \textbf{Zylindertransformation} \\
\midrule
Hadamard $H$ & $\frac{1}{\sqrt2}\begin{pmatrix}1&1\\1&-1\end{pmatrix}$ & $z\leftrightarrow r$ vertauschen, Phase um $\pi/2$ drehen: bewegt Pol auf Äquator \\
Phase $Z$ & $\begin{pmatrix}1&0\\0&-1\end{pmatrix}$ & $\pi$ zur Phase $\theta$ addieren: Rotation um die Zylinderachse \\
Bit-Flip $X$ & $\begin{pmatrix}0&1\\1&0\end{pmatrix}$ & $(z,r)$-Drehung um $\alpha = \pi K_{\mathrm{frak}}$ \\
\bottomrule
\end{tabular}
\end{center}

Im zylindrischen Formalismus ist jede Gatteroperation eine geometrische
Transformation des Torsionsvektors. Das Hadamard-Gatter dreht keine
Matrix --- es kippt die Konfiguration von der Längsachse auf die
Transversalebene. Das $Z$-Gatter rotiert die Windung um ihre eigene
Achse. Der Faktor $K_{\mathrm{frak}}$ im $X$-Gatter ist kein nachträglich
eingefügter Korrekturfaktor; er folgt aus der fraktalen
Raumzeitgeometrie und bedeutet: Jede Qubit-Rotation ist leicht kleiner
als $\pi$, weil die Raumzeit selbst minimal von der flachen Geometrie
abweicht. Das ist eine prüfbare Vorhersage: Alle $\pi$-Gatter sollten
systematisch um $\Delta\alpha\approx0{,}013\,\%$ vom idealen Wert
abweichen.\status{S}

\section{Quantenregister: exponentielles Wachstum durch Kombination}

Ein einzelnes Qubit hat unendlich viele Zwischenzustände, liefert beim
Auslesen aber immer einen der zwei Pole. Die eigentliche Stärke kommt
durch Kombination.

\textbf{Produktzustände: additive Klassik.} Zwei unabhängige Qubits
haben zusammen vier Basiszustände. Der allgemeine Zustand ohne
Verschränkung ist $(\alpha_A|0\rangle+\beta_A|1\rangle)\otimes
(\alpha_B|0\rangle+\beta_B|1\rangle)$ --- zwei getrennte Bloch-Kugeln,
kein Gewinn über zwei unabhängige Qubits. Der Zustandsraum ist additiv:
$2+2 = 4$ Parameter.

\textbf{Verschränkte Zustände: multiplikatives Wachstum.} Sobald die
Qubits verschränkt sind, kann der gemeinsame Zustand nicht mehr als
Produkt geschrieben werden. Der allgemeine Zweiqubitzustand
\begin{equation}
  |\Psi\rangle = \alpha|00\rangle+\beta|01\rangle+\gamma|10\rangle+\delta|11\rangle
\end{equation}
braucht vier komplexe Amplituden, also $2\cdot4-2 = 6$ reelle Parameter.
Verschränkte Qubits tragen gemeinsame Information, die in keinem der
Einzelqubits steckt. Für $n$ Qubits wächst der Hilbertraum exponentiell,
$d(n) = 2^n$, mit $2\cdot2^n-2$ reellen Parametern:

\begin{center}
\small
\begin{tabular}{rrrl}
\toprule
$n$ Qubits & Basiszustände & Reelle Parameter & Klassisch \\
\midrule
1 & 2 & 2 & 1 Bit \\
2 & 4 & 6 & 2 Bits \\
3 & 8 & 14 & 3 Bits \\
10 & 1024 & 2046 & 10 Bits \\
50 & $10^{15}$ & $2\times10^{15}$ & 50 Bits \\
300 & $10^{90}$ & $2\times10^{90}$ & $\approx$ Atome im Universum \\
\bottomrule
\end{tabular}
\end{center}

Das ist der eigentliche Quantenvorteil: Nicht ein einzelnes Qubit ist
mächtig, sondern das kollektive Wachstum eines verschränkten Registers.

\textbf{Bell-Zustände.} Die vier maximal verschränkten
Zweiqubitzustände bilden eine Orthonormalbasis des verschränkten
Unterraums:
\begin{align}
  |\Phi^\pm\rangle &= \tfrac{1}{\sqrt2}(|00\rangle\pm|11\rangle),\\
  |\Psi^\pm\rangle &= \tfrac{1}{\sqrt2}(|01\rangle\pm|10\rangle).
\end{align}
Misst man Qubit~A, ist das Ergebnis von Qubit~B festgelegt, egal wie
weit entfernt. Im FFGFT-Bild ist ein Bell-Zustand keine Fernwirkung
zwischen zwei Qubits, sondern eine topologisch verbundene
Torsionskonfiguration über zwei Resonatoren: Beide teilen einen
gemeinsamen Torsionsknoten, dessen Auflösung bei jeder lokalen Messung
beide Enden zugleich bestimmt. Das ist kein Signal mit
Überlichtgeschwindigkeit, sondern ein gemeinsamer geometrischer Zustand,
der bei der Präparation etabliert wurde (Kapitel~\ref{ch:verschraenkung}).
Die frühere Dämpfungsformel für die Bell-Korrelation, die den CHSH-Wert
um $\sim10^{-4}$ absenkte, ist durch die Zwei-Observablen-Lesart aus
Kapitel~\ref{ch:chsh_exp} ersetzt.

\textbf{Hardware-Validierung.} Die FFGFT-Vorhersagen für Bell-Zustände
wurden auf IBM Brisbane und IBM Sherbrooke (je 127 Qubits, Heavy-Hex-
Gitter) geprüft: Bell-Zustands-Fidelität $97{,}17\,\%$; FFGFT-Algorithmus
und Standard-QM stimmen innerhalb der Messgenauigkeit überein;
Wiederholvarianz $10^{-4}$ bis $10^{-3}$; Abweichungen in $P(00)$ und
$P(11)$ von $1$--$5\,\%$ sind Hardware-Rauschen.\status{K} Die
$\xi$-Abweichung von der idealen Bell-Korrelation liegt unterhalb der
Rauschschwelle und konnte nicht isoliert werden. Der vollständige
Python-Code ist im Repositorium verfügbar.

\section{Drei Wege zu mehr Zuständen}

Aus der Resonatorphysik ergeben sich drei grundlegend verschiedene
Strategien, mehr als zwei Zustände zu kodieren:

\begin{center}
\small\renewcommand{\arraystretch}{1.3}
\begin{tabular}{p{1.6cm}p{2.3cm}p{1.6cm}p{2.5cm}}
\toprule
\textbf{Strategie} & \textbf{Mehr Zustände durch} & \textbf{Wachstum} & \textbf{FFGFT-Bedeutung} \\
\midrule
Superposition (1 Qubit) & kontinuierliche Überlagerung & $\infty$, aber 1 Bit bei Messung & intermediäre Torsion \\
Register ($n$ Qubits) & Verschränkung & $2^n$ & topologisch verbundene Knoten \\
Qudit ($d>2$) & höhere Niveaus, Orbitale & $d$ pro Objekt & $d$ stabile Konfigurationen in einem Resonator \\
\bottomrule
\end{tabular}
\end{center}

Keine Strategie dominiert überall. Registerverschränkung skaliert am
stärksten, braucht aber $n$ physische Objekte und fehlerkorrigierte
Verschränkungsoperationen. Qudit-Kodierung gewinnt Information pro Objekt
($\log_2d>1$); zwei Ququarts ($d = 4$) tragen $4^2 = 16$ Zustände, vier
Qubits $2^4 = 16$ --- mathematisch äquivalent, physikalisch verschieden
robust. Superposition ist immer da: Sie ist die Grundlage der
Interferenz, der eigentlichen Triebkraft von Quantenalgorithmen.

\section{Quanteninterferenz: warum Superposition nützlich ist}

Superposition allein wäre eine bloße Überlagerung; das Auslesen lieferte
immer einen der zwei Pole. Was Superposition zur Ressource macht, ist die
Interferenz: Amplituden können sich konstruktiv verstärken oder
destruktiv auslöschen, je nach Phase.\status{E} Das ist das Prinzip
hinter allen nützlichen Quantenalgorithmen: Man präpariert das Register
in einer gleichmäßigen Zwischentorsion über alle Eingaben (Hadamard);
wendet eine Funktion an, die die Phasen der richtigen Antworten erhöht
und die der falschen erniedrigt; lässt Interferenz die falschen
Antworten in instabile Konfigurationen treiben und die richtige in eine
stabile; und das Auslesen treibt das System in den nächsten stabilen
Pol, der nach der Interferenz mit überwältigender Häufigkeit der
gesuchten Antwort entspricht. Ohne Interferenz: Monte Carlo. Mit
Interferenz: Grover ($\sqrt N$ statt $N$ Schritte), Shor (Faktorisierung
in polynomialer Zeit).

In der FFGFT ist die Phase keine abstrakte komplexe Zahl, sondern die
Rotationsstellung der intermediären Torsionskonfiguration. Konstruktive
Interferenz: Zwei Torsionsphasen liegen gleich und verstärken sich zu
einer geometrisch stabileren Konfiguration. Destruktive Interferenz:
Zwei entgegengesetzte Phasen löschen sich aus; die resultierende
Konfiguration hat keinen stabilen geometrischen Anker und wird beim
Auslesen nicht angesteuert. Ein Quantenalgorithmus ist in dieser
Sprache eine Folge von Torsionsrotationen, die die Phasenkonfiguration
des Registers so ausrichtet, dass die gesuchte Antwort in einer stabilen
Konfiguration endet.

\begin{laierspur}
Wasserwellen, die sich treffen, verstärken sich, wenn Berg auf Berg
trifft, und löschen sich aus, wenn Berg auf Tal trifft. Ein
Quantenrechner steuert Wahrscheinlichkeiten wie Wellen: Er lässt die
richtigen Antworten wachsen und die falschen verschwinden. Das geht nur,
weil Quantenzustände eine Phase haben, die klassischen Bits fehlt.
\end{laierspur}

\section{QND-Messung: Auslesen ohne Zerstören}

Die Auslesewechselwirkung treibt die intermediäre Konfiguration in
einen Pol --- für invasive Messverfahren. Es gibt eine Klasse von
Verfahren, die diesen Schritt nicht vollziehen: die
Quanten-Nichtdemolitionsmessung (QND). Sie liest eine Observable $\hat A$
aus, ohne den Eigenzustand zu verändern; die formale Bedingung ist
$[\hat A,\hat H_{\mathrm{Messung}}] = 0$.\status{E} Praktisch: Derselbe
Zustand kann wiederholt gemessen werden und liefert jedes Mal dasselbe
Ergebnis.

Das wichtigste Beispiel ist die Faraday- oder Kerr-Rotation: Ein linear
polarisierter Laserstrahl passiert den Quantenpunkt ohne Absorption und
dreht seine Polarisationsebene um $\theta_F\propto S_z$. Kein Photon wird
absorbiert, kein Energieübertrag in die Spinmode, der Spinzustand ist
vor und nach der Messung identisch.

\textbf{QND als empirischer Beleg für Determinismus.} Das Argument ist
direkt: Die Faraday-Rotation misst $S_z$ ohne den Zustand zu verändern;
wiederholte Messungen liefern immer dasselbe Ergebnis; also hatte der
Zustand diesen Wert vor jeder Messung, definit, nicht als abstrakte
Superposition, die erst durch Messung Realität wird. In der FFGFT ist
die Faraday-Rotation zustandserhaltend, weil eine stabile
Torsionskonfiguration durch eine dispersive, nichtresonante Kopplung
nicht aus ihrem Zustand getrieben werden kann. Das Photon sieht die
Torsion --- es dreht seine Polarisation entsprechend ---, überträgt aber
keine Energie, weil es nicht resonant ist.\status{S}

\begin{center}
\small\renewcommand{\arraystretch}{1.3}
\begin{tabular}{p{2.4cm}p{3cm}p{2.8cm}}
\toprule
\textbf{Verfahren} & \textbf{Wechselwirkung} & \textbf{Zustand danach} \\
\midrule
Resonante Fluoreszenz & Photon absorbiert, regt Übergang an & verändert \\
Spin-zu-Ladungs-Konversion & Elektron tunnelt aus dem Resonator & zerstört \\
Faraday/Kerr & Photon passiert dispersiv & unverändert: QND \\
\bottomrule
\end{tabular}
\end{center}

Die Standardlesart, der Zustand existiere erst durch Messung, ist mit
QND-Messungen schwer vereinbar. In der FFGFT ist die
Torsionskonfiguration eine geometrische Realität: stabil, definit,
präexistent. Die Messung liest sie ab; sie erschafft sie nicht.

\textbf{Grenzen der QND.} QND ist möglich für Observablen, die mit dem
Messoperator kommutieren, typisch $S_z$. Eine intermediäre Konfiguration
--- ein Zustand, der kein $S_z$-Eigenzustand ist --- wird durch eine
$S_z$-QND-Messung in einen Eigenzustand getrieben, deterministisch
abhängig vom Ausgangszustand; die Zwischentorsion ist danach weg.

\section{Die Potentiallandschaft}

Das vollständige Bild der Qubit-Zustände ist eine Potentiallandschaft
mit zwei stabilen Minima und einem Sattel dazwischen. Das Minimum
$|\uparrow\rangle$ ist eine stabile Konfiguration, tief genug, dass
schwache Kopplungen den Zustand nicht herausheben --- QND möglich. Das
Minimum $|\downarrow\rangle$ ist spiegelbildlich gleich tief. Der Abhang
ist die intermediäre Torsion: Die Phase bestimmt, wie weit der Zustand
von welchem Minimum entfernt ist; kein echtes Minimum, metastabil auf
der Kohärenzzeit $T_2$. Jeder Punkt dazwischen ist ein geometrisch
definierter Zustand, kontinuierlich parametriert, diskret in den
Endpunkten.

Die Auslesestärke entscheidet: Ein Zustand im stabilen Minimum bleibt
bei schwacher Kopplung (QND) und kann bei starker Kopplung herausgehoben
werden (invasiv). Ein Zustand auf dem Abhang rollt schon bei schwacher
Kopplung ins nächste Minimum --- welches, bestimmt die genaue Phase,
epistemisch unbekannt, deterministisch; bei starker Kopplung dasselbe,
nur schneller.

Die entscheidende Einsicht: Es gibt keine prinzipielle Grenze zwischen
\enquote{QND-möglich} und \enquote{QND-unmöglich}, nur stabilere und
weniger stabile Zustände in einer kontinuierlichen Landschaft. Der
Sprung in die Endposition ist nicht zufällig: Ein Zustand nahe
$|\uparrow\rangle$ landet in $|\uparrow\rangle$, einer nahe
$|\downarrow\rangle$ in $|\downarrow\rangle$, einer auf dem Sattel
abhängig von der feinen Phase. Die scheinbare Zufälligkeit der
Standard-QM ist die epistemische Unkenntnis dieser Phase. Das Bild gilt
für alle $d$ Niveaus eines Qudits: $d$ Minima mit $d-1$ Sattelpunkten.

\section{Zwei Ebenen der $\xi$-Hierarchie: Spin-Qubit und photonischer Resonator}

Die bisherigen Abschnitte haben das Qubit am Spin-Qubit entwickelt ---
einem Elektron in einem Nanometerkäfig. Es gibt eine qualitativ andere
Realisierungsklasse: den photonischen Resonator. Beide erzwingen
diskrete Zustände durch Geometrie, sitzen aber auf verschiedenen Stufen
der $\xi$-Hierarchie, mit Konsequenzen für Temperatur, Stabilität und
die Art der Information.

\textbf{Spin-Qubit: minimaler Resonator auf Stufe $N\approx8$.} Ein
einzelnes Elektron, ein Energieniveau, zwei Orientierungen. Der Abstand
der Spinzustände in einem typischen Feld von 1\,T ist
$\Delta E_Z = g^*\mu_BB\approx0{,}1$--$1$\,meV; die thermische Energie
bei Raumtemperatur ist $k_BT = 26$\,meV. Also $\Delta E_Z\ll k_BT$:
Thermisches Rauschen wirft das System ständig aus dem Minimum. Tiefe
Temperaturen unter 1\,K sind nicht optional, sondern zwingend --- nicht
um Superposition zu erhalten, sondern um überhaupt in einem der Minima
zu bleiben.

\textbf{Photonischer Resonator: kollektive Geometrie auf Stufe
$N\approx9$--$10$.} Ein Wellenleiter aus Lithiumniobat, ein
Mach-Zehnder-Interferometer, ein Ringresonator --- Resonatoren aus
$10^{18}$ bis $10^{23}$ Atomen in Gitterordnung. Der Resonator ist nicht
ein einzelnes Teilchen, sondern die makroskopische Geometrie des ganzen
Kristalls; das Photon gehört der kollektiven Mode. Die Diskretheit kommt
aus den Randbedingungen der makroskopischen Geometrie. Die Energie eines
optischen Photons ist $1$--$2$\,eV, also $E/k_BT\approx40$: Thermisches
Rauschen kann keine optischen Photonen erzeugen. Photonik funktioniert
bei Raumtemperatur nicht, weil sie robuster gebaut ist, sondern weil
sie auf einer höheren $\xi$-Stufe operiert: Die geometrische Robustheit
ist durch den Einkristall mit $10^{23}$ Gitterplätzen gegeben, die
thermische durch die eV-Skala --- vier Größenordnungen Abstand.

\textbf{Nicht auf zwei Zustände beschränkt.} In einem Mach-Zehnder-
Interferometer läuft das Photon durch zwei Pfade mit kontinuierlich
einstellbarer Phase; das Interferenzergebnis $\cos^2(\phi/2)$ ist eine
kontinuierliche Übertragungsfunktion. Ein photonischer Resonator kann
beide Modi betreiben: digital mit $\phi\in\{0,\pi\}$ und Transmission
1 oder 0, oder analog mit kontinuierlicher Phase --- unendlich viele
Zwischenzustände, alle bei Raumtemperatur stabil. Beim Spin-Qubit liegt
die analoge Information in einem thermisch instabilen Zwischenzustand,
$T_2$ bricht bei Raumtemperatur in Nanosekunden zusammen; beim
photonischen Resonator liegt sie in der Phase des Photons, und die ist
thermisch stabil.

\begin{center}
\small\renewcommand{\arraystretch}{1.3}
\begin{tabular}{p{2.6cm}p{2.8cm}p{3cm}}
\toprule
 & \textbf{Spin-Qubit} & \textbf{Photonischer Resonator} \\
\midrule
$\xi$-Stufe & $N\approx8$ (nm) & $N\approx9$--10 ($\mu$m--mm) \\
Resonator & 1 Elektron im Käfig & $10^{18}$--$10^{23}$ Atome kollektiv \\
Energieskala & $0{,}1$--1\,meV & 1--2\,eV \\
$\Delta E/k_BT$(300\,K) & $0{,}004$--$0{,}04$ & $\approx40$ \\
Kühlung nötig? & ja, $T<1$\,K & nein \\
Analoge Zustände & fragil ($T_2\sim$ns) & stabil \\
Diskretheit aus & Confinement (1 Teilchen) & Modengeometrie \\
Typische Anwendung & digitales Qubit & analoge Signalverarbeitung, photonisches Rechnen \\
\bottomrule
\end{tabular}
\end{center}

Spin-Qubit und photonischer Resonator sind in der FFGFT keine
grundlegend verschiedenen Systeme, sondern Resonatoren auf verschiedenen
Stufen derselben Skalierungsleiter. Jede Stufe trägt die Energie um den
Faktor $1/\xi\approx7500$ höher als die vorherige; von meV (Stufe 8) zu
eV (Stufe 9--10) ist der Faktor $10^3$--$10^4$, konsistent mit der
Skalierung.\status{S} Das erklärt auch, warum biologische Systeme
photonische Mechanismen (Photosynthese, Vogelnavigation) und keine
Spin-Qubits für Quanteneffekte nutzen: Die Natur arbeitet auf der
thermisch robusten Stufe.

\section{Die analoge Auflösungsgrenze}

Der photonische Resonator ist thermisch robust, aber nur, solange man
nicht zu fein auflöst. Wer analoge Zwischenwerte unterscheiden will,
teilt den Energieraum in $N$ Stufen mit $\Delta E_{\mathrm{Stufe}} =
E_{\mathrm{Photon}}/N$, und die Bedingung $\Delta E_{\mathrm{Stufe}}\gg
k_BT$ ergibt $N\ll E_{\mathrm{Photon}}/k_BT$. Bei Raumtemperatur und
optischen Photonen: $N_{\max}\approx40$, entsprechend
$\log_2 40\approx5{,}3$\,Bit pro Photon. Das photonische System entkommt
dem thermischen Problem nicht --- es verschiebt es auf eine feinere
Skala. Mehr Auflösung erfordert Kühlung oder eine noch höhere
Energieskala.

Das ist die älteste Ingenieursweisheit der Nachrichtentechnik,
geometrisch begründet: Die Hochfrequenztechnik klettert seit hundert
Jahren dieselbe Hierarchie hoch. Bei jeder Frequenzstufe verbessert sich
das Signal-Rausch-Verhältnis, weil $E/k_BT$ mit jeder Stufe wächst.

\begin{center}
\footnotesize\renewcommand{\arraystretch}{1.25}
\begin{tabular}{p{2.2cm}p{1.3cm}p{1.6cm}p{1.2cm}p{2.9cm}}
\toprule
\textbf{Technologie} & \textbf{Frequenz} & \textbf{$E/k_BT$} & \textbf{$N$-Stufe} & \textbf{Regime} \\
\midrule
AM-Rundfunk & 1\,MHz & $1{,}5\times10^{-7}$ & $\approx5{,}5$ & thermisch dominiert \\
UKW & 100\,MHz & $1{,}5\times10^{-5}$ & $\approx6{,}4$ & rauschbegrenzt \\
Mikrowelle & 10\,GHz & $1{,}5\times10^{-3}$ & $\approx7{,}3$ & beherrschbar \\
5G/mmWave & 30\,GHz & $5\times10^{-3}$ & $\approx7{,}5$ & Übergang \\
THz & 1\,THz & $0{,}15$ & $\approx7{,}8$ & $E\sim k_BT$: schwierig \\
Nahinfrarot & 200\,THz & 31 & $\approx8{,}8$ & thermisch robust \\
Glasfaser & 300\,THz & 46 & $\approx9{,}0$ & $N_{\max}\approx46$ \\
UV & 1000\,THz & 154 & $\approx9{,}4$ & $N_{\max}\approx150$ \\
\bottomrule
\end{tabular}
\end{center}

Von $N = 6$ (AM) nach $N = 9$ (Optik) überbrücken drei Stufen den Faktor
$(1/\xi)^3\approx4\times10^{11}$; tatsächlich liegt die Glasfaserfrequenz
$10^8$-fach über AM-Rundfunk. Die Glasfaser ist in diesem Bild keine
neue Technologie, sondern die Fortsetzung desselben Trends auf die
nächste Stufe. Und die Auflösungsgrenze gilt auf jeder Stufe: Wer bei
optischen Wellenlängen mehr als 40 analoge Stufen auflösen will, muss
kühlen --- wie der Radioingenieur 1930 seinen Empfänger abschirmte. Die
universelle Grenze lautet
\begin{equation}
  N_{\max}(T,N_\xi)\approx\frac{E_0\,(1/\xi)^{N_\xi-N_0}}{k_BT}.
\end{equation}

\textbf{Oberhalb der Optik.} Die Optik ist nicht das Ende der Leiter,
aber das Ende der nutzbaren Zone für geführte photonische Übertragung.
Im nahen UV (3--10\,eV, $N\approx9{,}3$) beginnt Photoionisation: Der
Resonator wird durch sein eigenes Signal beschädigt. Bei weichem Röntgen
(100\,eV--1\,keV, $N\approx10$) nähert sich die Wellenlänge dem
Gitterabstand, Bragg-Streuung an einzelnen Gitterebenen, die
Wellenleitung bricht zusammen. Bei hartem Röntgen fliegt das Photon
durch die Gitterlücken, Materie wird teils transparent, es gibt keinen
Brechungsindex und keinen Resonator. Bei Gammaenergien wechselwirkt das
Photon nur noch mit Atomkernen --- eine andere $\xi$-Stufe. Oberhalb
1\,MeV ($N\approx11{,}3$) setzt Paarerzeugung ein: die natürliche
Obergrenze. In der FFGFT gibt es keine Kanonenkugeln und keine Wellen als
getrennte Konzepte, nur Torsionskonfigurationen auf verschiedenen Skalen;
Wellenleitung ist eine Skalenresonanz zweier Torsionsmuster, des
Kristalls auf Stufe 7--8 und des Photons auf Stufe 9, und sie bricht
zusammen, wenn die Skalen inkompatibel werden.

\begin{keypoint}[Was dieses Kapitel festhält]
Das Qubit ist eine Mode im zylindrischen Phasenraum; Gatter sind
Torsionsrotationen; Register wachsen exponentiell durch topologische
Verbindung; Interferenz ist Phasenüberlagerung von Torsionen; QND-Messung
belegt die Präexistenz stabiler Pole; die Potentiallandschaft macht den
Messausgang deterministisch bei epistemisch unbekannter Phase. Spin-Qubit
und Photonik sind Resonatoren auf verschiedenen $\xi$-Stufen, und die
thermische Auflösungsgrenze gilt auf jeder Stufe gleich.
\end{keypoint}

\quelle{Dok.~175 (Qubit-Zustandsräume, März 2026), Dok.~034 (zylindrischer Formalismus), Dok.~173 (Resonatorbedingungen), Dok.~174 (Torsion, Faraday-Auslese), Dok.~230}

% QMB_n09_algorithmen.tex — Kapitel 9: Quantenalgorithmen
% Inhaltliche Grundlage: Dok. 073, 147, 176, 316, 343; eigene Darstellung
\chapter{Quantenalgorithmen: Deutsch, Grover, Shor und die Grenzen der Periode}
\label{ch:algorithmen}

\section{Die vier Standardalgorithmen}

Vier Algorithmen bilden das Lehrstück der Quantenrechnung. Sie zeigen,
was Quanteninterferenz leisten kann und was sie nicht leisten kann.

\subsection*{Deutsch: ein Bit in einer Abfrage}

Der Deutsch-Algorithmus beantwortet die Frage, ob eine unbekannte
Funktion $f:\{0,1\}\to\{0,1\}$ konstant oder balanciert ist, mit einer
einzigen Abfrage --- klassisch braucht man zwei.

In der Standard-QM: Man präpariert $|0\rangle$, wendet Hadamard an,
fragt das Orakel $U_f$ ab, wendet Hadamard an und misst. Ergebnis $0$:
konstant; Ergebnis $1$: balanciert. Der Trick ist das Phasenkickback:
$U_f|x\rangle = (-1)^{f(x)}|x\rangle$. Das Hadamard danach macht die
globale Phasendifferenz sichtbar.

In der FFGFT: Dieselbe Sequenz auf dem zylindrischen Phasenraum. Das
Orakel dreht die Phase des Torsionsvektors; Hadamard vertauscht $z$ und
$r$ und dreht um $\pi/2$. Das Messresultat ist identisch. Es gibt keinen
algorithmischen Unterschied --- nur der konzeptuelle Rahmen ist ein
anderer: Die Phase ist keine abstrakte komplexe Zahl, sondern die
Rotationsstellung der Torsionskonfiguration.\status{K}

\subsection*{Grover: Suche in $\sqrt N$ Schritten}

Grover sucht in einer unsortierten Datenbank mit $N$ Einträgen in $O(\sqrt N)$
Schritten statt $O(N)$ klassisch. Der Mechanismus: Man präpariert das
Register in einer gleichmäßigen Superposition aller $N$ Einträge,
markiert den gesuchten Eintrag durch Phasenkickback
($U_\mathrm{Grover}:|x\rangle\to(-1)^{\delta_{x,x_0}}|x\rangle$) und
spiegelt dann an der gleichmäßigen Superposition (Inversionsspiegel um
den Mittelwert). Jeder Grover-Schritt erhöht die Amplitude des gesuchten
Eintrags auf Kosten aller anderen. Nach $\pi/4\cdot\sqrt N$ Schritten
ist die Wahrscheinlichkeit des richtigen Eintrags nahe 1.

FFGFT-Einordnung: Grover-Iterationen sind in der FFGFT eine Folge von
Torsionsrotationen, die die Phase des gesuchten Eintrags systematisch
ausrichten. Die anderen Einträge landen in instabilen Konfigurationen;
die Auslesewechselwirkung treibt das System deterministisch in den
ausgerichteten Pol.\status{K} Die Vorhersagen sind identisch mit der
Standard-QM.

\subsection*{Shor: Faktorisierung in polynomialer Zeit}

Shors Algorithmus faktorisiert $N$ durch Periodenfindung. Sei $a$ mit
$\gcd(a,N) = 1$; die Funktion $f(x) = a^x\bmod N$ ist periodisch mit
Periode $r$. Die Quantenfouriertransformation (QFT) macht die Periode in
den Messwahrscheinlichkeiten sichtbar; die Kettenbruchnäherung extrahiert
$r$ aus dem Messwert. Aus $r$ folgen die Faktoren über $\gcd(a^{r/2}\pm1, N)$.

Zentrale Aussage des Korpus: Der Algorithmus zerfällt in zwei Vorgänge,
die nichts miteinander zu tun haben:\status{K}

\begin{center}
\small
\begin{tabular}{p{2cm}p{3cm}p{3cm}}
\toprule
 & \textbf{Aufbereitung} & \textbf{Transformation} \\
\midrule
Aufgabe & Modulare Exponentiation $a^x\bmod N$ & Periodenextraktion \\
Natur & arithmetisches Problem & Wellenoperation (Interferenz) \\
Aufwand & teuer ($O(n^3)$ Gatter) & billig ($O(n^2)$ Gatter) \\
\bottomrule
\end{tabular}
\end{center}

Die Aufbereitung macht $\approx95\,\%$ der Gatterkosten aus; die QFT nur
$\approx5\,\%$.\status{K} Shor spezifizierte nur den klassischen Prozess
$(a,1)\to(a,a^x\bmod N)$ --- die Implementierung der modularen
Exponentiation in Superposition ist das eigentliche Ingenieusproblem und
hängt nicht an quantenmechanischen Eigenschaften.

\section{$\phi$-QFT und Äquivalenzsatz}

Dok.~147 entwickelt eine alternative Fouriertransformation, die Phasen
$2\pi/\varphi^k$ statt $2\pi/2^k$ verwendet, wobei $\varphi$ der
Goldene Schnitt ist:
\begin{equation}
  \phi\text{-QFT}:|x\rangle\mapsto\frac{1}{\sqrt{Q_\varphi}}\sum_{y=0}^{Q_\varphi-1}e^{2\pi ixy/Q_\varphi}|y\rangle,
\end{equation}
mit $Q_\varphi = \varphi^n$ statt $Q = 2^n$.

\textbf{Hauptsatz:} Für Shors Algorithmus zur Faktorisierung von
$N < 2^{20}$ mit Fehlerwahrscheinlichkeit $\delta < 10^{-6}$ gilt:\status{K}
\begin{equation}
  P_{\mathrm{success}}(\text{Standard-QFT})\le P_{\mathrm{success}}(\phi\text{-QFT})\le P_{\mathrm{success}}(\text{Standard-QFT})+\xi.
\end{equation}
Die $\phi$-QFT ist zur Standard-QFT äquivalent, mit einer maximalen
Abweichung von $\xi\approx1{,}3\times10^{-4}$ --- vernachlässigbar.
Der Äquivalenzsatz zeigt: $\xi$ hat in $\mathbb{Z}/N\mathbb{Z}$ keine
strukturelle Funktion. Die Periode $r$ ist durch $\mathrm{lcm}(p-1,q-1)$
bestimmt, eine individuelle arithmetische Eigenschaft von $N$, nicht ein
universeller geometrischer Parameter.

\section{Die strukturellen Grenzen der Periode}

Drei Grenzen setzen dem, was Quantenalgorithmen leisten können, präzise
Schranken.

\textbf{1. Weyl-Gesetz und Zähldichte.}\status{K} Die Eigenwertdichte
jedes kompakten Zustandsraums wächst nach einem Potenzgesetz
$N(\lambda)\sim\lambda^{d/2}$. Die Primzahlzählfunktion $\pi(T)$ und die
Riemannschen Nullstellen wachsen dagegen wie $T\ln T$. Das ist keine
Potenz. Ein endlicher Zustandsraum kann ein solches Spektrum nicht
vollständig tragen. Die Weyl-Obstruktion: Ein $n$-Qubit-Register ist ein
kompakter Zustandsraum; arithmetische Zähldichten wachsen schneller als
jede kompakte Geometrie erlaubt.

\textbf{2. Aufbereitung als Flaschenhals.}\status{K} Die QFT macht
Periodizitäten sichtbar, die in den präparierten Zuständen bereits
vorhanden sind. Der limitierende Schritt ist die Zustandspräparation
--- die Auswertung von $a^x\bmod N$ --- nicht die Transformation selbst.
Kein geometrischer Parameter $\xi$ kann die arithmetische Komplexität
der Aufbereitung reduzieren: Sie ist eine Funktion der Primstruktur von
$N$, nicht der Geometrie des Trägers.

\textbf{3. Kein Quantenvorteil für RSA (Stand 2026).}\status{X} Ein
Vorteil muss auf der Shor-Faktorisierungskurve sichtbar werden: Der
Schwellenwert der Fehlerkorrektur muss demonstriert und gehalten werden;
die Aufbereitung muss in Superposition ausgeführt werden; das Ergebnis
muss mit RSA-relevanten Schlüssellängen (2048 Bit) repliziert werden.
Nichts davon ist Stand 2026 realisiert. Google Willow unterschreitet
zwar den Fehlerkorrektur-Schwellenwert für logische Qubits --- aber in
einem Regime weit unterhalb der für Shor nötigen Skalierung.

\section{Riemann-Zeta und FFGFT-Träger}

Dok.~316 und Dok.~343 klären die Beziehung zwischen der Riemann-Zeta
und dem FFGFT-Träger.

\textbf{Die Riemann-Zeta als Faktor der Torus-Spektral-Zeta.}\status{K}
Die geschlossene Form der $\mathbb{Z}^4$-Spektral-Zeta lautet:
\begin{equation}
  Z_{\mathbb{Z}^4}(s) = 8\,(1-4^{1-s})\,\zeta(s)\,\zeta(s-1).
\end{equation}
Das ist eine Identität. Die Nullstellen der Torus-Spektral-Zeta haben
drei Quellen: die Riemannschen Nullstellen (aus $\zeta(s)$), ihre
Verschobenen (aus $\zeta(s-1)$) und die Pol-Null $s = 1$. Die
Riemann-Hypothese ist damit eine Aussage über ein geometrisches Objekt
des Rahmens --- aber eine, die die Geometrie nicht beantworten kann:
Der Torus kann die Nullstellen als Faktor tragen, aber nicht erzeugen.

\textbf{$\mathbb{Z}_3$-Orbifold-Projektion.}\status{B} Die
$Z_3$-Projektion von $T^4/\mathbb{Z}_3$ transformiert die volle
Riemann-Zeta in eine Dirichlet-$L$-Funktion:
\begin{equation}
  \zeta_{T^4/\mathbb{Z}_3}(s) = L(s,\chi_0) = (1-3^{-s})\,\zeta(s).
\end{equation}
Das Euler-Produkt dieser $L$-Funktion ist nach der Eintrittstiefe
$k = \mathrm{ord}_p(3)$ hierarchisch geordnet; die physikalisch
sichtbaren Primzahlen $\{2,3,5,7,11,13\}$ sind genau jene mit $k\le6$.

\textbf{Galois-Struktur und Neutrinos.}\status{B} Die vier
Frobenius-Konjugationsklassen primitiver Elemente in $\mathrm{GF}(27)^*$
werden durch die FFGFT-Sektorpaarung $k\leftrightarrow-k$ zu zwei Klassen
gepaart: $\{f_3,f_4\}$ (Dirac, geladene Leptonen) und $\{f_1,f_2\}$
(Majorana, $\nu_1,\nu_2$ aktiv; $\nu_3$ in Dirac-Orbit $O_4$). Vorhersage:
$\nu_3$ hat keine Majorana-Phase.

\textbf{Weyl-Grenze.}\status{K} Keine kompakte Geometrie kann die
Nullstellen der Riemann-Zeta als Spektrum tragen. Die Nullstellendichte
$N(T)\sim T\ln T$ verlangt eine Dilatationsstruktur --- die ausgerollte
Domäne. Die Geometrie zeigt die Adresse an, liefert aber nicht die Antwort.
Drei naheliegende geometrische Zugänge sind ausgeschlossen: Keine
kompakte Geometrie liefert das $T\ln T$-Wachstum; $\xi$ verschiebt die
kritische Linie nicht; die $\xi$-Leiter diskretisiert die Nullstellen nicht.

\section{Deterministische Implementierung am PC}

Die Zustandsbrücke $(z,r,\theta)\leftrightarrow(\alpha,\beta)$ ist
bijektiv (Kapitel~\ref{ch:hilbertraum}). Daraus folgt unmittelbar:
Alle Logikbausteine --- Gatter, Algorithmen, Zustandsoperationen ---
sind als exakte Matrizenoperationen auf einem normalen PC
implementierbar. Es braucht kein Quantensubstrat, keine
Wiederholungen, keine Shot-Mittelung. Jedes Ergebnis ist in einem
einzigen Durchlauf vollständig bestimmt.

Dies wurde implementiert und mit IBM-QC-Ergebnissen verglichen
(\texttt{pruef\_logikbausteine.py}):

\begin{description}[leftmargin=2em,style=sameline,itemsep=0.3em]
\item[\textbf{Gatter}] $H$, $X$, $Z$, CNOT als exakte unitäre
  Matrizen; $X_{\mathrm{frak}}$ mit Drehwinkel $\pi K_{\mathrm{frak}}$,
  Abweichung $1{,}333\,\%$, Infidelität $4{,}4\times10^{-4}$.
  Ein Matrixprodukt --- kein Sampling.

\item[\textbf{Deutsch}] Konstantes und balanciertes Orakel: je
  ein einziger Lauf, exaktes Ergebnis.

\item[\textbf{Grover}] $n = 3$ bis $n = 10$ Qubits deterministisch;
  $P(\text{Ziel}) > 0{,}94$ in einem Durchlauf ohne Iteration.
  Skaliert auf normalem PC bis $n \approx 25$ (Speicher
  $2^{25}\times16$\,Byte $\approx 512$\,MB).

\item[\textbf{Bell-Zustand}] $P(00) = P(11) = 0{,}5$ exakt in
  einem Lauf. IBM-Kingston: Fidelität $0{,}9876$ über
  $50\times2048$ Shots. Der QC brauchte 100.000 Messungen;
  der PC: einen einzigen Durchlauf.

\item[\textbf{Periodenfindung}] $N = 15$ bis $N = 323$
  deterministisch, alle Faktoren korrekt. Kein Sampling,
  kein QC.\status{K}
\end{description}

Die deterministischen Ergebnisse stimmen mit den IBM-QC-Resultaten
innerhalb der Messgenauigkeit überein --- nicht annähernd, sondern
exakt. Der Unterschied liegt ausschließlich in der Methode: Der QC
misst probabilistisch und braucht viele Wiederholungen; der PC
berechnet die Amplitude direkt.

\section{Klassischer PC versus Quantenrechner}

\textbf{Deterministischer Einzeldurchlauf.}\status{K}
Die FFGFT-Zustandsabbildung $(z,r,\theta)\leftrightarrow(\alpha,\beta)$
ist bijektiv (Kapitel~\ref{ch:hilbertraum}). Daraus folgt: Alle
algorithmischen Schritte --- Periodenfindung, Faktorextraktion,
Gatteroperationen --- sind auf dem PC in einem einzigen Durchlauf
vollständig bestimmt. Kein Wiederholen, keine Shot-Mittelung, keine
Statistik. Der IBM-Kingston-Vergleich brauchte auf der QC-Seite
50~Läufe à 2048~Shots, weil der QC probabilistisch misst; die
FFGFT-Simulation am PC: ein einziger Lauf, identisches Ergebnis.

\textbf{Aufbereitung: PC schneller als QC.}\status{K}
Die modulare Exponentiation $a^x\bmod N$ ist auf dem PC direktes
klassisches Rechnen ohne Overhead. Der Quantenschaltkreis für
dieselbe Aufgabe ist durch Reversibilität, Uncomputation und
Toffoli-Gatter nachweislich \emph{teurer} als die klassische
Implementierung (Dok.~176). Shor hat die Aufbereitung in
Superposition nicht spezifiziert; sie ist und bleibt ein
arithmetisches Problem in jeder Technologie gleich aufwändig.

\textbf{Klassische Verfahren übertreffen die Periodensuche.}\status{K}
Die FFGFT-Periodensuche sequenziell: $O(N)$, gemessener
Wachstumsexponent der mittleren Ordnung $0{,}95$. Pollard~$\rho$
braucht 3--7 Schritte für dieselbe Aufgabe, Probeteilung 58--180.
Die klassische Resonanzsuche scheitert zudem doppelt: $O(N)$
Schritte \emph{und} eine Frequenzauflösung von $1/N^2$, die bei
RSA-Größen unterhalb jeder Zahlendarstellung liegt
(\texttt{s1\_periodensuche.py}, Dok.~176).

\textbf{Verschränkung trägt den Rechenvorteil nicht.}\status{K}
Gottesman-Knill (1998): Jeder Clifford-Schaltkreis ($H$, $S$,
CNOT) ist klassisch in Polynomzeit simulierbar --- unabhängig von
der erzeugten Verschränkung. Ein GHZ-Zustand mit $n = 1000$
Qubits ist maximal verschränkt und braucht klassisch nur $2n^2$
Bits Speicher. Was den Quantenvorteil trägt, ist ungeklärt;
Kontextualität und Nicht-Stabilisiertheit sind Kandidaten, keiner
ist etabliert.

\textbf{Die einzige Ausnahme: QFT in Superposition.}\status{K}
Der einzige echte Quantenvorteil liegt ausschließlich darin,
dass die QFT die Periode über alle $x$ gleichzeitig in
Superposition ausliest --- $O((\log N)^3)$ statt $O(N)$. Das
kann ein klassischer PC nicht replizieren. Dies ist die
\emph{einzige} Stelle, an der der QC schneller wäre.

\textbf{Derzeit nicht erreichbar.}\status{K}
Für RSA-2048 wären $4{,}1\times10^6$ physikalische Qubits und
$8{,}6\times10^9$ Gatter nötig --- eine Ressourcengrenze,
prinzipiell überwindbar, aber Stand~2026 um Größenordnungen
nicht erreicht. Google~Willow unterschreitet den
Fehlerkorrektur-Schwellenwert für logische Qubits, liegt aber
weit unterhalb der für Shor nötigen Skalierung. Die QFT-Ausnahme
existiert theoretisch; sie ist praktisch irrelevant.

\section{Experimentelle Validierung}

IBM Kingston (28. Mai 2026), 50 unabhängige Läufe, 2048 Shots pro Lauf:
Bell-Zustands-Fidelität $0{,}9876\pm0{,}0028$. Die Wiederholbarkeitsvarianz
ist statistisch verträglich mit gewöhnlichem Quanten-Shot-Noise (Verhältnis
beobachtete/erwartete Varianz: $1{,}02$). Eine frühere Auswertung von 3
Läufen hatte fälschlicherweise $40\times$ Determinismus gegenüber der QM
berichtet; das war Stichproben-Artefakt.\status{K}

Die $\xi$-Korrektur an CHSH-Werten liegt bei $\Delta\sim10^{-5}$ pro
Messung, weit unter NISQ-Rauschen. Die häufig zitierten kumulativen
Werte für $N = 73$ Qubits (Dok.~022 und Dok.~147) und die elementare
Abweichung $\xi/(2\pi)$ sind nicht direkt vergleichbar und dürfen nicht
gleichgesetzt werden (Kapitel~\ref{ch:chsh_exp}).

\quelle{Dok.~073 (Algorithmen in T0, 2025), Dok.~147 §3 ($\phi$-QFT-Äquivalenz), §4 (Bell-Tests), §5 (Shor), §8 (IBM-Validierung, Mai 2026), Dok.~176 (Aufbereitung vs.\ Transformation), Dok.~316 (Riemann, Weyl), Dok.~343 (Zeta-Galois)}

% QMB_n10_spin_qubit.tex — Kapitel 10: Quantenpunkt, Spin, Stabilität
% Inhaltliche Grundlage: Dok. 173 (März 2026), Dok. 174, Dok. 178; eigene Darstellung
\chapter{Quantenpunkt, Spin und die Grenze der Superposition}
\label{ch:spin}

\section{Was ein Quantenpunkt ist: Geometrie erzwingt Diskretheit}

Ein Quantenpunkt ist ein künstliches Atom --- ein Hohlraum so klein,
dass seine Geometrie diskrete Energieniveaus erzwingt. Nicht die Zahl
der Atome bestimmt die Niveaus, sondern die Form des Hohlraums: wie die
Töne einer Flöte hängen von der Länge des Rohrs ab, nicht von der Zahl
der Atome im Metall. Die Energieniveaus des eingesperrten Teilchens:
\begin{equation}
  E_n = \frac{n^2\pi^2\hbar^2}{2m_{\mathrm{eff}}R^2},\qquad n = 1,2,3,\dots
\end{equation}
Der Faktor $\pi^2$ folgt aus der Geometrie: Eine stehende Welle, die am
Rand verschwindet, erzwingt $k\cdot R = \pi$ für die Grundmode. Die
Diskretheit kommt aus der Geometrie, nicht aus einem Postulat.

In der FFGFT ist ein Elektron keine Punktladung in einem Kasten, sondern
eine stabile Torsionskonfiguration des zugrundeliegenden Feldes. Die
diskreten Niveaus sind die erlaubten geometrischen Konfigurationen dieser
Feldverdrehung in einem begrenzten Raum. Zwischen den Niveaus gibt es
keine stabile Konfiguration; das Feld existiert nur in definierten
Mustern, nicht in beliebigen Zwischenzuständen.

Der fundamentale Parameter $\xi = 4/30000$ erzeugt eine universelle
Längenhierarchie $L_N = L_0\cdot(1/\xi)^N$ beginnend bei der
Fundamentallänge $L_0 = \xi\ell_P$. Stufe $N = 8$ entspricht
$L_8\approx22$\,nm; der exzitonische Bohr-Radius für CdSe liegt
bei $N\approx7{,}85$ --- derselben Skalenstufe. Die Skala der Exzitonen
ist in der FFGFT keine empirische Größe, sondern geometrisch
erzwungen.\status{S}

\begin{center}
\small
\begin{tabular}{cll}
\toprule
$N$ & $L_N$ & Physikalische Domäne \\
\midrule
0 & $2{,}15\times10^{-39}$\,m & Fundamentale FFGFT-Länge \\
1 & $1{,}62\times10^{-35}$\,m & Planck-Länge \\
6 & $3{,}8\times10^{-16}$\,m & Subnuklear \\
7 & $2{,}9\times10^{-12}$\,m & Atomare Skala \\
\textbf{8} & $\mathbf{2{,}2\times10^{-8}}$\,\textbf{m} & Quantenpunkt- und Exzitonbereich \\
9 & $1{,}6\times10^{-4}$\,m & Zelluläre Skala \\
10 & $\approx1$\,m & Makroskopisch \\
\bottomrule
\end{tabular}
\end{center}

\section{Zwei Bedingungen für Quantenpunkt-Verhalten}

Zwei Bedingungen müssen beide erfüllt sein:

\textbf{Bedingung 1 --- Geometrische Perfektion} (absolut, temperatur-
unabhängig): Der Resonator muss präzise genug sein, um stehende Wellen
zu erzwingen. Ein fehlerbehaftetes Gitter zeigt kein Confinement-
Quantenverhalten, egal wie klein, egal bei welcher Temperatur. Bedingung 1
ist die härtere; sie kann nicht durch Kühlen ersetzt werden.

\textbf{Bedingung 2 --- Niveauabstand gegen Wärme} (größen- und
temperaturabhängig): Der Abstand $\Delta E = 3E_1$ zwischen den untersten
zwei Niveaus muss die thermische Energie übersteigen: $\Delta E > k_BT$.

Aus Bedingung 2 folgt die maximale Resonatorgröße:
\begin{equation}
  R_{\mathrm{bulk}}(T) = \pi\hbar\sqrt{\frac{3}{2m_{\mathrm{eff}}k_BT}}.
\end{equation}
Für Silizium bei 300\,K: $R_{\mathrm{bulk}}\approx15$\,nm --- genau dort,
wo klassische Transistoren 2012 aufgehört haben zu skalieren.

\begin{center}
\small
\begin{tabular}{lrrrrr}
\toprule
\textbf{Material} & $m_{\mathrm{eff}}/m_e$ & \multicolumn{4}{c}{$R_{\mathrm{bulk}}$ bei} \\
\cmidrule(lr){3-6}
 & & 4\,K & 77\,K & 300\,K & 310\,K \\
\midrule
GaAs & 0{,}067 & 221\,nm & 50\,nm & 25{,}5\,nm & 25{,}1\,nm \\
CdSe & 0{,}130 & 159\,nm & 36\,nm & 18{,}3\,nm & 18{,}0\,nm \\
Silizium & 0{,}190 & 131\,nm & 30\,nm & 15{,}2\,nm & 14{,}9\,nm \\
\bottomrule
\end{tabular}
\end{center}

Kühlen auf 4\,K verschiebt die Grenze für Silizium von 15\,nm auf
131\,nm --- ein Faktor~9. Damit werden makroskopisch wirkende Objekte zu
Quantenobjekten: Ein 50-nm-Siliziumkristall ist bei Raumtemperatur
klassisch, bei 4\,K quantisiert. Ein makroskopischer Körper hat nicht
keine Niveaus, sondern $10^{23}$ so dichte, dass $k_BT$ sie alle
überdeckt. Die Diskretheit verschwindet nicht; sie wird so feinkörnig,
dass sie als Bandstruktur erscheint.

\section{Transistoren und die $\xi$-Stufentabelle}

Als die Halbleiterindustrie 2012 auf FinFET umstellen musste, weil
klassische planare Transistoren aufgehört hatten zu skalieren, hatten die
Kanallängen die $\xi$-Stufe $N = 8{,}00$ unterschritten:

\begin{center}
\small
\begin{tabular}{p{3.5cm}ccp{2.8cm}}
\toprule
\textbf{Technologie} & \textbf{Kanallänge} & \textbf{$N$-Stufe} & \textbf{Regime} \\
\midrule
65\,nm (2005) & 65\,nm & 8{,}12 & klassisch \\
22\,nm FinFET (2012) & 22\,nm & \textbf{8{,}00} & Quanteneinflüsse beginnen \\
10\,nm (2016) & 10\,nm & 7{,}91 & Tunneln messbar \\
5\,nm (2020) & 5\,nm & 7{,}84 & QD-Regime \\
3\,nm GAA (2023) & 3\,nm & 7{,}78 & echter QD-Kanal \\
2\,nm (2025) & 2\,nm & 7{,}73 & $\approx$10 Atomlagen \\
\bottomrule
\end{tabular}
\end{center}

\section{Biologische Quanteneffekte}

Lebende Zellen nutzen Quanteneffekte bei 37\,°C. Drei Klassen mit
verschiedenen Bedingungen: Confinement-Quantisierung braucht beide
Bedingungen; Tunneln und Spin-Kohärenz haben eigene Kriterien.
Bedingung~1 (geometrische Präzision) gilt für alle drei ohne Ausnahme.

\begin{center}
\small
\begin{tabular}{p{1.8cm}p{2.2cm}p{2.2cm}p{3cm}}
\toprule
\textbf{Effekt} & \textbf{Bed.1: Geometrie} & \textbf{Bed.2: $\Delta E > k_BT$} & \textbf{Mechanismus} \\
\midrule
Photosynthese (FMO) & \checkmark\ Proteinhülle & grenzwertig ($J/k_BT\approx0{,}7$) & Exzitonkohärenz, $J\approx19$\,meV \\
Vogelnavigation & \checkmark\ Kryptochrom & kein Confinement & Spinverschränkung, Relaxationszeit $>$ Reaktionszeit \\
Enzymtunneln & \checkmark\ Aktivzentrum & kein Confinement & Quantentunneln durch Barriere \\
\bottomrule
\end{tabular}
\end{center}

Der biologische Betriebspunkt bei 37\,°C liegt nicht wegen einer
thermischen Energiegrenze dort; $R_{\mathrm{bulk}}$ ändert sich zwischen
0\,°C und 57\,°C um weniger als 10\,\%. Was bei 42--45\,°C kollabiert,
ist die geometrische Präzision: Proteine denaturieren, die
Torsionskonfiguration bricht zusammen. 37\,°C ist der Punkt, an dem
biologische Resonatoren (Proteinkavitäten $\sim2$--4\,nm) maximal chemisch
aktiv sind, ohne ihre Konfiguration zu verlieren.

\section{Spin beeinflussen: fünf Methoden}

Spin-up und Spin-down sind das absolute Minimum eines Quantenresonators:
zwei erlaubte Orientierungen einer Torsionskonfiguration. Fünf
Methoden, diese Orientierung zu drehen:

\textbf{Magnetfeld (Zeeman-Effekt).} Das statische Feld $B$ spaltet die
Spinzustände um $\Delta E_Z = g^*\mu_BB$ auf; ein resonanter
Hochfrequenzpuls bei der Larmor-Frequenz dreht den Spin gezielt. In der
FFGFT legt das Feld eine bevorzugte Richtung im Torsionsfeld fest und
bricht die Symmetrie zwischen Up und Down energetisch auf.\status{E}

\textbf{Elektrisches Feld (EDSR).} Elektrische Felder koppeln über die
Spin-Bahn-Kopplung: Das Elektron sieht in seinem Ruhesystem ein
effektives Magnetfeld, das den Spin dreht. Vorteil: Gate-Elektroden
erzeugen Felder auf wenigen Nanometern lokal und schalten in unter einer
Nanosekunde. Für integrierte Qubit-Arrays ist das die bevorzugte
Methode.\status{E}

\textbf{Zirkulär polarisiertes Licht.} $\sigma^+$-Photonen tragen
Drehimpuls $+\hbar$ und übertragen ihn bei resonanter Anregung auf das
Exziton und damit auf den Elektronenspin. Rein optische Spinpräparation
ohne äußeres Feld; Auslesen über die Polarisation des Fluoreszenzlichts.
In der FFGFT: Das Photon trägt seinen Drehimpuls als topologische
Eigenschaft der elektromagnetischen Feldtorsion; bei der Absorption
rotiert sie die Elektronentorsion.\status{E}

\textbf{Austauschwechselwirkung.} Wenn zwei Elektronen in benachbarten
Quantenpunkten räumlich überlappen, koppeln ihre Spins mit $H_J =
J(t)\mathbf{S}_1\cdot\mathbf{S}_2$. Durch Pulsieren des Gate-Feldes
(das den Überlapp und damit $J(t)$ steuert) führt man gezielte
Zwei-Qubit-Operationen aus --- die Basis für Spin-Qubit-Gatter in
Halbleitersystemen (Loss-DiVincenzo, 1998).\status{E}

\textbf{Hyperfein-Wechselwirkung.} Das Elektron wechselwirkt mit den
Kernspins der $\sim10^4$--$10^6$ umliegenden Atome. Das fluktuierende
Kernspinbad zerstört die Spinpolarisation normalerweise in $T_2^*\sim
10$--100\,ns --- ein Problem. Kontrolliert eingesetzt, etwa durch
dynamische Entkopplung (schnelle Pulsfolgen, die Kernspin-Fluktuationen
mitteln), verlängert sich $T_2$ auf Mikrosekunden. In Phosphor-in-Silizium
wird der Kernspin selbst als hochstabiles Qubit genutzt
($T_1 > 1$\,s bei tiefen Temperaturen).\status{E}

\section{Spin auslesen: fünf Wege}

\textbf{Spin-zu-Ladungs-Konversion.} Das Potential wird so gewählt,
dass nur Spin-down energetisch erlaubt ist, aus dem Quantenpunkt zu
tunneln; das tunnelnde Elektron erzeugt einen messbaren Strompuls im
benachbarten Einzelelektronen-Transistor. Genauigkeiten $>99{,}5\,\%$ bei
Auslesezeiten unter 1\,\textmu{}s; ausreichend für fehlerkorrigierte
Quantencomputer nach dem Surface-Code-Schwellenwert.\status{E}

\textbf{Pauli-Blockade.} Zwei benachbarte Quantenpunkte, je ein Elektron.
Versucht man, beide in denselben Punkt zu laden, geht das bei
antiparallelen Spins (beide passen ins Grundniveau), nicht bei parallelen
(Pauli-Verbot). Messbar als Leitwert-Unterschied. In der FFGFT: Zwei
gleiche Torsionskonfigurationen können nicht denselben Resonator belegen,
nicht wegen eines Verbots, sondern weil dieselbe Torsionswindung im
Hohlraum keine zweite stabile Konfiguration derselben Art zulässt. Das
Pauli-Prinzip ist eine geometrische Aussage.

\textbf{Faraday- und Kerr-Rotation.} Ein linear polarisierter
Laserstrahl passiert den Quantenpunkt ohne Absorption und dreht seine
Polarisationsebene um $\theta_F\propto S_z$. Das ist eine
Quantum-Non-Demolition-Messung (QND); der Spinzustand ist unverändert.
Zeitauflösung im Pikosekundenbereich; einzelne Spins nachweisbar.
Wiederholte Messungen liefern dasselbe Ergebnis.\status{E}

\textbf{Resonante Fluoreszenz.} Das Photon wird absorbiert und regt den
Übergang an --- invasiv, der Zustand ist danach verändert oder zerstört.

\textbf{Abbildende Methoden.} Für Ensembles: MOKE (magneto-optischer
Kerr-Effekt), Zeitaufgelöste Faraday-Rotation, Spin-Rausch-Spektroskopie.

\section{Dekohärenz: die Grenze der Superposition}

Zwei Zeitskalen beschreiben den Zerfall von Quantenzuständen.
$T_1$ (Relaxationszeit) ist die Zeit, bis das System aus dem angeregten
Zustand in den Grundzustand gefallen ist. $T_2$ (Kohärenzzeit) ist die
Zeit, bis die Phase der Superposition zerstört ist; stets $T_2\le2T_1$.

\begin{center}
\small
\begin{tabular}{p{3.2cm}p{1.3cm}p{1.3cm}p{3.3cm}}
\toprule
\textbf{System} & \textbf{$T_1$} & \textbf{$T_2$} & \textbf{Hauptmechanismus} \\
\midrule
Elektron-Spin GaAs & $\sim1$\,ms & $\sim10$\,ns & Kernspin-Bad \\
Elektron-Spin Si/SiGe & $>1$\,s & $\sim10$\,ms & ${}^{28}$Si (kernspinfrei) \\
Elektron-Spin P:Si & $>1$\,s & $>100$\,ms & isotopische Reinigung \\
Exziton CdSe-QD & $\sim1$\,ns & $\sim10$\,ps & Phonon-Streuung \\
Supraleiter-Transmon & $\sim100$\,\textmu{}s & $\sim100$\,\textmu{}s & dielektrische Verluste \\
\bottomrule
\end{tabular}
\end{center}

Der Vergleich zwischen Exziton ($T_2\sim10$\,ps) und Spin-Qubit in Si
($T_2\sim10$\,ms) ergibt neun Größenordnungen. Mit dem Higgs-Kopplungs-
parameter $\xi_{\mathrm{Higgs}}\approx1{,}038\times10^{-5}$ überbrücken
zwei Hierarchiestufen exakt $9{,}97$ Größenordnungen --- übereinstimmend
mit dem gemessenen Verhältnis auf $<5\,\%$.\status{S} Die Deutung: Der
Exziton-Resonator koppelt direkt an seine Umgebung (Phononen derselben
Stufe), der Spin-Qubit an ein entferntes Bad (Kernspins auf Stufe
$N\approx7$), also zwei Stufen schwächer.

Silizium mit angereichertem ${}^{28}$Si eliminiert das Kernspin-Bad fast
vollständig. Was übrig bleibt, ist die intrinsische Geometrie des
Resonators. Die lange Kohärenzzeit in ${}^{28}$Si ist ein Beleg dafür,
dass geometrische Perfektion (Bedingung~1) die entscheidende ist.\status{K}

\section{Superposition: fragil und nie der stabile Grundzustand}

Hier liegt der eigentliche konzeptuelle Kern. Die Standardtheorie sagt:
Ein einzelnes Elektron befindet sich in echter Superposition
$\alpha|\uparrow\rangle+\beta|\downarrow\rangle$ bis zur Messung. Auf
jedem Kühlschrank klebt ein Permanentmagnet. Zwischen diesen Tatsachen
liegt eine Spannung.

Ein Permanentmagnet ist die makroskopische Summe von Spinausrichtungen.
Die Austauschwechselwirkung zwingt benachbarte Spins zur parallelen
Ausrichtung; das ist energetisch günstiger. Diese kollektive Ausrichtung
bildet Weiss-Domänen. Im Permanentmagneten sind die Domänen durch ein
äußeres Feld dauerhaft eingefroren. Was der Permanentmagnet direkt
beweist: Spin hat bevorzugte, stabile Ausrichtungen; diese Stabilität
ist makroskopisch robust bei Raumtemperatur über Jahre; kollektive
Superposition tritt nie auf.

Die Standardtheorie erklärt das mit Dekohärenz: Die Wechselwirkung mit
der Umgebung zerstört die Superposition so schnell, dass sie praktisch
nicht beobachtbar ist. Das ist technisch korrekt, aber eine nachträgliche
Erklärung: Sie beginnt mit Superposition als Grundzustand und erklärt
dann, warum dieser Grundzustand immer sofort zerstört wird.

Die FFGFT beginnt umgekehrt: Stabilität ist der Grundzustand.
Superposition ist die begrenzte, fragile Ausnahme mit messbaren
Zeitskalen und einer geometrischen Ursache.

\begin{center}
\small
\begin{tabular}{p{3.2cm}p{1.3cm}p{1.3cm}p{1.3cm}p{2.3cm}}
\toprule
\textbf{System} & \textbf{$T_1$} & \textbf{$T_2$} & \textbf{Temp.} & \textbf{Bemerkung} \\
\midrule
Exziton CdSe-QD & $\sim1$\,ns & $\sim10$\,ps & 4\,K & Superposition: Pikosekunden \\
Elektron-Spin GaAs & $\sim1$\,ms & $\sim10$\,ns & mK & Superposition: Nanosekunden \\
Elektron-Spin Si & $>1$\,s & $\sim10$\,ms & mK & Superposition: Millisekunden \\
Kernspin P:Si & $>1$\,s & $>100$\,ms & mK & bestes bekanntes System \\
\textbf{Permanentmagnet} & \textbf{Jahre} & --- & \textbf{300\,K} & \textbf{Stabilität ohne Kühlung} \\
\bottomrule
\end{tabular}
\end{center}

Die Tabelle zeigt: Kein einziges experimentelles System zeigt
Superposition als dauerhaft stabilen Zustand. Selbst im besten Qubit
(Kernspin P:Si bei mK) zerfällt die Superposition nach Sekunden. Der
Permanentmagnet hält seinen stabilen Spinzustand über Jahre. Stabilität
ist energetisch bevorzugt; Superposition ist metastabil.

\textbf{Das Einzelelektron.} Die FFGFT-Position: Ein einzelnes Elektron
hat zu jedem Zeitpunkt eine definite Torsionskonfiguration --- entweder
Spin-up oder Spin-down. Was als Superposition $\alpha|\uparrow\rangle+
\beta|\downarrow\rangle$ beschrieben wird, ist epistemische Unkenntnis
über diesen definitiven Zustand, keine ontologische Gleichzeitigkeit. Der
Beweis kommt aus der QND-Messung (Faraday-Rotation, Kapitel~\ref{ch:qubit}):
Ein Spinzustand kann wiederholt gemessen werden, ohne ihn zu verändern,
und liefert immer dasselbe Ergebnis. Wäre das Elektron wirklich in
Superposition, würde eine QND-Messung bei jeder Wiederholung zufällige
Ergebnisse liefern.\status{S}

\textbf{Vom Einzelspin zum Permanentmagneten.} Der Übergang ist in der
FFGFT kein Rätsel, sondern derselbe Mechanismus auf verschiedenen Skalen:
Stabilität ist geometrisch begründet, Superposition ist metastabile
Zwischenkonfiguration, Dekohärenz ist der Prozess der die intermediäre
Torsion in eine der stabilen Extremkonfigurationen drängt. Im
Permanentmagneten haben sich $10^{23}$ Torsionskonfigurationen kollektiv
auf dieselbe Orientierung eingestellt und stabilisieren sich gegenseitig.
Die makroskopische Ausrichtung ist nicht trotz, sondern wegen der
geometrischen Kopplung stabil.

\quelle{Dok.~173 (Quantenpunkt, März 2026), Dok.~174 (Spin und Qubit-Zustände), Dok.~178 (Spin-Stabilität)}

% QMB_n11_hardware.tex — Kapitel 11: Hardware
% Inhaltliche Grundlage: Dok. 083, 161, 183, 184; eigene Darstellung
\chapter{Hardware: Coherent Ising Machine, Google Willow, p-Bit, Photonenchip}
\label{ch:hardware}

\section{Coherent Ising Machine: Optimierung als physikalische Relaxation}

Die Coherent Ising Machine (CIM) ist ein optischer Analogcomputer, der
kombinatorische Optimierungsprobleme durch physikalische Relaxation in
den Grundzustand eines Ising-Hamiltonians löst. Das fundamentale Prinzip:
Das System \emph{fällt} in seine Lösung, es berechnet sie nicht.

\subsection*{Der Ising-Hamiltonian als universelle Optimierungssprache}

Jedes NP-harte kombinatorische Problem lässt sich als Minimierung des
Ising-Hamiltonians formulieren:
\begin{equation}
  H_{\mathrm{Ising}} = -\tfrac12\sum_{i\ne j}J_{ij}\sigma_i\sigma_j
  - \sum_i h_i\sigma_i,
\end{equation}
wobei $J_{ij}$ die Kopplungsstärken und $h_i$ externe Felder sind.
Der Grundzustand --- die Spin-Konfiguration mit minimalem $H_{\mathrm{Ising}}$
--- entspricht der optimalen Lösung. Die Abbildung ist exakt, keine
Approximation: Travelling Salesman, MaxCut, Portfolio-Optimierung,
Proteinstrukturvorhersage, RSA-Faktorisierung sind sämtlich als
Ising-Minimierung formulierbar.

\subsection*{Physikalische Implementierung}

Die CIM implementiert das Ising-Modell durch ein Netzwerk zeitgemulti-
plexter degenerierter optischer parametrischer Oszillatoren (DOPO):
Spin $+1$ ist ein Laserpuls mit Phase $0°$, Spin $-1$ ein Puls mit
Phase $180°$; die Kopplungen $J_{ij}$ werden durch elektronisches
Feedback realisiert. Das System entwickelt sich deterministisch gemäß
den Feldgleichungen des parametrischen Oszillators und findet dabei mit
hoher Wahrscheinlichkeit den Grundzustand. Bis zu 100\,000 optische
Spins sind realisiert (Stand 2024).\status{E}

Eine CIM scheint stochastisch zu suchen; verschiedene Läufe liefern
manchmal verschiedene Lösungen. Aber physikalisch folgt das System
\emph{deterministisch} den Feldgleichungen. Die scheinbare Zufälligkeit
ist ausschließlich Unkenntnis über die genauen Anfangsbedingungen des
Quantenvakuums --- eine direkte Parallele zur FFGFT-Messung.

\subsection*{Sechs Strukturelle Parallelen zur FFGFT}

Zwischen CIM und FFGFT lassen sich sechs Parallelen identifizieren:

\textbf{Parallele 1: Determinismus hinter scheinbarer Stochastik.} In
der CIM ist die Lösung physikalisch erzwungen; in der FFGFT ist das
Messergebnis durch die Feldkonfiguration geometrisch bestimmt. Beide
ersetzen fundamentalen Probabilismus durch versteckten Determinismus.

\textbf{Parallele 2: Grundzustand als geometrisches Minimum.} Die CIM
relaxiert in das Energieminimum des Ising-Hamiltonians; die FFGFT
beschreibt das Universum als im Grundzustand seiner eigenen Geometrie
mit $\xi = 4/30000$. Optimierung ist Relaxation, nicht Iteration.

\textbf{Parallele 3: Qubit-Struktur.} CIM-Spin $\pm1$ ist die Phase
des Laserpulses; FFGFT-Qubit $|\tilde T_\uparrow,m_\downarrow\rangle$
und $|\tilde T_\downarrow,m_\uparrow\rangle$ ist durch $\tilde T\cdot m = 1$
intrinsisch erzwungen. Der FFGFT-Spin muss nicht extern auf $\pm1$
begrenzt werden; er ist es durch die Geometrie.

\textbf{Parallele 4: Die Kopplungsmatrix.} In der klassischen CIM sind
$J_{ij}$ extern programmiert; in einer FFGFT-basierten Ising-Maschine
würden sie geometrisch aus dem Zeitfeld emergieren:
$J_{ij}^{(\mathrm{FFGFT})} = \int\tilde T_i\cdot\tilde T_j\cdot G(x_i,x_j)\,d^3x$.
Das bedeutet: keine Kodierungsfehler möglich, da die Kopplungen reale
physikalische Wechselwirkungen wären.\status{S}

\textbf{Parallele 5: Fraktale Korrektur als Annealing.} Der Korrekturfaktor
$K_{\mathrm{frak}}^{3/2}$ spielt in der FFGFT die analoge Rolle des
Annealing-Temperaturparameters: Er verhindert, dass das System in einem
lokalen Minimum der geometrischen Energie verbleibt.

\textbf{Parallele 6: Toroidale Topologie.} Die FFGFT beschreibt den
Träger als $T^4/\mathbb{Z}_3$. Toroidale Randbedingungen hätten für
eine physikalische Ising-Maschine einen Vorteil: kein Randproblem, alle
Spins haben gleich viele Nachbarn, natürliche Implementierung
translationsinvarianter Kopplungen.

\section{Google Willow: topologische Fehlerkorrektur und Selbstmessung}

Googles Willow-Quantenprozessor erzielte 2024 und 2025 drei Durchbrüche,
die aus FFGFT-Sicht besondere Bedeutung haben.

\subsection*{Fehlerkorrektur unter dem Schwellenwert (Dezember 2024)}

Willow (105 supraleitende Transmon-Qubits, $\sim10$\,mK) zeigt erstmals,
dass die logische Fehlerrate exponentiell sinkt, wenn die Code-Distanz
$d$ erhöht wird.\status{E} Ein $d = 7$-Block erzielt eine Fehlerrate um
Faktor $2{,}14$ unter dem $d = 5$-Block. Physikalische Einzel-Qubit-
Fehlerraten liegen bei $0{,}05\,\%$ --- halb so viel wie auf dem
Vorgänger Sycamore. Die Fehlerkorrektur basiert auf einem Surface Code
mit topologisch geschützten logischen Qubits.

FFGFT-Einordnung: Mehr physische Qubits reduzieren den logischen Fehler
exponentiell --- das ist der experimentelle Beleg, dass topologisch stabile
Konfigurationen im Quantensystem existieren und erreichbar sind, nicht
durch externe Rauschunterdrückung, sondern durch geometrische Ordnung.
Das ist die Richtung, auf die die FFGFT mit ihren Modenklassen des
$D_4$-Gitters zeigt (Kapitel~\ref{ch:gitter}).\status{S}

\subsection*{Floquet topologische Ordnung (September 2025)}

Will et al.\ realisierten auf Willow eine Floquet-topologisch geordnete
Phase --- eine Quantenphase, die im thermischen Gleichgewicht verboten
ist, aber unter periodischem äußerem Antrieb entsteht.\status{E} Auf bis
zu 58 Qubits entstanden: chirale Kantenmoden (Informationsausbreitung in
ausgezeichneter Richtung), anyonische Anregungen (e-Anyonen werden zu
m-Anyonen und umgekehrt), und ein interferometrisch messbarer
topologischer Bulk-Invariant.

FFGFT-Einordnung: Die nicht-gleichgewichtliche Ordnung trägt eine
$\mathbb{Z}_2$-Symmetrie auf dem zweidimensionalen Gitter --- strukturell
verwandt mit der $\mathbb{Z}_3$-Trialität des FFGFT-Trägers, aber nicht
identisch (Dimension, Symmetriegruppe und Trägergeometrie verschieden).
Die Floquet-topologische Ordnung ist ein experimentelles Analogon der
Modenstruktur auf $T^4/\mathbb{Z}_3$.\status{S}

\subsection*{Quantum Echoes (Oktober 2025)}

Out-of-Time-Order-Korrelatoren (OTOC) auf Willow zeigen harmonische
Ausbreitung von Quanteninformation durch konstruktive Interferenz --- ein
Echo-Muster, das periodisch zurückkehrt. FFGFT-Einordnung: Im
Torusformalismus sind Moden auf einer geschlossenen Geometrie
definiert; Informationsausbreitung kehrt zur Quelle zurück --- das Echo
ist eine Konsequenz der Kompaktheit, nicht ein Zufall.

\section{Das p-Bit: zwischen Bit und Qubit}

Das p-Bit (probabilistisches Bit) besetzt den Zwischenbereich zwischen
dem klassischen Bit (tiefe Potenzialsenke, $U\gg k_BT$, stabil über
Jahre) und dem Qubit (metastabile Superposition, $U\ll k_BT$, zerfällt in
Pikosekunden bis Millisekunden). Es ist eine magnetische Nanostruktur auf
der geometrischen Skala $L_8\approx22$\,nm, bei der die Torsionsbarriere
thermisch überwindbar wird.

\subsection*{Physikalische Grundlage}

Die Energiebarriere ist die Stoner-Wohlfarth-Energie $U = K_u\cdot V$,
mit der magnetischen Anisotropieenergie $K_u$ und dem Volumen $V$. Das
Schaltverhalten folgt dem Néel-Arrhenius-Gesetz:
$\tau = \tau_0\exp(U/k_BT)$. Bei 22\,nm und Raumtemperatur liegt $\tau$
im Bereich von Mikrosekunden bis Millisekunden --- kein stabiles Bit,
aber auch kein kohärentes Qubit.

Die FFGFT-Einordnung: Das Schaltverhalten ist deterministisch auf der
sub-Planck-Zeitskala $t_0 = t_P/7500\approx3{,}6\times10^{-47}$\,s;
was als thermische Stochastizität erscheint, ist die makroskopisch
unaufgelöste Abtastung dieser deterministischen Dynamik. Die scheinbare
Zufälligkeit ist epistemisch, nicht ontologisch.\status{S}

\subsection*{Rekonfigurierbare Funktion}

Durch Wahl der magnetischen Vorspannung kann dasselbe p-Bit als
stochastischer Schalter, als erzwungener Schalter oder als analoger
Komparator betrieben werden. Das macht p-Bits zu Bausteinen für
probabilistische Schaltkreise (p-Bits-Netzwerke), die wie Boltzmann-
Maschinen trainiert werden können und NP-harte Probleme ähnlich wie CIMs
durch physikalische Relaxation angehen.

FFGFT-Einordnung: Das p-Bit ist kein Quantenobjekt (kein kohärenter
Überlagerungszustand), sondern eine thermisch fluktuierende klassische
Mode. Im Torusformalismus entspricht es einer Mode, deren Phase $\theta$
durch thermisches Rauschen vollständig randomisiert ist --- der Grenzfall
maximaler Dekohärenz auf dem Träger.

\subsection*{Warum 22 nm: die geometrische Schwelle}

Die kritische Größe für ein p-Bit bei Raumtemperatur entspricht der
FFGFT-Stufe $N = 8{,}00$, also $L_8\approx22$\,nm. Das ist dieselbe
Schwelle, bei der klassische Transistoren 2012 die Quantengrenze
erreichten (Kapitel~\ref{ch:spin}). Drei unabhängige Phänomene konvergieren
auf dieselbe geometrische Skala: Transistorgrenze, exzitonischer
Bohr-Radius, p-Bit-Schwelle.\status{S}

\section{Photonenchip: Raumtemperatur-Quantenphotonik}

Chinas CHIPX/Touring-Quantum-Chip (2025) ist ein monolithischer
6-Zoll-Dünnfilm-Lithiumniobat (TFLN)-Wafer mit über 1000 optischen
Komponenten. Kerndaten: $1000\times$-Speedup für parallele Tasks,
$100\times$ geringerer Energieverbrauch, stabil bei Raumtemperatur,
$12000$ Wafer pro Jahr Produktion.\status{E}

\textbf{Warum Raumtemperatur funktioniert:} Optische Photonen haben
$E\approx1$--2\,eV; $E/k_BT(300\,\text{K})\approx40$. Thermisches
Rauschen kann keine optischen Photonen erzeugen; das System sitzt weit
unterhalb der thermischen Schwelle (Kapitel~\ref{ch:qubit}). Der TFLN-
Einkristall mit $\sim10^{23}$ Gitterplätzen ist geometrisch extrem
robust.\status{K}

\textbf{FFGFT-Einordnung:} Masselose Moden (Photonen) sind in der FFGFT
Dreier-Orbits in $GF(27)^*$ --- strukturell von den massiven Fermionen
(Fixpunkte $\{+1,-1\}$) getrennt (Frobenius-Trennung, Dok.~339).\status{B}
Die photonische Plattform eignet sich für schlupflochfreie Bell-Tests und
ist damit der natürliche Prüfstand für die $\xi$-Abweichung
$\Delta\mathrm{CHSH}\approx\xi/(2\pi)\approx2\times10^{-5}$ --- wenn
Detektionseffizienz und Kohärenzzeiten ausreichen.

\textbf{Vorgeschlagene FFGFT-Optimierungen:} Drei Ansätze sind konzeptuell
skizziert, alle noch nicht realisiert: (1) Topogie-Compiler, der
topologische Qubit-Platzierung nach der $D_4$-Gitterstruktur vornimmt
und SWAP-Gatter um 30--50\,\% reduzieren könnte; (2) harmonische Resonanz
mit Qubit-Frequenzen auf $\xi$-Vielfachen; (3) Zeitfeld-Modulation zur
aktiven Kohärenzerhaltung. Status: [S] für alle drei.

\quelle{Dok.~083 (Photonenchip 2025), Dok.~161 (Coherent Ising Machine), Dok.~183 (Google Willow, April 2026), Dok.~184 (p-Bit), Dok.~339 (Frobenius-Trennung), Dok.~343 (Auflösungsboden)}

% QMB_n12_grenze.tex — Kapitel 12: Grenzen des Quantenrechnens
% Inhaltliche Grundlage: Dok. 075, 076, 147, 316, 343; eigene Darstellung
\chapter{Grenzen des Quantenrechnens: RSA, Weyl und der Auflösungsboden}
\label{ch:grenze}

\section{Was Quantencomputer leisten und was nicht}

Dieses Kapitel zieht drei Grenzen, die der Korpus mit Begründung
festlegt. Sie sind keine spekulativen Einschränkungen, sondern
Konsequenzen mathematischer Sätze oder empirischer Befunde.

\section{Grenze 1: RSA-Faktorisierung und der arithmetische Träger}

RSA (Rivest, Shamir, Adleman 1977) beruht auf der Schwierigkeit der
Faktorisierung großer Semiprimzahlen $N = p\cdot q$. Shors Algorithmus
löst das Problem in polynomialer Zeit --- in Prinzip. Die Frage ist, was
der Geometrieparameter $\xi$ dabei leisten kann.

Das Ergebnis des Korpus ist klar: \textbf{$\xi$ hat in $\mathbb{Z}/N\mathbb{Z}$
keine Funktion.}\status{K} Die Periode $r$ der modularen Potenzfunktion
$f(x) = a^x\bmod N$ ist durch $r = \mathrm{lcm}(p-1,q-1)$ bestimmt ---
eine individuelle arithmetische Eigenschaft von $N$, nicht ein
geometrischer Parameter. Shors Algorithmus nutzt diese Periodenstruktur;
sie ist eine Eigenschaft der Primzahlen, nicht der FFGFT-Geometrie.

Die Tests in Dok.~075 und Dok.~076 verwenden Rasterparameter $\sigma$
mit Werten wie $1/42$, $1/100$, $1/1000$ für Periodensuchverfahren.
Diese Größen sind numerische Abtastweiten, keine physikalischen
Parameter; sie dürfen nicht mit $\xi = 4/30000$ gleichgesetzt werden.
Das Resonanzbild ist prinzipiell richtig --- Periodenfindung ist ein
spektrales Problem ---, aber die Generatoren der Periode sind die
Primzahlen, nicht ein einstellbarer Geometrieparameter.

Ein Vorteil gegenüber klassischen Algorithmen muss auf der
Shor-Faktorisierungskurve sichtbar werden: Der Fehlerkorrektur-
Schwellenwert muss für relevante Schlüssellängen (2048 Bit) gehalten
werden; die Aufbereitung muss in Superposition ausgeführt werden; ein
RSA-relevantes $N$ muss erfolgreich faktorisiert werden. Nichts davon
ist Stand 2026 realisiert.\status{X}

\section{Grenze 2: Weyl-Obstruktion und Zähldichte}

Der Weyl-Satz bestimmt, wie viele Eigenwerte ein kompakter Zustandsraum
unterhalb einer Schranke $\lambda$ haben kann:
$N(\lambda)\sim C\cdot\lambda^{d/2}$
--- ein Potenzgesetz in der Dimension $d$.\status{K}

Die Primzahlfunktion $\pi(T)$ und die Riemannschen Nullstellen wachsen
dagegen wie $T\ln T$ --- keine Potenz, sondern quasi-logarithmisch. Das
bedeutet: Ein endlicher kompakter Zustandsraum (ein $n$-Qubit-Register
ist ein solcher) kann ein solches Spektrum nicht vollständig tragen.
Die Weyl-Obstruktion schließt drei geometrische Zugänge zur
Riemann-Hypothese aus:\status{K} keine kompakte Geometrie kann
$T\ln T$-Wachstum aus einem Laplace-Spektrum liefern; $\xi$ verschiebt
die kritische Linie nicht; die $\xi$-Leiter diskretisiert die
Nullstellen nicht.

Shors Schaltkreismodell ist digital und von der Weyl-Grenze direkt nicht
betroffen. Aber ein physikalischer Quantencomputer ist ein Resonanzsystem;
solange das Schwellentheorem für relevante Skalierungen nicht demonstriert
ist, liegt er innerhalb der Resonanzgrenzen.\status{S}

\section{Grenze 3: Der Auflösungsboden}

Die $\xi$-Korrektur zu Bell-Korrelationen liegt bei $\Delta\mathrm{CHSH}
\approx\xi/(2\pi)\approx2\times10^{-5}$ pro Messung.\status{S} Das ist
der Auflösungsboden: Unterhalb dieser Schwelle ist keine FFGFT-spezifische
Abweichung von der Standard-QM messbar. Alle NISQ-Prozessoren (NISQ =
Noisy Intermediate-Scale Quantum) operieren mit physikalischen Fehlerraten
$\sim10^{-2}$, drei Größenordnungen über dem $\xi$-Signal.

Das bedeutet für den Quantenrechner: Die $\xi$-Korrekturen zu Algorithmen
--- Bell-Dämpfung in Shor, modifizierte Gatterwinkel durch $K_{\mathrm{frak}}$
--- sind unter NISQ-Rauschen begraben.\status{K} Weder eine Verlangsamung
noch eine Beschleunigung von Quantenalgorithmen durch $\xi$ ist mit
aktueller Hardware messbar. Das ist kein Misserfolg, sondern eine
Präzisionsaussage.

Zukünftige Hochpräzisions-Bell-Tests (schlupflochfrei, Detektions-
effizienz $>99\,\%$, optische Plattform) könnten die Grenze erreichen.
Das photonische System (Kapitel~\ref{ch:hardware}) ist der natürliche
Kandidat.

\section{Was bleibt}

Von den vier Parallelstrukturen zwischen FFGFT und Quantenalgorithmen
bleiben nach dieser Grenzziehung:

Die Übersetzbarkeit (Kapitel~\ref{ch:hilbert}) ist vollständig und
symmetrisch. Jede QM-Rechnung ist in FFGFT formulierbar; jede FFGFT-Aussage
über Quantensysteme ist in der QM abbildbar. Das ist keine Einschränkung,
sondern eine Stärke: Praktiker können ihren vertrauten Formalismus
behalten.

Der Mehrwert der FFGFT liegt nicht in anderen algorithmischen Ergebnissen,
sondern in der Herkunft der Eingabewerte: Massen, Kopplungen und
$\alpha$ aus derselben Geometrie, die den Quantenformalismus trägt. Das
ist die strukturelle Asymmetrie, die Kapitel~\ref{ch:hilbert} festgehalten
hat.

\quelle{Dok.~075 (RSA-Verfahren), Dok.~076 (RSA-Periodentests), Dok.~147 §8 (IBM-Validierung), Dok.~316 (Riemann, Weyl-Grenze), Dok.~343 (Auflösungsboden)}


% --- Teil IV ---
\part{Zustandsraum}
% QMB_n13_hilbertraum.tex — Kapitel 13: Hilbertraum-Übersetzung
% Inhaltliche Grundlage: Dok. 230 (Mai 2026), Dok. 231; eigene Darstellung
\chapter{Zwei Sprachen, ein Inhalt: FFGFT und Hilbertraum}
\label{ch:hilbert}

\section{Die Behauptung}

Die FFGFT beschreibt Quantensysteme in der Sprache von Moden auf einem
Torus. Die Standard-Quantenmechanik beschreibt sie in der Sprache von
Vektoren in einem Hilbertraum. Dieses Kapitel zeigt, dass beide Sprachen
denselben Inhalt tragen: Jede Rechnung, die in der einen möglich ist,
ist in der anderen möglich, mit demselben Ergebnis. Die Übersetzung ist
keine Näherung und keine Spezialisierung, sondern eine bijektive
Zuordnung von Zuständen, Operatoren, Gattern, Produkträumen und
Messvorschriften.\status{B}

Das ist methodisch die wichtigste Aussage dieses Buches. Sie bedeutet,
dass die FFGFT die Quantenmechanik nicht ersetzt, sondern erweitert:
Beide Formalismen liefern bis auf eine $\xi$-Korrektur der Größenordnung
$10^{-5}$ dieselben messbaren Vorhersagen. Sie unterscheiden sich in der
Deutung einiger Grundbegriffe --- Born-Regel, Nichtlokalität, der Status
von $\hbar$, die Rolle der Raumzeit, Singularitäten. Wer in
Hilbertraum-Sprache zu Hause ist, kann die FFGFT benutzen, ohne seine
Werkzeuge zu wechseln. Die Brücke ist die Übersetzungstabelle.

\section{Wo wir uns geometrisch befinden}

Bevor die Tabellen kommen, ein Bild der Geometrie, in der die Übersetzung
stattfindet. Wir werden mit Zylinderkoordinaten rechnen, und der Zylinder
ist nicht die volle FFGFT-Geometrie, sondern eine bequeme Reduktion.

\paragraph{Die volle Geometrie ist vierdimensional.} Der Raum, auf dem
alle Aussagen der FFGFT leben, ist der Torus $T^4 = T^3\times S^1$: drei
räumliche Wicklungsrichtungen und eine zeitliche. Die Wellenfunktion
einer Mode $\varphi_{n,l,j}$ ist eine Funktion auf diesem Torus. Aus
seiner Geometrie folgen die Massen, die Feinstrukturkonstante, die
kosmologischen Größen --- alles, was die FFGFT von der reinen
Quantenmechanik unterscheidet.

\paragraph{Ein einzelnes Qubit braucht weniger.} Ein Zwei-Niveau-System
zu fester Zeit, ohne Massendynamik, aktiviert die meisten dieser
Wicklungsrichtungen nicht. Zwei Drehfreiheitsgrade bleiben: einer, der
kodiert, wo zwischen $|0\rangle$ und $|1\rangle$ wir uns befinden (das
wird $z$), und einer für die Phase (das wird $\theta$). Der natürliche
Raum dafür ist ein 2-Torus $T^2 = S^1\times S^1$.

\paragraph{Vom 2-Torus zum Zylinder.} Ist der Hauptradius $R$ des
2-Torus groß gegen den Schlauchradius, sieht die Geometrie lokal wie ein
Zylinder $\mathbb{R}\times S^1$ aus. Eine der beiden Schleifen ist
aufgeschnitten: Aus zwei Kreisen wird eine Gerade und ein Kreis. In
Zylinderkoordinaten $(z,r,\theta)$ ist $z$ deshalb eine lineare Achse
von $-1$ bis $+1$, nicht mehr zyklisch.

\paragraph{Vom Zylinder zur Bloch-Sphäre.} Zieht man den Zylinder oben
und unten zu je einem Punkt zusammen, entsteht eine Sphäre $S^2$: die
Bloch-Sphäre der Standard-Quantenmechanik. Sie ist die zweite,
weitergehende Reduktion.

\begin{center}
\small\renewcommand{\arraystretch}{1.3}
\begin{tabular}{p{2.9cm}cp{3.6cm}}
\toprule
\textbf{Geometrie} & \textbf{Dim.} & \textbf{Verwendung} \\
\midrule
4-Torus $T^4$ & 4 & volle FFGFT: Massen, $\alpha$, Kosmologie \\
3-Torus $T^3$ & 3 & Modenraum bei fester Zeit \\
2-Torus $T^2$ & 2 & Einzelqubit ohne Reduktion \\
Zylinder $\mathbb{R}\times S^1$ & 2 & Limes $R\to\infty$, lokale Näherung \\
Bloch-Sphäre $S^2$ & 2 & Pole zusammengezogen; Standard-QM \\
\bottomrule
\end{tabular}
\end{center}

\paragraph{Was die Reduktionen kosten.} Jede Stufe verliert etwas. Vom
4-Torus zum 3-Torus geht die Zeitwicklung verloren und damit jede
Aussage über Massen. Vom 3-Torus zum 2-Torus geht eine räumliche
Wicklung verloren und damit die Verbindung zur dritten Leptongeneration
und zu vielen Spinstrukturen. Vom 2-Torus zum Zylinder geht eine
kompakte Schleife verloren --- das ist die schwächste Reduktion, und
genau hier sitzt der fraktale Korrekturfaktor
$K_{\mathrm{frak}} = 1-100\xi\approx0{,}9867$, der in den
Bit-Flip-Drehungen auftaucht. $K_{\mathrm{frak}}$ ist die Erinnerung des
Zylinders an seine Torusherkunft.

Die Übersetzungstabellen dieses Kapitels gelten für den Qubit-Sektor,
also für die unteren drei Zeilen. Die Standard-QM arbeitet auf der
Bloch-Sphäre, die FFGFT auf dem Zylinder oder dem 2-Torus. Die
Übersetzung zwischen diesen Stufen ist verlustfrei, sofern man die
kleinen geometrischen Korrekturen ($K_{\mathrm{frak}}$, $\Delta$CHSH)
mitführt. Der 4-Torus liegt eine Stufe höher und liefert, was der
Hilbertraum nicht liefern kann: die Werte der Eingabeparameter.

\begin{laierspur}
Man kann eine Landkarte auf verschiedene Weisen zeichnen: als Globus,
als abgerollten Zylinder, als flaches Blatt. Alle zeigen dieselben
Städte an denselben Stellen. Aber nur der Globus zeigt, dass man von
Osten aus wieder im Westen ankommt. Die Bloch-Sphäre der
Quantenmechanik ist das flache Blatt; der Torus der FFGFT ist der Globus.
Für die Navigation zwischen zwei Städten reicht das Blatt. Für die
Frage, warum die Welt so groß ist, wie sie ist, braucht man den Globus.
\end{laierspur}

\section{Zustände}

In der Standard-QM ist ein Qubit ein Vektor in $\mathbb{C}^2$:
\begin{equation}
  |\psi\rangle = \alpha|0\rangle + \beta|1\rangle,\qquad
  |\alpha|^2+|\beta|^2 = 1,\qquad \alpha,\beta\in\mathbb{C}.
\end{equation}
Mit globaler Phasenfreiheit ist das ein Punkt auf der Bloch-Sphäre
(Hopf-Faserung $S^3\to S^2$).

In der FFGFT ist dasselbe Qubit ein Punkt $(z,r,\theta)$ auf dem
Zylinder:
\begin{equation}
  z\in[-1,1],\quad r\in[0,1],\quad \theta\in[0,2\pi),\qquad z^2+r^2 = 1.
\end{equation}
$z$ ist die Projektion auf die Berechnungsachse, $r$ die radiale
Superpositionsamplitude, $\theta$ die Phase.

Die Brücke zwischen beiden Darstellungen:
\begin{equation}
  \boxed{\;
  \alpha = \sqrt{\tfrac{1+z}{2}}\;e^{\,i\theta/2},\qquad
  \beta = \sqrt{\tfrac{1-z}{2}}\;e^{-i\theta/2}.\;}
  \label{eq:state-bridge}
\end{equation}
Umgekehrt:
\begin{equation}
  z = |\alpha|^2-|\beta|^2,\qquad
  r = 2|\alpha\beta|,\qquad
  \theta = \arg\alpha-\arg\beta.
\end{equation}
Die Konsistenz prüft man in vier Zeilen: Die Normierung
$|\alpha|^2+|\beta|^2 = \frac{1+z}{2}+\frac{1-z}{2} = 1$ stimmt; die
Zylinderbedingung $z^2+r^2 = (|\alpha|^2-|\beta|^2)^2+4|\alpha|^2|\beta|^2 = 1$
stimmt; $|0\rangle$ entspricht $(1,0,\ast)$ und $|1\rangle$ entspricht
$(-1,0,\ast)$, mit beliebiger Phase an den Polen.

Der begriffliche Unterschied liegt nicht in den Zahlen. In der QM ist
$|\alpha|^2$ eine Wahrscheinlichkeit. In der FFGFT ist $r^2 = 1-z^2$
eine Energiedichte der radialen Konfiguration. Beide Größen sind bei
großer Statistik numerisch gleich; konzeptuell ist die FFGFT ein
deterministisches geometrisches Modell, dessen statistische Vorhersagen
mit der Born-Regel zusammenfallen.

\section{Operatoren}

Die drei Pauli-Matrizen erzeugen alle Ein-Qubit-Operationen:
\begin{equation}
  \sigma_x = \begin{pmatrix}0&1\\1&0\end{pmatrix},\quad
  \sigma_y = \begin{pmatrix}0&-i\\i&0\end{pmatrix},\quad
  \sigma_z = \begin{pmatrix}1&0\\0&-1\end{pmatrix},
\end{equation}
mit der Algebra $\sigma_i\sigma_j = \delta_{ij}\mathbb{1}+i\epsilon_{ijk}\sigma_k$.

Auf dem Zylinder entsprechen ihnen drei geometrische Drehungen:
$\sigma_z$ ist die Rotation um die Zylinderachse, also die Phasenrotation
$\theta\to\theta+\pi$; $\sigma_x$ ist die Drehung in der $(z,r)$-Ebene um
den Winkel $\pi$, der Bit-Flip; $\sigma_y$ ist die dazu orthogonale
Drehung, die Bit-Flip und Phasenrotation kombiniert.

\begin{center}
\small\renewcommand{\arraystretch}{1.5}
\begin{tabular}{cp{2.6cm}p{3.6cm}}
\toprule
\textbf{Operator} & \textbf{Hilbertraum} & \textbf{FFGFT-Geometrie} \\
\midrule
$\sigma_z$ & $\mathrm{diag}(1,-1)$ & Phasenrotation $\theta\to\theta+\pi$ \\
$\sigma_x$ & $\begin{pmatrix}0&1\\1&0\end{pmatrix}$ & $(z,r)$-Drehung um $\pi K_{\mathrm{frak}}$ \\
$\sigma_y$ & $\begin{pmatrix}0&-i\\i&0\end{pmatrix}$ & Drehung mit Phasenkippung \\
$H$ & $\frac{1}{\sqrt2}\begin{pmatrix}1&1\\1&-1\end{pmatrix}$ & $z\leftrightarrow r$ vertauschen, $\theta\to\theta+\pi/2$ \\
$T$ & $\mathrm{diag}(1,e^{i\pi/4})$ & Phasenrotation $\theta\to\theta+\pi/4$ \\
\bottomrule
\end{tabular}
\end{center}

Die einzige FFGFT-spezifische Abweichung ist der Faktor
$K_{\mathrm{frak}}\approx0{,}9867$ in der $\sigma_x$-Drehung. Er folgt aus
der fraktalen Raumzeitdimension $D_f = 2{,}999867$ und ist eine prüfbare
Vorhersage: Alle $\pi$-Gatter weichen systematisch um
$\Delta\alpha\approx0{,}013\,\%$ vom idealen Wert ab.\status{S} Im
Hilbertraum-Formalismus wäre das eine Ad-hoc-Korrektur, die als
Hardware-Rauschen eingebaut werden müsste; in der FFGFT folgt sie aus
der Geometrie.

\section{Gatter}

Jedes unitäre Ein-Qubit-Gatter $U\in U(2)$ ist eine Bloch-Drehung:
\begin{equation}
  U = e^{i\alpha}R_{\hat n}(\phi)
    = e^{i\alpha}\bigl(\cos\tfrac{\phi}{2}\,\mathbb{1}
      - i\sin\tfrac{\phi}{2}\,\hat n\cdot\vec\sigma\bigr).
\end{equation}
In der FFGFT ist das die Rotation des Punktes $(z,r,\theta)$ um die Achse
$\hat n$ um den Winkel $\phi$. Die Übersetzung ist exakt.

Das CNOT-Gatter
\begin{equation}
  \mathrm{CNOT} =
  \begin{pmatrix}1&0&0&0\\0&1&0&0\\0&0&0&1\\0&0&1&0\end{pmatrix}
\end{equation}
übersetzt sich als Vorschrift: Wenn $z_A = -1$ (Kontrollqubit auf
$|1\rangle$), führe den Bit-Flip auf Qubit B aus. Das ist eine geometrische
Konditionierung --- die Drehung des einen Zylinderpunkts hängt vom
$z$-Wert des anderen ab. Topologisch ist das eine Verkettung der beiden
Zylinder, und das ist die Verschränkung.

Die Menge $\{H,T,\mathrm{CNOT}\}$ ist nach dem Solovay-Kitaev-Theorem
universell für die Quantenmechanik.\status{E} Da jedes dieser Gatter in der
FFGFT als geometrische Operation darstellbar ist, ist die FFGFT
vollständig universell für Quantenberechnung: Es gibt keine QM-Aussage,
die nicht in FFGFT-Sprache formuliert werden kann.\status{B}

\section{Produkträume und Verschränkung}

Für $n$ Qubits ist der Hilbertraum
$\mathcal{H}_n = \bigotimes_{i=1}^n\mathbb{C}^2$ mit Dimension $2^n$,
und ein allgemeiner Zustand lautet
\begin{equation}
  |\Psi\rangle = \sum_{x\in\{0,1\}^n}\alpha_x|x\rangle,\qquad
  \sum_x|\alpha_x|^2 = 1.
\end{equation}
In der FFGFT entsprechen $n$ Qubits $n$ Modenkonfigurationen
$(z_i,r_i,\theta_i)$ auf einer gemeinsamen Torusgeometrie. Verschränkung
ist eine topologische Verbindung: Verschränkte Qubits sind Knoten der
Modenkonfiguration, die nicht auf Produktform reduziert werden können.
Der Bell-Zustand $|\Phi^+\rangle = \frac{1}{\sqrt2}(|00\rangle+|11\rangle)$
ist eine topologisch verbundene Zylinderkonfiguration, in der die
einzelnen $(z_i,r_i,\theta_i)$ nicht unabhängig sind.

Beide Formalismen reproduzieren das exponentielle Wachstum $\dim = 2^n$:
im Hilbertraum durch das Tensorprodukt, in der FFGFT durch $n$
unabhängige Modenkonfigurationen mit topologischen Verbindungen. Die
Zahl der reellen Parameter pro Zustand ist in beiden Fällen $2\cdot2^n-2$.

Die CHSH-Vorhersage: Die Standard-QM gibt $|\langle\mathrm{CHSH}\rangle|
\le2\sqrt2$. Die FFGFT sagt eine kleine, geometrisch motivierte Abweichung
voraus, $\Delta\mathrm{CHSH}\approx\xi/(2\pi)\approx10^{-5}$ pro
Messung.\status{S} Sie liegt unter der aktuellen Hardware-Rauschschwelle;
die Unterscheidung von der kumulativen Größe aus Dok.~022 ist in
Kapitel~\ref{ch:chsh_exp} ausgeführt.

\section{Messung}

Eine Messung in der Berechnungsbasis liefert in der QM den Wert $0$ mit
Wahrscheinlichkeit $|\alpha|^2$ und den Wert $1$ mit Wahrscheinlichkeit
$|\beta|^2$; der Zustand kollabiert auf $|0\rangle$ oder $|1\rangle$.

In der FFGFT ist Messung kein probabilistischer Kollaps, sondern eine
geometrische Wechselwirkung, die den Punkt $(z,r,\theta)$ auf das
nächstgelegene stabile Extrem $z = \pm1$ treibt. Das ist eine
deterministische Bewegung im Modenraum --- wegen der unbekannten Phase
$\theta$ und kleinster Variationen der Anfangsbedingungen statistisch
aber ununterscheidbar vom Verhalten der Born-Regel. Bei großer Statistik
konvergiert die Messverteilung:
\begin{equation}
  P_{\mathrm{FFGFT}}(z = +1) = \tfrac{1+z}{2} = |\alpha|^2 = P_{\mathrm{QM}}(0).
\end{equation}
Die Formalismen liefern identische statistische Vorhersagen.\status{B}
Der Streit über lokalen Realismus löst sich in der FFGFT geometrisch,
weil die Verschränkung als topologische Verbindung und nicht als
nichtlokale Korrelation dargestellt wird (Kapitel~\ref{ch:bell_ffgft}).

\section{Zeitentwicklung}

Die Zeitentwicklung der QM ist die Schrödinger-Gleichung
$i\hbar\,\partial_t|\psi\rangle = H|\psi\rangle$. In der FFGFT ist sie die
deterministische Bewegung der Modenkonfiguration $\varphi_{n,l,j}$ auf
$T^3\times\mathbb{R}$. Die Modengleichung reduziert sich im
Niederenergielimes $E\ll E_\xi$ auf die Schrödinger-Gleichung:
\begin{equation}
  i\hbar\,\frac{\partial\varphi_{n,l,j}}{\partial t}
  \;\xrightarrow{\;E\ll E_\xi\;}\;
  -\frac{\hbar^2}{2m}\nabla^2\varphi_{n,l,j}+V\varphi_{n,l,j}.
\end{equation}
Bei höheren Energien treten $\xi$-Korrekturen auf, die im
Hilbertraum-Formalismus als Modifikationen des Hamiltonoperators
erscheinen.

\section{Was übersetzbar ist und was nicht}

Alle Hilbertraum-Aussagen sind in FFGFT abbildbar: Zustände, Operatoren
und Gatter (mit oder ohne $K_{\mathrm{frak}}$), Tensorprodukte,
Verschränkung, Bell- und GHZ-Zustände, Messung, unitäre Evolution und
alle Algorithmen. Umgekehrt entspricht jede $(z,r,\theta)$-Konfiguration
eindeutig einem $(\alpha,\beta)$-Vektor, jede geometrische Drehung einem
unitären Operator, jede topologische Verbindung einem verschränkten
Vektor. Jede Berechnung, die in einem Formalismus durchführbar ist, ist
im anderen durchführbar --- mit identischem Ergebnis.

Die Übersetzbarkeit gilt für Aussagen über Quantensysteme. Sie gilt
nicht für Aussagen über die Eingabewerte der Quantenmechanik:
\begin{itemize}
  \item Der Parameter $\xi = 4/30000$ ist im Hilbertraum ein externer
        Eingabewert. In der FFGFT folgt er aus der Geometrie
        (Kapitel~\ref{ch:grundlage}).
  \item Die Fermionmassen erscheinen im Hilbertraum als Yukawa-Kopplungen.
        In der FFGFT folgen sie aus $m_i = r_i\,\xi^{p_i}\,v$.
  \item Die Feinstrukturkonstante ist in der FFGFT ableitbar, im
        Hilbertraum eingegeben.
  \item $K_{\mathrm{frak}}$ erscheint in FFGFT-Gattern als
        Geometrieabweichung; im Hilbertraum müsste er als Rauschen
        modelliert werden.
\end{itemize}
Das ist die strukturelle Asymmetrie: Der Hilbertraum hat alle Werkzeuge
für QM-Berechnungen, aber keine Erklärung für die Eingabewerte. Die FFGFT
liefert die Werkzeuge und die Eingabewerte aus einer einzigen Geometrie.
In diesem Sinn ist die Quantenmechanik eine Untermenge der FFGFT: die
FFGFT, eingeschränkt auf die Frage, wie sich Systeme verhalten, ohne die
Frage, warum die Konstanten diese Werte haben.

Die FFGFT erhebt keinen Anspruch, die einzige korrekte Darstellung zu
sein. Sie ist eine Darstellung, die zusätzlich die Werte der
fundamentalen Konstanten liefert. Wer diesen Mehrwert nicht braucht,
kann sie ignorieren; die Standard-QM bleibt gültig.

\section{Operationale Äquivalenz, interpretative Divergenz}

Die vollständige Übersetzbarkeit darf nicht als ontologische Identität
gelesen werden. Operational stimmen FFGFT und Standard-QM in jedem
messbaren Ergebnis bis auf $\xi$-Korrekturen überein. Interpretativ
weichen sie an fünf Stellen ab, an denen die Standardlesart Aussagen
trifft, die in der FFGFT als Artefakte der gewählten Trägerannahmen
gelesen werden:

\begin{enumerate}
  \item \textbf{Born-Regel als irreduzibler Zufall.} Standardlesart:
        $|\psi|^2$ ist ontologische Wahrscheinlichkeit, der Messausgang
        fundamental zufällig. FFGFT: $|\psi|^2$ ist die korrekte Verteilung
        der Beobachtungen, entsteht aber aus einer deterministischen
        Polübergangsdynamik. Die Stochastizität sitzt beim Messapparat,
        nicht im Naturprozess.
  \item \textbf{Nichtlokalität als Eigenschaft der Natur.} Standardlesart
        der Bell-Korrelationen: Fernwirkung im $\mathbb{R}^3$. FFGFT:
        topologische Verbindung auf $T^4$, geometrische Nähe in einer
        geschlossenen Geometrie.
  \item \textbf{$\hbar$ als Primitive.} Standardlesart: fundamentale
        Naturkonstante. FFGFT: SI-Konversionsfaktor zwischen Energie- und
        Zeitmaß, der aus der Abschlussbedingung $S = h$ auf $T^4$ folgt.
        Numerisch identisch, ontologisch abgeleitet.
  \item \textbf{Raumzeit als Bühne.} Standardlesart der Relativitäts-
        theorie: gegebene Mannigfaltigkeit. FFGFT: aus der
        $T^4$-Identifikationsgeometrie abgeleitet. Fragen wie
        \enquote{was war vor dem Urknall} sind Artefakte der
        $\mathbb{R}^3\!+\!t$-Annahme.
  \item \textbf{Singularitäten als Realitäten.} Standardlesart: Schwarzes
        Loch und Urknall als echte Singularitäten. FFGFT: Da $T^4$
        geschlossen ist, sind sie Endpunkte des Trägermodells, nicht der
        Physik. Außerhalb der Singularitätsregion stimmen die Vorhersagen
        überein.
\end{enumerate}

Diese fünf Punkte sind keine empirischen Behauptungen gegen die
Standard-QM. Sie sind Interpretationsdifferenzen, die sichtbar werden,
sobald nach dem ontologischen Träger gefragt wird. Wer nur mit den Zahlen
rechnet, kann sie ignorieren. Wer wissen will, warum die Zahlen gerade
diese sind, für den werden sie zu prüfbaren strukturellen Aussagen.

Innerhalb der heutigen Messpräzision --- NISQ-Niveau $\sim10^{-2}$,
optische Präzisionsexperimente $\sim10^{-12}$ bei Bell-Tests --- liefern
beide dieselben Vorhersagen. Die $\xi$-Korrektur $\sim10^{-5}$ ist mit
gegenwärtiger Hardware nicht auflösbar; das ist auf IBM Heron~r2
empirisch geprüft (Kapitel~\ref{ch:chsh_exp}).

\section{Feldtheoretische Erweiterung}

Die Übersetzung lässt sich vom Qubit-Sektor auf Quantenfeldzustände
ausdehnen. Der Fock-Raum der Standard-QFT --- die direkte Summe aller
$n$-Teilchen-Räume --- entspricht dem Modenspektrum des Trägers, wobei
die Besetzungszahlen als Wicklungszahlen erscheinen. Erzeugungs- und
Vernichtungsoperatoren sind Übergänge zwischen benachbarten
Wicklungszuständen derselben Mode; Kommutatoren folgen aus der
Periodizität. Das erzeugt keine neuen Vorhersagen im perturbativen
Bereich, strukturiert aber den Aufbau des Standardmodells geometrisch um:
Die Teilchensorten sind Modenklassen des Gitters, nicht unabhängige
Felder mit je eigener Kopplung.\status{S}

\begin{keypoint}[Was dieses Kapitel festhält]
Zustände, Operatoren, Gatter, Verschränkung, Messung und Zeitentwicklung
sind zwischen FFGFT und Hilbertraum bijektiv übersetzbar. Die einzige
formale Abweichung ist $K_{\mathrm{frak}}$ in den $\pi$-Gattern; die
einzige statistische ist $\Delta$CHSH. Der Mehrwert der FFGFT liegt nicht
in anderen Rechenergebnissen, sondern darin, dass die Eingabewerte der
Quantenmechanik aus derselben Geometrie folgen.
\end{keypoint}

\quelle{Dok.~230 (Übersetzbarkeit, Mai 2026), Dok.~231 (feldtheoretische Erweiterung), Dok.~147 §1.2, Dok.~175 §3--4, Dok.~202 §22}

% QMB_n14_gitter_spektral.tex — Kapitel 14: Gitter im Hilbertraum, Spektraltheorie
% Inhaltliche Grundlage: Dok. 314 (Gitter im Hilbertraum), Dok. 322 (Spektraltheorie), Dok. 327; eigene Darstellung
\chapter{Das Gitter im Hilbertraum}
\label{ch:gitter}

\section{Werden die Gitter flach?}

Die FFGFT rechnet geometrisch auf $T^4$ mit dem $D_4$-Gitter und
quantenmechanisch auf $\mathcal{H} = L^2(T^4)\otimes\mathbb{C}^3$. Beim
Übergang stellt sich eine naheliegende Frage: Werden die Gitter dabei
flach? Bleibt von Kusszahl~24, Packungsdichte und Trialität etwas
übrig, oder ist im Hilbertraum jedes Gitter dasselbe?

Die Antwort: Der Torus ist ohnehin flach --- die Krümmung war nie der
Unterschied ---, und die Gitterstruktur geht nicht verloren, sondern
\textbf{wandert ins Spektrum und in die Faser}. Im Hilbertraum wird das
Gitter zur Indexmenge der Fouriermoden. Damit ist es keine Kugelpackung
mehr; in diesem Sinn wird jedes Gitter flach. Aber die Information
verschwindet nicht: Die Entartungen des Laplace-Operators sind die
Thetareihe des Gitters, die Kusszahl wird zur Entartung der ersten
Anregung, die Trialität wird zur Orbifold- und Faserstruktur.
\begin{equation}
  \boxed{\;\text{Geometrie}\;\longrightarrow\;\text{Spektrum}\;+\;\text{Faser}\;}
\end{equation}

Zwei Anker halten diese Übersetzung stabil. Erstens ist $D_4$ bis auf
Skalierung selbstdual: Das duale Gitter $D_4^*$ ist wieder ein
$D_4$-Gitter, die Fouriertransformation führt nicht aus der Familie
heraus. Zweitens ändert die Faser am Gitter nichts: Die
$\mathbb{C}^3$-Faser hängt an jedem Punkt und trägt die
$\mathbb{Z}_3$-Struktur, das Modengitter bleibt das duale Gitter.

\textbf{Zur Notation.} In der Gittertheorie ist $\Lambda$ das übliche
Zeichen für ein Gitter. Im FFGFT-Korpus ist $\Lambda$ belegt (kosmologische
Konstante, Abschlussskala). Dieses Kapitel verwendet deshalb $\Gamma$
für das Gitter und $\Gamma^*$ für das duale Gitter; die Kantenlängen des
Torus heißen $R_{123}$ (drei Raumrichtungen) und $R_4$ (vierte Richtung),
ihr Verhältnis $r = R_4/R_{123}$.

\section{Das Spektrum trennt, was die Krümmung nicht trennt}

Die vollständige Schalenzählung bis $|k|^2 = 12$, für das kubische
Gitter $\mathbb{Z}^4$ und für $D_4$ (alle ganzzahligen $k$ mit gerader
Koordinatensumme):

\begin{center}
\small
\begin{tabular}{rrrc}
\toprule
Norm $|k|^2$ & $\mathbb{Z}^4$ & $D_4$ & $D_4$-Split ($k_4{=}0$ / $k_4{\ne}0$) \\
\midrule
1 & 8 & 0 & --- \\
2 & 24 & 24 & 12 / 12 \\
3 & 32 & 0 & --- \\
4 & 24 & 24 & 6 / 18 \\
5 & 48 & 0 & --- \\
6 & 96 & 96 & 24 / 72 \\
7 & 64 & 0 & --- \\
8 & 24 & 24 & 12 / 12 \\
9 & 104 & 0 & --- \\
10 & 144 & 144 & 24 / 120 \\
11 & 96 & 0 & --- \\
12 & 96 & 96 & 8 / 88 \\
\bottomrule
\end{tabular}
\end{center}

Drei Ablesungen. \textbf{Zwei flache Tori sind spektral
unterscheidbar:} $\mathbb{Z}^4$ beginnt bei Norm~1 mit Entartung~8,
$D_4$ bei Norm~2 mit Entartung~24. Bei $D_4$ fehlen die ungeraden Normen
vollständig --- die Paritätsbedingung schließt sie aus. Was die
Krümmung nicht unterscheidet, unterscheidet das Spektrum sofort.\status{K}
\textbf{Unabhängige Kontrolle:} Die $D_4$-Reihe $24,24,96,24,144,96,\dots$
ist die bekannte Thetareihe $\theta_{D_4} = (\theta_3^4+\theta_4^4)/2$;
die Skriptrechnung reproduziert sie koeffizientenweise. \textbf{Die
$12+12$-Zerlegung ist spektral adressierbar:} Die erste Anregung trägt
exakt zwölf Moden mit $k_4 = 0$ (die räumliche FCC-Hälfte) und zwölf mit
$k_4\ne0$ (die Wicklungshälfte). Bei höheren Schalen ist die Aufspaltung
asymmetrisch --- $6/18$ bei Norm~4, $8/88$ bei Norm~12. Die
$12/12$-Symmetrie ist eine Eigenschaft der ersten Schale, kein
Allgemeinsatz; wo der Korpus mit $24 = 12+12$ argumentiert, gilt das genau
dort, wo es gemacht wird.

\section{Anisotropie: wo $\tilde T\cdot m = 1$ im Spektrum ablesbar wird}

Wird die vierte Richtung mit anderem Radius kompaktifiziert, lautet der
Eigenwert $E(k) = k_1^2+k_2^2+k_3^2+(k_4/r)^2$, und die erste Schale
spaltet auf:

\begin{center}
\small
\begin{tabular}{rll}
\toprule
$r$ & $E$ (Wicklung, 12 Moden) & $E$ (FCC, 12 Moden) \\
\midrule
$1{,}0$ & $2{,}000$ & $2{,}000$ (entartet: 24) \\
$1{,}1$ & $1{,}826$ & $2{,}000$ \\
$1{,}5$ & $1{,}444$ & $2{,}000$ \\
$2{,}0$ & $1{,}250$ & $2{,}000$ \\
\bottomrule
\end{tabular}
\end{center}

Die FCC-Moden bleiben bei $E = 2$ fest, die Wicklungsmoden folgen exakt
$E_{\mathrm{Wicklung}} = 1+1/r^2$. Die Aufspaltung ist ein direktes Maß
für das Radienverhältnis --- die Stelle, an der die Zeit-Masse-Struktur
im Spektrum ablesbar wird: Massenrichtung und Raumrichtungen tragen
verschiedene Anregungsenergien, sobald sie verschieden lang sind.\status{K}

Alle Niveaus haben die Form $E = a+b/r^2$ mit ganzzahligen $a,b$. Im
Bereich $1<r\le2{,}5$ gibt es unter allen Niveaus bis Norm~8 exakt zwei
Kreuzungsradien: Bei $r = \sqrt2$ kreuzt die reine Wicklung
$(0,0,0,\pm2)$ die FCC-Zwölf, bei $r = \sqrt3$ die Wicklungs-Zwölf. Die
Niveaufolge ändert sich also nicht kontinuierlich, sondern an zwei
ausgezeichneten Wurzelradien; jede Ordnungsaussage über die niedrigsten
Anregungen zerfällt in genau drei Regime. Jenseits von $\sqrt3$ ist der
Grundton des Spektrums eine reine Wicklung.

\section{Die Reziprozität ist modenweise exakt --- als Ortsrelation}

Die vierte Richtung ist die zeit-masse-duale; ausgerollt wird auf Zeit.
Zwei Prüffragen: Wie entstehen Zeitperioden, und gilt $\tilde T\cdot m = 1$
nur im Mittel?

Das Ausrollen der kompakten vierten Richtung ist die universelle
Überlagerung $\mathbb{R}\to S^1$; die Decktransformationsgruppe ist
$\mathbb{Z}$, die abgerollte Koordinate zählt Zyklusdurchläufe. Eine
$k_4$-Mode wird dabei zur periodischen Funktion mit Periode
$2\pi R_4/k_4$. Da $\pi_1(S^1) = \mathbb{Z}$ abelsch ist, ist diese
Buchhaltung verlustfrei.

Im Deformationsspektrum ist $E(k) = k_1^2+k_2^2+k_3^2+(k_4/r)^2$ wörtlich
$E^2 = p^2_{\mathrm{Raum}}+m^2$ mit $m_k = k_4/R_4$: \textbf{Die
Wicklungszahl ist die Massenzahl.} Damit gilt
\begin{equation}
  \lambda_4\cdot m = \frac{2\pi R_4}{k_4}\cdot\frac{k_4}{R_4} = 2\pi
  \qquad\text{exakt, für jede Mode.}
\end{equation}
Die Reziprozität gilt nicht nur im Mittel --- sie ist modenweise
identisch erfüllt, weil Periode und Masse aus derselben Quantenzahl und
demselben Radius gebaut sind.\status{K}

\textbf{Aber die Periode ist zunächst ein Ort, keine Dauer.} Die Größe
$\lambda_4 = 2\pi R_4/k_4$ ist eine Bogenlänge auf dem Kreis: der
Abstand in $x_4$, nach dem sich die Mode wiederholt. Die Relation
$\lambda_4\cdot m = 2\pi$ ist die de-Broglie-Relation auf der
Massenrichtung --- Wellenlänge mal Wellenzahl ---, keine Zeitaussage.
Dauer entsteht erst aus der ausgerollten Koordinate: Ein Punkt der
ausgerollten Zeit ist ein Paar $t = w\cdot2\pi R_4+x_4$ aus Umlaufzahl
$w\in\mathbb{Z}$ und Ort auf dem Zyklus. Die Dauer steckt vollständig
im Zähler $w$, der Ort im Rest. Auf dem Zyklus gibt es keinen
ausgezeichneten Anfang; jeder Moment ist eine Stelle, die wiederkehrt.

\textbf{Die zeitliche Ablesung trägt einen $\gamma$-Faktor.} Fragt man
nach der Schwingungsdauer in der ausgerollten Zeit, ist
$T_{\mathrm{osc}} = 2\pi/E$ mit $E = \sqrt{p^2+m^2}$ zuständig, und dann
\begin{equation}
  T_{\mathrm{osc}}\cdot m = \frac{2\pi}{\gamma}\le2\pi,
  \qquad\text{Gleichheit genau bei } p_{\mathrm{Raum}} = 0.
\end{equation}
Exakt $2\pi$ ist die Ortsablesung; die Zeitablesung fällt nur für die
ruhende Mode mit ihr zusammen.

\textbf{Im Gemisch wird die Gleichung zur Ungleichung.} Überlagert ein
Zustand mehrere $k_4$-Schalen mit Gewichten $p_k$, gilt
$\langle T\rangle\langle m\rangle = 2\pi\langle k\rangle\langle1/k\rangle$
und nach Jensen $\langle k\rangle\langle1/k\rangle\ge1$, mit Gleichheit
genau dann, wenn eine einzige Schale besetzt ist. Numerisch: scharfe
Schale exakt $1$; $k = 1,2$ gleichbesetzt $1{,}125$; Gleichverteilung
$k = 1\dots50$ liefert $2{,}295$; thermische Verteilungen zwischen
$1{,}00002$ und $2{,}6$; zweitausend Zufallsverteilungen nie unter~1.
Die Vermutung dreht sich um: Die Bedingung für $\tilde T\cdot m = 1$ ist
nicht Gleichverteilung, sondern \textbf{Schärfe} --- genau eine besetzte
Massenschale. Der Exzess misst die Massenstreuung des Zustands.

Damit stehen drei Aussagen in drei Richtungen: Die Ortsablesung ist
immer $2\pi$; die Zeitablesung einer bewegten Mode ist $2\pi/\gamma$;
das Massengemisch liegt bei $\ge2\pi$. Die exakte Relation gilt für die
ruhende, scharfe Mode --- dort und nur dort verschwinden beide
Korrekturfaktoren zugleich. Bewegung drückt das Produkt herunter,
Massenstreuung hebt es an. Das ist kein Widerspruch zur nativen
FFGFT-Formulierung, sondern die Kaluza-Klein-Ablesung ihrer
Gültigkeitsbedingungen.

\begin{laierspur}
Eine Gitarrensaite hat eine Grundfrequenz. Spannt man sie und zupft
sie an, schwingt sie mit genau dieser Frequenz --- die Relation
\enquote{Länge mal Frequenz} ist exakt. Aber wenn man mehrere Saiten
zugleich anschlägt, oder wenn die Gitarre an einem vorbeifährt, hört
man etwas anderes: einen Akkord, oder einen Dopplerton. Die FFGFT-Relation
$\tilde T\cdot m = 1$ ist die Saite in Ruhe; alles andere sind
Korrekturen, die man genau angeben kann.
\end{laierspur}

\section{Eingerollt und ausgerollt: die Schließungsgabelung ist spektral unsichtbar}

Der Korpus stellt die Schließungsfrage als Gabelung: Der Defekt
$d = 100\xi = 1/75$ je Umlauf lässt drei Fälle für die ausgerollte Zeit
zu. Mit der Zyklenlänge $\tau_c = 2\pi/H_0 = 91{,}9$\,Gyr: Fall~A ohne
Defekt schließt nach einem Umlauf; Fall~B mit eingefrorenem $\xi$ nach
75 Umläufen (rationale Rotationszahl $74/75$, teilerfremd), also nach
$6892$\,Gyr; Fall~C mit laufendem $\xi$ nie --- die Windung wird zur
logarithmischen Spirale.

Diese Zahlen sind Fristen, keine Wiederholungen. Sie sagen nicht, dass
sich etwas wiederholt, sondern innerhalb welcher Frist eine Wiederkehr
stattfinden müsste, wenn der jeweilige Fall zuträfe. Geschlossen ist die
Geometrie; ob die Geschichte zurückkehrt, ist davon unabhängig. Der
\enquote{Urknall} ist in dieser Lesart ein Ort auf dem Zyklus --- die
Region kleinster Massen und langsamster Uhren ---, kein Ereignis; ein
geschlossener Kreis hat keinen Anfangspunkt.

Naheliegend wäre, den Defekt als Twist der Randbedingung zu lesen
($k_4\to k_4+1/75$). Das ist falsch, und der Grund ist strukturell:
\textbf{Im eingerollten Zustand gibt es keine fraktalen Korrekturen.}
Die Eindeutigkeit der Mode auf dem Kreis erzwingt $k_4\in\mathbb{Z}$
--- ganzzahlig oder gar nicht. Eine Wicklungszahl ist topologisch; sie
hat keinen Platz für $1/75$. Korrekturen brauchen etwas, das sich
akkumuliert, und akkumulieren kann nur ein durchlaufener Weg. Schon die
Formulierung \enquote{Defekt je Umlauf} setzt das Zählen von Umläufen
voraus, und Umläufe zählt man nur in der ausgerollten Koordinate.

\begin{center}
\small
\begin{tabular}{p{4.2cm}ll}
\toprule
Größe & Zustand & fraktale Korrektur \\
\midrule
Thetareihe, Entartungen, $24 = 12+12$ & eingerollt & nein \\
$\mathrm{Aut}(D_4)$, Trialität, Orbifold & eingerollt & nein \\
Zirkulant, Container $4\times6$ & eingerollt & nein \\
Casimir, Wärmekern, Kongruenzen & eingerollt & nein \\
$\lambda_4\cdot m = 2\pi$ & eingerollt & nein \\
\midrule
Zyklenlänge, Defekt, Fall A/B/C & ausgerollt & ja \\
Umlaufzahl $w$, Geschichte & ausgerollt & ja \\
$T_{\mathrm{osc}}$, $\gamma$-Faktor & ausgerollt & ja \\
\bottomrule
\end{tabular}
\end{center}

Zwei Konsequenzen. Erstens hat die Gabelung keine spektrale Signatur:
Das Modenspektrum ist ein Objekt des eingerollten Zustands; ob die Bahn
nach einem, nach 75 oder nach keinem Umlauf schließt, ändert an den
Eigenwerten nichts. Zweitens sind alle folgenden Abschnitte
fallunabhängig: Orbifold, Zirkulant und Container sind Aussagen über die
kompakte Geometrie, und dort greift der Defekt nicht.\status{K}

\section{Umorganisationen: drei Ebenen und eine Grenze}

Flache Gitter lassen sich auf drei Ebenen umorganisieren, die sehr
Verschiedenes kosten. \textbf{Basiswechsel} (kostenlos): Jede
unimodulare Transformation aus $GL(4,\mathbb{Z})$ beschreibt dasselbe
Punktgitter mit anderen Erzeugern; das Spektrum ändert sich um nichts.
\textbf{Gittersymmetrien} (kostenlos, strukturtragend): Umorganisationen,
die das Gitter auf sich abbilden --- hier sitzt der zentrale Befund des
nächsten Abschnitts. \textbf{Umbau zu einem anderen Gitter} (kostet
Physik): diskret über Klebevektoren entlang der Kette
$D_4\subset\mathbb{Z}^4\subset D_4^*$ mit jeweils Index~2, oder
kontinuierlich durch Deformation im Modulraum. Jede Stufe verschiebt
Entartungen; $D_4$ ist darin ein ausgezeichneter Punkt --- Kusszahl 24,
in vier Dimensionen nachweislich nicht überbietbar.

Die Grenze der Spektralaussage: Ab Dimension~4 existieren Paare flacher
Tori, die isospektral, aber nicht isometrisch sind (Schiemann). Der
Allgemeinsatz \enquote{Spektrum bestimmt das Gitter} ist in 4D falsch.
Für $D_4$ selbst unkritisch --- die Kusszahl 24 zeichnet es eindeutig
aus ---, aber als Regel darf es nicht gebucht werden.

\section{Die Trialität ist eingebaut, nicht angeheftet}

Am Gitter nachgerechnet, durch Brute-Force-Zählung aller Basistupel mit
identischer Gram-Matrix:
\begin{center}
\small
\begin{tabular}{lr}
\toprule
$|\mathrm{Aut}(D_4)|$ & $1152 = |W(F_4)|$ \\
$|\mathrm{Aut}(\mathbb{Z}^4)|$ & $384 = 2^4\cdot4!$ \\
$|\mathrm{Aut}(D_4)|/|W(D_4)|$ & $1152/192 = 6 = |S_3|$ \\
\bottomrule
\end{tabular}
\end{center}
Der Quotient $6 = |S_3|$ weist die Trialität direkt an der
Gitter-Automorphismengruppe nach, nicht nur am Dynkin-Diagramm: $D_4$ ist
das einzige Gitter der $D_n$-Familie mit drei gleichberechtigten
Dynkin-Armen; bei $D_5,D_6,\dots$ schrumpft die äußere Symmetrie auf
$\mathbb{Z}_2$.\status{K}

Die 80 Elemente der Ordnung~3 zerfallen in drei Klassen. Gerechnet wird
exakt: Die Trialitätselemente sind halbzahlig --- die Klebevektormischung
$(\frac12,\frac12,\frac12,\frac12)$ geht ein; Rundung auf ganze Zahlen
zerstört sie, deshalb rechnet man mit $2A$.
\begin{center}
\small
\begin{tabular}{lrp{4.5cm}}
\toprule
Klasse & Anzahl & Charakter \\
\midrule
Spur $-2$, halbzahlig & 16 & fixpunktfrei in $\mathbb{R}^4$; Eigenwerte $\{\omega,\bar\omega\}$ doppelt \\
Spur $1$, ganzzahlig & 32 & Fixebene, frei auf den 24 Minimalvektoren \\
Spur $1$, halbzahlig & 32 & 6 Fixvektoren $+$ 6 Dreierorbits \\
\bottomrule
\end{tabular}
\end{center}

\textbf{Der Orbifold ist Gittersymmetrie, keine Zusatzannahme.} Für
jedes der 16 Elemente mit Spur $-2$ gilt exakt $|\det(1-A)| = 9$; der
Quotient $T^4/\langle A\rangle$ ist ein $\mathbb{Z}_3$-Orbifold mit genau
neun Fixpunkten, alle in Drittelkoordinaten. Die $T^4/\mathbb{Z}_3$-Struktur
muss nicht gesetzt werden; sie ist in $D_4$ bereits vorhanden.\status{K}

\textbf{Auf jedem Modenorbit wirkt die $\mathbb{Z}_3$ als Zirkulant.}
Die unitäre Wirkung auf $L^2(T^4)$ permutiert die drei ebenen Wellen
eines Dreierorbits zyklisch; die Permutationsmatrix hat Eigenwerte
$\{1,\omega,\omega^2\}$, Phasen $0^\circ$, $+120^\circ$, $-120^\circ$
--- strukturell exakt die Zirkulantdiagonalisierung der
$\mathbb{C}^3$-Faser. Die erste Schale liefert acht Kopien:
$24 = 8\times3 = 8\times(\chi_0\oplus\chi_1\oplus\chi_2)$.

\section{Der Negativbefund: 5 teilt 1152 nicht}

Die Ordnungsverteilung in $\mathrm{Aut}(D_4)$, exakt bestimmt:
\begin{center}
\small
\begin{tabular}{lrrrrrrr}
\toprule
Ordnung & 1 & 2 & 3 & 4 & 6 & 8 & 12 \\
\midrule
Anzahl & 1 & 139 & 80 & 228 & 464 & 144 & 96 \\
\bottomrule
\end{tabular}
\end{center}
Summe 1152. \textbf{Ordnung 5: null Elemente} --- kein Rechenergebnis,
sondern ein Satz: $1152 = 2^7\cdot3^2$, und 5 teilt das nicht
(Lagrange).\status{K}

Konsequenz für die Koide-Phase $\theta = 2/9$: Die Fünffachdrehung, aus
der $p_0 = 2/9$ folgt, ist mit der $D_4$-Gitterstruktur prinzipiell
unverträglich. Der Test wurde zusätzlich direkt geführt: Alle
Invarianten der linearen Gitterwirkung an den neun Orbifold-Fixpunkten
haben das Nennerspektrum $\{1,2,3,6\}$ --- keine Neuntel. Das entwertet
die Herleitung von $2/9$ nicht, sondern grenzt sie ab: Die Phase gehört
einer Schicht \emph{über} dem Gitter (der $A_5$-Struktur), das Gitter
trägt die $\mathbb{Z}_3$. Die Arbeitsteilung ist damit exakt benannt.

\section{Die Verdopplung ist dennoch vorbereitet}

Widerspricht der Negativbefund der Verdopplung $6 = 3\oplus3'$? Nein,
in drei Stufen.

\textbf{Das Verbot trifft die falsche Ebene.} $5\nmid1152$ verbietet
Ordnung-5-Gitterautomorphismen. Die Verdopplung lebt auf der Faser
($\mathbb{C}^3\to\mathbb{C}^6$); $A_5$ wird nirgends als Raumsymmetrie
gefordert. Die Kompatibilitätsbedingung ist allein, dass die
Schnittstelle beider Ebenen existiert: Elementordnungen in $A_5$ sind
$\{1,2,3,5\}$, in $\mathrm{Aut}(D_4)$ sind sie $\{1,2,3,4,6,8,12\}$; der
Schnitt ist $\{1,2,3\}$, kanonisch also $C_3$. Genau die stellt die
Trialität bereit. Die Schnittstelle ist nicht nur erlaubt, sie ist die
einzig mögliche --- und sie ist besetzt.

\textbf{Das Gitter bereitet die Verdopplung vor.} Die zentrale
Involution $-1\in\mathrm{Aut}(D_4)$ kommutiert mit der Trialitäts-$\mathbb{Z}_3$
und paart die acht Dreierorbits der ersten Schale in vier disjunkte
Paare. Kein Orbit ist selbstgepaart, beweisbar: $-k = A^jk$ erforderte
den Eigenwert $-1$ für ein Element der Ordnung~3, unmöglich. Also
\begin{equation}
  24 = 8\times3 = 4\times6\qquad(\text{Orbit}\oplus\text{Spiegelorbit}).
\end{equation}
Jeder Sechserblock trägt unter $C_3$ den Charakterinhalt
$2\times(\chi_0\oplus\chi_1\oplus\chi_2)$ --- exakt die Restriktion von
$3\oplus3'$. Die erste Anregungsschale liefert vier fertige Container.\status{K}

\textbf{Die Grenze, sauber gebucht.} Die Etikettierung --- welcher Faktor
$3$ und welcher $3'$ ist --- ist gitterblind. Die $A_5$-Charaktere der
beiden Tripletts unterscheiden sich nur auf den 5er-Klassen ($\varphi$
gegen $1-\varphi$); auf der 3er-Klasse sind beide null. Die Unterscheidung
hängt exakt an der Struktur, die im Gitter nicht existiert. Bilanz: Das
Gitter liefert Container und $C_3$-Inhalt, die $A_5$-Schicht liefert die
Etiketten, $\varphi$ und die Fünffach --- und damit $2/9$.

\section{Ankerlose Ergebnisse}

Alle folgenden Größen sind dimensionslos --- Verhältnisse, Zählungen,
Kongruenzen. Keine SI-Skala geht ein.

\textbf{Zwei Kongruenzsätze.} Alle vollständigen Schalen bis Norm~120
geprüft: $N(n)\equiv0\pmod3$ und $N(n)\equiv0\pmod6$ ausnahmslos. Beides
ist Satz, nicht Empirie: Die Trialitäts-$\mathbb{Z}_3$ ist auf
$\Gamma\setminus\{0\}$ fixpunktfrei, also zerfällt jede Schale in
Dreierorbits; zusammen mit der $-1$-Paarung ohne Selbstpaarung in
Sechserblöcke.\status{K}

\textbf{Charakteraufgelöstes Spektrum.} Folge daraus: exakte
Gleichverteilung auf jeder Schale. Der invariante Sektor des Orbifolds
hat exakt $N(n)/3$ Zustände pro Schale; die $\mathbb{Z}_3$-Projektion
kostet auf keiner Schale Asymmetrie. Ehrliche Konsequenz:
Gleichverteilung und Sechsercontainer sind keine Auszeichnung der ersten
Schale, sondern Satz für alle. \textbf{Die Struktur ist überall --- und
kann darum nichts selektieren.} Das Muster ist skaleninvariant präsent;
welche Schale physikalisch besetzt ist, entscheidet etwas anderes.

\textbf{Casimir: das dichteste Gitter verliert.} Die
$\zeta$-regularisierte Vakuumenergie eines skalaren Feldes, beide Tori
auf Kovolumen~1 normiert: $\mathbb{Z}^4$-Torus $E = -0{,}9321$,
$D_4$-Torus $E = -0{,}8686$. Der $\mathbb{Z}^4$-Torus liegt tiefer. Der
Grund ist die Spektrallücke: Das dichteste Gitter hat bei gleichem
Volumen die größte erste Eigenfrequenz ($\sqrt2$ gegen 1) und damit die
weniger negative Nullpunktsenergie. Bemerkenswert die Ordnungsumkehr:
Für $s>0$ minimiert $D_4$ die Epstein-Zeta, bei $s = 0$ sind alle Gitter
gleich, bei $s = -1/2$ ist die Ordnung invertiert. Interpretationsgrenze:
Das gilt für das skalare, periodische Feld; \enquote{die Natur wählt
$\mathbb{Z}^4$} folgt daraus nicht.

\textbf{Wärmekern: die Beobachtbarkeitsgrenze beziffert.} Für beide Tori
gilt $t^2\theta(t)\to1$ (identischer Weyl-Term bei gleichem Volumen); der
Gitterunterschied ist exponentiell klein: relative Differenz
$1{,}2\times10^{-2}$ bei $t = 0{,}5$, $1{,}8\times10^{-13}$ bei
$t = 0{,}1$, $4\times10^{-27}$ bei $t = 0{,}05$. Jede Observable, die nur
auf den führenden Wärmekern-Koeffizienten beruht, ist gitterblind.

\section{Was hier noch nackt ist --- und was nicht}

Dieses Kapitel rechnet mit dem glatten Laplace-Operator. Zu klären ist,
was die FFGFT an fraktaler Struktur dagegenhält, und dabei sind zwei
Größen sauber zu trennen: $D_f^{\mathrm{Raum}} = 3-\xi$ ist lokal und
wirkt mit $6{,}7\times10^{-5}$; $K_{\mathrm{frak}} = 1-100\xi$ ist
kumulativ über den RG-Lauf von der Planck- zur elektroschwachen Skala
(17 Größenordnungen) und wirkt mit $1{,}33\times10^{-2}$. Das
Verhältnis der beiden Effekte ist exakt 200. Die fraktale Dimension der
FFGFT ist die räumliche, nicht eine vierdimensionale; eine Größe
$D_f = 4-\xi$ kennt der Korpus nicht.

Ein Laplace-Operator auf dem kompakten Torus ist lokal. Eine kumulative
Größe kann er nicht tragen; $K_{\mathrm{frak}}$ kommt nicht über den
Operator herein, sondern über die Brückenkonstanten bei der
Skalenumrechnung. Die Konsequenz zerfällt in drei ungleiche Teile.

\textbf{Unberührt: alles Ankerlose.} $K_{\mathrm{frak}}$ greift bei
absoluten Größen und kürzt sich bei Verhältnissen derselben Ebene. Damit
bleiben unberührt: alle Zählungen ($24$, $8$, $N(n)$, $8\times3$,
$4\times6$), die Kongruenzen, die Charakteraufteilung, alle Verhältnisse
$E/E_{\min}$ und die Kreuzungsradien. Diese Aussagen sind
fraktal-invariant.

\textbf{Ein Faktor: alle absoluten Energien.} Mit einem Anker wird
$E = \hbar H_0|k|$. Diese Werte tragen den Faktor: erste $D_4$-Schale
$3{,}230\times10^{-52}$\,J nackt, $3{,}187\times10^{-52}$\,J mit
$K_{\mathrm{frak}}$; durchgehend $-1{,}33\,\%$. Der Faktor sitzt nicht im
Spektrum, sondern im Anker: Die dimensionslosen Eigenwerte bleiben
unverändert, korrigiert wird die Umrechnung in SI.

\textbf{Der Operator, nicht die Umrechnung: die Störungsrechnung.} Bei
$D_f^{\mathrm{Raum}} = 3-\xi$ ist der Operator kein exakt glatter
Laplace mehr, und exakte Symmetrie ist genau das, was die Entartungen
trägt. Maßgeblich ist allein der lokale Effekt, $\Delta E/E\approx
6{,}7\times10^{-5}$ --- zwei Größenordnungen unter dem, was eine
Abschätzung über $K_{\mathrm{frak}}$ ergäbe. Die 24-fache Entartung
wäre dann ein Multiplett mit Feinstruktur. In welche Untermultipletts es
zerfällt, entscheidet die Symmetrie der Störung, und zwar über die
Darstellungstheorie, nicht die Orbitstruktur:

\begin{center}
\small
\begin{tabular}{lp{3.6cm}}
\toprule
Störung respektiert & Multiplettstruktur der 24 \\
\midrule
volle $\mathrm{Aut}(D_4)$ & $9+8+4+2+1$ \\
Trialität und $-1$ & $8\times1+8\times2$ \\
nur $\mathbb{Z}_3$ & $8\times1+8\times2$ \\
nur $-1$ & $24\times1$ \\
nichts davon & $24\times1$ \\
\bottomrule
\end{tabular}
\end{center}

Die naheliegende Erwartung ist falsch. Man vermutet, die Orbits blieben
zusammen. Tatsächlich folgen die Multiplettgrößen den irreduziblen
Darstellungen der respektierten Symmetriegruppe: Jede ortsabhängige
Störung bricht die Translationsinvarianz, und die trug die
24er-Entartung, nicht die Punktsymmetrie; $\mathbb{Z}_3$ ist abelsch, alle
Irreps eindimensional, die Dubletts sind die antiunitäre Paarung der
konjugierten Charaktere; $-1$ allein erzwingt gar nichts; selbst der
maximal symmetrische Störer hält die 24 nicht zusammen. Qualitativ ist
von der Entartung nur so viel stabil, wie die Störungssymmetrie an
Irrep-Dimension hergibt; quantitativ liegt die Störung bei $10^{-5}$,
eine Feinstruktur weit unterhalb jeder Ablesbarkeit, und die
Containerstruktur bleibt praktisch intakt.\status{K}

Bilanz: Nackt sind nur die SI-Absolutwerte, und zwar über den Anker,
nicht über den Operator. Entartungen, Verhältnisse, Kongruenzen,
Container, Casimir- und Wärmekern-Verhältnisse sind nicht nackt, sondern
korrekt: $K_{\mathrm{frak}}$ greift dort strukturell nicht.

\section{Der Spektraloperator $\hat F$}

Was bisher am Gitter gerechnet wurde, erhält seine rigorose
Hilbertraum-Einbettung durch den Spektraloperator. In der linearen
Algebra hat eine symmetrische Matrix reelle Eigenwerte und eine
Orthonormalbasis aus Eigenvektoren. In unendlichdimensionalen
Hilberträumen verallgemeinert die Spektraltheorie das: Für einen
selbstadjungierten Operator existiert ein eindeutiges spektrales Maß mit
$T = \int\lambda\,dE(\lambda)$, und die Selbstadjungiertheit garantiert
ein reelles Spektrum --- die Bedingung dafür, dass Observablen messbare
Werte haben.\status{E}

Die FFGFT setzt einen fraktalen Operator
\begin{equation}
  \hat F\equiv\bigoplus_{n=1}^\infty S_n\circ L_n
\end{equation}
mit Skalenoperatoren $S_n = r_nU_n$ (Kontraktionsraten $r_n\in(0,1)$,
unitäre Phasendrehungen $U_n$) und logischen Projektoren $L_n$ mit
$L_nL_m = L_{\min(n,m)}$ --- die fraktale Inklusionsstruktur: Projektionen
auf feinere Skalen enthalten die auf gröbere.\status{SETZUNG} Da die
Fundamentallänge $L_0 = \xi\ell_P$ die Skalenleiter endlich macht, ist
$\sum r_n<\infty$; $\hat F$ ist beschränkt auf ganz $\mathcal{H}$, die
Symmetrie folgt aus der $\mathbb{Z}_3$-Paarung $(k,-k)$, die Defektindizes
sind $(0,0)$: Es existiert genau eine selbstadjungierte
Realisierung.\status{B} Damit ist die in Dok.~282 offen geführte Frage
(Registerbrücke R82) geschlossen; auf dem kompakten Träger gilt:
beschränkt und symmetrisch impliziert selbstadjungiert.

Die Eigenwertgleichung $\hat F\psi_\lambda = \lambda\psi_\lambda$
definiert die Resonanzmoden. Der Geometrieparameter ist kein freier
Parameter, sondern ein Spektralwert:
\begin{equation}
  \xi = \lambda_{\min}(\hat F_{D_4}) = \frac{4}{30000},
\end{equation}
der kleinste Eigenwert des $D_4$-Suboperators --- die fundamentale
Frequenz des fraktalen Raumes. Das erklärt strukturell, warum $\xi$ als
Basis in allen Massenformeln erscheint: Es ist die unterste Mode des
Geometrieoperators.\status{K}

Der Zustandsraum hat Tensorproduktstruktur
$\mathcal{H} = \mathcal{H}_{\mathrm{geom}}\otimes\mathcal{H}_{\mathrm{spin}}
\otimes\mathcal{H}_{\mathrm{flavor}}$ mit
$\mathcal{H}_{\mathrm{geom}} = L^2(T^4/\mathbb{Z}_3,d\mu_f)$,
$\mathcal{H}_{\mathrm{spin}} = \mathbb{C}^6$ und
$\mathcal{H}_{\mathrm{flavor}} = \ell^2(\mathbb{N}^2)$ mit Wicklungszahlen
$(n_\theta,n_\phi)$ als Quantenzahlen. Das fraktale Maß $d\mu_f$ entsteht
durch sukzessive Fraktalisierung des Lebesgue-Maßes in 100 Schritten
mit laufendem $\xi^{(k)} = \xi(1-k/100)$; nach der Rekursion sind die
Koeffizienten $n^\xi\approx1+\xi\ln n$ nahezu 1, und die fraktale
Dimension $D_f = 3-\xi$ emergiert aus dem Rekursionsprozess, nicht aus
einer Setzung.

Die $\mathbb{Z}_3$-Orbifold-Projektion $P_{\mathbb{Z}_3} = \frac13\sum_k
\omega^k\tau^k$ schneidet den invarianten Unterraum aus $L^2(T^4)$ aus.
An den Fixpunkten gilt die Randbedingung
$\psi_\chi(x_0+y) = \chi\cdot\psi_\chi(x_0+g_*(y))$ mit dem Charakter
$\chi$ und der linearisierten Wirkung $g_*$. Die drei Neutrinoflavours
sind an drei dieser Fixpunkte lokalisiert; die unterschiedlichen
Positionen erzeugen unterschiedliche Eigenwerte des
Massenoperators.\status{K}

\begin{keypoint}[Was dieses Kapitel festhält]
Im Hilbertraum wird das Gitter flach, aber nicht folgenlos: Was die
Geometrie trug, tragen Spektrum und Faser. Die Trialität ist
Gittersymmetrie; der $\mathbb{Z}_3$-Orbifold mit neun Fixpunkten steckt
in $D_4$; Fünffachsymmetrie ist am Gitter unmöglich und gehört der
$A_5$-Faserschicht; $\tilde T\cdot m = 1$ ist modenweise exakt als
Ortsrelation und gilt zeitlich nur für die ruhende, scharfe Mode.
Eingerollte Größen tragen keine fraktale Korrektur; $K_{\mathrm{frak}}$
wirkt allein über die Skalenbrücke. $\hat F$ ist selbstadjungiert, $\xi$
sein kleinster $D_4$-Eigenwert.
\end{keypoint}

\quelle{Dok.~314 (Gitter im Hilbertraum, 2026), Dok.~322 (Spektraltheorie), Dok.~327 (Selbstadjungiertheit, R86), Dok.~133 ($K_{\mathrm{frak}}$), Dok.~285/291/293/311/313}

% QMB_n15_zeit_superposition.tex — Kapitel 15: Zeit im Zustandsraum, Superposition ohne Zeit
% Inhaltliche Grundlage: Dok. 307 (Juli 2026), Dok. 334 (Aug. 2026), Dok. 333; eigene Darstellung
\chapter{Wo die Zeit sitzt: Zustandsraum und Superposition}
\label{ch:zeit}

\section{Die Gabelung}

Am Anfang jeder quantenmechanischen Modellbildung steht eine Wahl, die
meist unausgesprochen bleibt und dennoch alles Weitere festlegt: Ist
die Zeit eine Richtung \emph{im} Zustandsraum, oder läuft sie
\emph{daneben} her, als äußerer Parameter?

Die Frage klingt technisch, ist aber ontologisch. Sie entscheidet nicht,
welche Mathematik verwendet wird --- beide Antworten arbeiten mit
selbstadjungierten Operatoren, Spektren und Projektoren ---, sondern
\emph{wo} die Zeit sitzt. Und genau diese Verortung entscheidet, ob die
physikalischen Größen aus der Struktur \emph{fallen} oder ihr
\emph{angeheftet} werden müssen.

Die Frage reicht zu den Wurzeln der Quantentheorie zurück. Schon 1932
erkannte Wolfgang Pauli, dass ein Zeitoperator, der kanonisch zum
Hamiltonoperator konjugiert ist, auf grundsätzliche Schwierigkeiten
stößt. Das Pauli-Theorem hat die Entwicklung der Quantenmechanik
geprägt und die Zeit als äußeren Parameter scheinbar unumstößlich
gemacht. Spätere Arbeiten --- Aharonov und Bohm, Garrison und Wong ---
zeigten, dass unter bestimmten Bedingungen, etwa bei nicht
halbbeschränkten Hamiltonoperatoren oder kompakten Konfigurationen,
Zeitoperatoren doch konstruierbar sind. Diese Ausnahme ist für die FFGFT
der Regelfall.

\section{Die ausgerollte Sicht: Zeit als äußerer Parameter}

Die vertraute Formulierung siedelt die Zeit außerhalb des Zustandsraums
an. Der Hilbertraum ist der Raum der Zustände \emph{zu einem
Zeitpunkt}, typischerweise $\mathcal{H} = L^2(\text{Konfigurationsraum})$.
Die Zeit ist der Parameter $t$, entlang dessen sich der Zustand
entwickelt, vermittelt durch $e^{-i\hat Ht}$. Ort, Impuls, Spin sind
Operatoren auf $\mathcal{H}$. Die Zeit ist keiner.

Dass sie keiner sein kann, ist kein Versehen, sondern ein Resultat:
Paulis Einwand zeigt, dass zu einem nach unten beschränkten
Hamiltonoperator kein selbstadjungierter, kanonisch konjugierter
Zeitoperator existiert.\status{E} In der ausgerollten Sicht bleibt die
Zeit deshalb notwendig die Bühne, nicht das Stück.

Die Konsequenz ist strukturell: Da die Zeit außen steht, enthält der
Zustandsraum allein keine dynamischen Skalen. Massen, Kopplungen,
Lebensdauern sind in der reinen Raumstruktur nicht vorhanden --- sie
müssen von außen eingeführt werden. Im Standardmodell geschieht das
durch Yukawa-Terme und die freien Parameter, die deren Stärke
festlegen. Die Struktur liefert den Rahmen; die Zahlen kommen aus der
Messung.

Das Problem ist alt. Weyl und Wigner erkannten, dass die Massen der
Elementarteilchen gruppentheoretisch als Casimir-Invarianten erscheinen
--- als strukturelle Label, nicht als angeheftete Parameter. Dennoch
blieb die Herkunft der Massenwerte offen. Der Higgs-Mechanismus erklärt
die Masse der Eichbosonen, führt aber selbst neue freie Parameter ein.
Die Kritik an der ausgerollten Sicht ist nicht, dass sie falsch wäre,
sondern dass sie die Herkunft der Massenskalen unbeantwortet lässt und
auf externe Eingaben verweist.

\subsection*{Eine repräsentative Konstruktion dieser Art}

Viele Modellprogramme, die eine geometrische Herkunft der
Teilcheneigenschaften anstreben, folgen bei aller Verschiedenheit
demselben Schema: Man wählt einen diskreten Träger (ein Gitter, eine
Cut-and-Project-Struktur, einen periodischen Graphen); definiert darauf
eine selbstadjungierte Operatorfamilie; extrahiert über
Spektralprojektoren Sektoren; und ordnet diesen Sektoren nachträglich
physikalische Bedeutung zu --- Ladung, Spin, Masse, Letztere meist über
eine parametrische Abbildung mit ein oder zwei freien Parametern.

Ein solches Vorgehen kann mathematisch vollständig kontrolliert sein.
Die entscheidende Beobachtung betrifft nicht seine Strenge, sondern
seine Verortung der Zeit: Der Träger ist ein räumliches Objekt, die
Operatorevolution läuft über einen äußeren Parameter. Die Sektoren sind
räumliche Eigenräume, keine Zeitmoden. Die Masse ist in ihnen nicht
enthalten; sie wird angeheftet.

Und hier greift ein Argument, das sich selbst trägt: \textbf{Eine
parametrische Abbildung mit freien Parametern kann jede endliche
Werteliste treffen.} Ist die Abbildung nicht vorab und universell
fixiert, besitzt sie keinen prädiktiven Gehalt --- sie benennt die
bekannten Eingaben um. Um Zahlen zu liefern, muss eine solche
Konstruktion an externe Werte gekoppelt werden: an die Messung, und
damit indirekt an das Standardmodell selbst. Das ist redlich, sofern es
seine Grenze deklariert: dass es keine Massen ableitet, sondern eine
Hypothesenarchitektur bereitstellt. Die Schwierigkeit dieses Wegs ist
kein Zufall, sondern die Folge der ausgerollten Verortung der Zeit.

Ein konkretes Beispiel: Die Graphen-Modellierung definiert
Quasiteilchen auf hexagonalen Gittern. Die Anregungen zeigen lineare
Dispersion und lassen sich als masselose Fermionen lesen --- aber die
Einführung einer Masse erfordert stets zusätzliche Terme (Haldane-Term,
Kekulé-Distorsion) mit freien Parametern. Die Geometrie liefert die
Kinematik, nicht die Massenskala.

\section{Die eingerollte Sicht: Zeit als kompakte Dimension}

Die Alternative belässt die Zeit nicht außen und eliminiert sie auch
nicht, sondern zieht sie \emph{in} den Zustandsraum --- als geschlossene,
periodische Dimension. An die Stelle von $L^2(\text{Raum})$ mit äußerem
$t$ tritt ein Raum, in dem eine kompakte Zeitrichtung Teil der Geometrie
ist: $\mathcal{H} = L^2(T^4)\otimes\mathbb{C}^3$, wobei eine der
Torusrichtungen die Zeit trägt. Die Bijektivität dieser Darstellung
(Typ III) ist im Korpus belegt (Kapitel~\ref{ch:hilbert}).

Auf einer geschlossenen Dimension sind die zulässigen Zustände stehende
Moden mit diskreten Windungszahlen. Die Quantisierung ist dann nicht
postuliert, sondern geometrisch erzwungen: Nur bestimmte Moden passen
auf die geschlossene Struktur. Die Masse wird zum Eigenwert eines
Operators auf $\mathcal{H}$ selbst --- nicht angeheftet, sondern
strukturell enthalten. Die Verhältnisse der Massen liegen fest, weil die
zulässigen Moden festliegen.\status{B}

Die Relation $\tilde T\cdot m = 1$ ist die konkrete Ausführung dieses
Gedankens: Jede massive Größe \emph{ist} ihr eigener Zeitzyklus. Die
Entkompaktifizierung $S^1_m\to\mathbb{R}_t$ ist maßerhaltend, also
verlustfrei. Dass eine kompakte Dimension einen diskreten Modenturm
erzeugt, dessen Massen sich wie das Inverse der Kompaktifizierungsperiode
verhalten, ist kein exotischer Zug, sondern der Mechanismus, den Kaluza
und Klein für eine kompakte Raumdimension vorgeführt haben --- hier auf
die Zeit angewandt.

Paulis Einwand trifft die eingerollte Sicht nicht. Er gilt für eine
offene, nach unten unbeschränkte Zeit konjugiert zu einem beschränkten
Hamiltonoperator. Eine kompakte Zeit verhält sich wie eine
Winkelvariable: Ihr konjugiertes Spektrum ist diskret, wie beim
Drehimpuls zum Winkel, und das Pauli-Argument entfällt.\status{B} Die
geschlossene Zeit ist nicht nur zulässig, sie ist der natürliche Ort
einer diskreten Modenstruktur.

Diese Einsicht hat Wurzeln: Witten wies darauf hin, dass kompakte
Zeitdimensionen in der Stringtheorie auftreten; Hawking diskutierte die
Implikationen geschlossener Zeit für die Kosmologie; Bars entwickelte
die zweizeitige Physik. Und die eingerollte Sicht berührt die
thermodynamische Zeit: Ist die Zeit kompakt, ist die Pfeilrichtung nicht
fundamental, sondern eine emergente Eigenschaft der nichtkompakten
Raumanteile --- eine natürliche Erklärung makroskopischer
Irreversibilität ohne Annahmen über Anfangsbedingungen.

\begin{laierspur}
Man kann eine Uhr auf zwei Weisen bauen. Entweder man hängt ein
Pendel an einen Rahmen und lässt die Zeit von außen laufen --- dann
muss man die Länge des Pendels von Hand einstellen. Oder man baut die
Uhr als geschlossenen Kreis, in dem jede Schwingung wieder in sich
zurückläuft --- dann bestimmt der Kreis selbst, welche Schwingungen
möglich sind. Die FFGFT baut die Uhr auf die zweite Art. Deshalb muss
sie die Massen nicht einstellen: Sie fallen aus der Geometrie.
\end{laierspur}

\section{Der Kern der Abwägung}

\textbf{Was die ausgerollte Sicht gewinnt:} den nahtlosen Anschluss an
die Maschinerie der Quantenfeldtheorie; die Freiheit von einer starken
ontologischen Festlegung; die Vertrautheit der offenen Zeitachse.
\textbf{Was sie kostet:} Die Struktur allein legt keine Massen fest.
Dynamische Werte müssen von außen kommen; Vorhersagen hängen an
externer Kalibrierung; ein Modell, das sich als eigenständig versteht,
koppelt für seine Zahlen doch wieder an die Messung.

\textbf{Was die eingerollte Sicht gewinnt:} Die Moden und ihre
Verhältnisse fallen aus der Struktur, ohne parametrisches Anheften; die
diskrete Quantisierung ist geschenkt statt postuliert. \textbf{Was sie
kostet:} eine starke ontologische Behauptung --- Zeit als geschlossene
Dimension ---, die der gewohnten Intuition der laufenden Zeit
widerspricht. Und sie hebt nicht jede äußere Anbindung auf: Die absolute
Skala braucht weiterhin eine Einheitenbrücke, die SI-Umrechnung über
feste Konversionsfaktoren.

Der entscheidende Befund ist aber dieser: \textbf{Beide Wege benutzen
denselben Operatorformalismus.} Der Unterschied liegt nicht in der
Mathematik, sondern allein darin, ob die Zeit innen oder außen sitzt.
Man kann dieselbe Konstruktion --- selbstadjungierter Operator,
Spektralsektoren --- beibehalten und nur die Zeit in den Raum ziehen;
dann werden die Sektoren zu Moden, und die Masse, die zuvor angeheftet
werden musste, wird zum Eigenwert der Struktur. Die Wahl der
Zeitverortung ist eine Rahmenwahl, die der Modellbildung vorausgeht und
deren interpretative Reichweite festlegt.

\section{Wer die Zeit schon einmal nach innen geholt hat}

Der Gedanke, die Zeit nicht als äußeren Parameter zu behandeln, ist
keine Neuheit, sondern eine lange verfolgte Linie. Es ist redlich, die
eingerollte Sicht in dieser Linie zu verorten und ihre Besonderheit
abzugrenzen.

\textbf{Relationale Zeit.} Page und Wootters lassen die Zeit innerhalb
des Zustandsraums entstehen: als Verschränkung zwischen einem
Uhr-Subsystem und dem Rest; der Gesamtzustand ist stationär,
\enquote{Evolution ohne Evolution}. Moreva et al. und Marletto und
Vedral haben den Mechanismus präzisiert und experimentell illustriert;
Gambini und Pullin entwickelten die Montevideo-Interpretation.

\textbf{Zeitlose Quantengravitation.} Wheeler--DeWitt: Die universelle
Wellenfunktion ist zeitlos, $\hat H\Psi = 0$; das daraus folgende
\enquote{Problem der Zeit} ist die Feststellung, dass es keine äußere
Zeit gibt. Kuchař und Isham lieferten die Übersichten. Rovelli:
relationale Quantenmechanik, mit Connes die thermische-Zeit-Hypothese.
Barbour: die konsequente Elimination der Zeit zugunsten zeitloser
Konfigurationen.

\textbf{Kompakte Zeit in der Stringtheorie.} Bars und Kollegen mit
zweizeitiger Physik und erweiterter Symmetriegruppe.

\textbf{Die Modenseite.} Wigner: Masse ist eine Casimir-Invariante, ein
Label irreduzibler Darstellungen der Poincaré-Gruppe --- bereits ein
struktureller Eigenwert. Kaluza und Klein: Eine kompaktifizierte
Dimension erzeugt einen diskreten Modenturm mit gequantelten Ladungen
--- die Blaupause dafür, dass eine eingerollte Dimension Moden
hervorbringt.

Die Abgrenzung ist präzise: Der überwiegende Teil dieser Tradition
macht die Zeit relational, emergent oder zeitlos --- sie nimmt die Zeit
als Fundamentales \emph{heraus}. Die eingerollte Sicht nimmt die Zeit
\emph{hinein}: als kompakte geometrische Dimension, die Moden trägt.
Sie teilt mit Page--Wootters und Rovelli die Ablehnung der äußeren
Parameterzeit, unterscheidet sich aber im positiven Vorschlag.

\section{Mathematische Ausführung}

Der Hilbertraum der eingerollten Sicht ist
$\mathcal{H} = L^2(T^1\times M)\otimes\mathcal{H}_{\mathrm{int}}$, mit
einer kompakten Zeitdimension $T^1$ der Länge $T$ und dem gewöhnlichen
Raum $M$. Eine Funktion darauf entwickelt sich in eine Fourierreihe:
\begin{equation}
  \Psi(t,\mathbf{x}) = \sum_{n\in\mathbb{Z}}e^{i n t/T}\,\psi_n(\mathbf{x}).
\end{equation}
Die diskreten Frequenzen $\omega_n = 2\pi n/T$ entsprechen den
Massenwerten in natürlichen Einheiten. Der Operator, der sie erzeugt,
ist $\hat E = -i\,\partial_t$ mit dem Spektrum $\{2\pi n/T\}$. Die
Quantisierung der Massen ist eine direkte Konsequenz der Periodizität
der Zeit; nichts muss postuliert werden.

\textbf{Yukawa-Kopplung als Auswahlregel.} Die Kopplung zwischen Moden
ist ein Matrixelement des Wechselwirkungsoperators:
\begin{equation}
  g_{ijk}\propto\int_{T^1}e^{i(n_i-n_j-n_k)t/T}\,dt
  = T\cdot\delta_{n_i,\,n_j+n_k}.
\end{equation}
Das ist eine exakte Erhaltung des Modenindex in der Zeitrichtung. Die
Yukawa-Kopplungen des Standardmodells erscheinen als Auswahlregeln für
die Modenindizes, nicht als freie Parameter.\status{B}

\textbf{Ladungsquantisierung.} Die kompakte Zeitrichtung verhält sich
wie eine fünfte Dimension im Sinne von Kaluza--Klein, aber zeitartig.
Die Ladung ist $Q\propto n/R_T$ mit dem Radius $R_T$ der kompakten Zeit
--- Ladungsquantisierung als Konsequenz der Periodizität.

\textbf{Massenhierarchie.} Der Korpus beschreibt die Leptonmassen in
drei ineinander übersetzbaren Darstellungen: erstens grundlegend über
Quantenzahlen $(n_i,l_i,j_i)$ mit $\xi_i = \xi_0\,f(n_i,l_i,j_i)$ und
$E_i = 1/\xi_i$, wobei die Absolutwerte die fraktale Korrektur
$K_{\mathrm{frak}}$ tragen; zweitens in Leiterform
$m_i = r_i\,\xi^{\pi_i}\,v$ mit rationalen $r_i$ und der
Generationshierarchie $\pi = \frac32, 1, \frac23$; drittens als
$\mathbb{Z}_3$-Zirkulant, $\sqrt{m_k} = M\bigl(1+\sqrt2\cos(\frac29+2\pi k/3)\bigr)$,
mit der Koide-Relation $Q = 2/3$ und dem Winkel $\theta = 2/9$ als
Ausgabe der Diagonalisierung. Die dimensionslosen Verhältnisse sind
exakt; Abweichungen entstehen erst beim Übergang zu absoluten
SI-Werten, aus der SI-Umrechnung, den Brückenkonstanten $v$ und $E_0$,
der fraktalen Korrektur und der endlichen Messgenauigkeit. Das ist ein
Strukturbeweis: Die Verhältnisse folgen aus der Geometrie, die
Absolutwerte aus der Brücke.\status{K}

\section{Kritische Würdigung}

\textbf{Ontologischer Status.} Die kompakte Zeit ist im Korpus real und
geometrisch, nicht bloß Rechentechnik. Der Rand $\partial T^4 = \emptyset$
macht die Frage nach einem Anfangszeitpunkt nicht falsch beantwortet,
sondern falsch gestellt. Die Zeit ist diskret: Die FFGFT besitzt ein
minimales, von $\xi$ festgelegtes Zeitinkrement $d_1 = 100\xi_0 = 1/75$,
und die zulässigen Zustände sind diskrete Moden. Wheelers
\enquote{Zeit als Illusion} teilt die eingerollte Sicht nicht: Sie
nimmt die Zeit als real, aber geometrisch anders als gewohnt.

\textbf{Die absolute Skala: fixiert, nicht frei.} Die Massenverhältnisse
folgen aus der Modenstruktur und sind exakt. Die absolute Skala ist
kein freier Parameter: Das Verhältnisgerüst ruht auf $\xi_0 = 4/30000$,
die Absolutwerte tragen ihre SI-Korrektur über feste Brückenkonstanten
--- $v$ in den Massen, $E_0 = 7{,}398$\,MeV in $\alpha = \xi E_0^2$,
feste Zahlenfaktoren in $G$. Die Planck-Größen sind kein Eingang,
sondern folgen aus $\xi$: $\xi\to\{m_e,\dots\}\to G\to\ell_P$, und
$\hbar$ selbst ist geometrisch. Die einzige exakte $\xi$-Potenzrelation
ist $L_0/\ell_P = \xi_0$, $L_0\approx2{,}15\times10^{-39}$\,m. Der
Zeitzyklus ist nicht an die Planck-Skala gekoppelt --- die Kausalrichtung
ist umgekehrt. Es bleibt kein zu eichender Modul offen; die
Dilaton-Analogie zur Stringtheorie trägt hier nicht.

\textbf{Kausalität.} Erzeugt eine periodische Zeit geschlossene
zeitartige Kurven? Der Träger ist ein Torus mit $\pi_1\neq0$ und stabilen
Windungen, kein einfacher Kreis; die relevante Verbindung ist die
topologische Verbindung auf $T^4$ (Kapitel~\ref{ch:verschraenkung}).
Die klassische CTC-Problematik der Gödel-Lösungen betrifft den
dekompaktifizierten Lorentz-Sektor; die kompakte Masse-Zeit-Richtung ist
spektral, nicht kausal geschlossen.

\textbf{Experimentelle Überprüfung.} Die eingerollte Sicht sagt keinen
verborgenen schweren Modenturm oberhalb der Beschleunigerreichweite
voraus. Das diskrete Spektrum der kompakten Zeit \emph{ist} über
$\tilde T\cdot m = 1$ bereits das beobachtete Teilchenspektrum; es gibt
keinen zusätzlichen Kaluza-Klein-Turm. Prüfbar sind: die Teilchenmassen
selbst als $\xi$-Ableitungen (Leptonstruktur, Koide auf $10^{-5}$,
$g-2$); die Modenindex-Erhaltung als Kopplungsauswahlregel; die
CHSH-Amplitude der Ordnung $\xi$; die fraktale Dimension $D_f$ und die
Fundamentallänge $L_0$. Für eine saubere Falsifikation reichen
Verhältnistests allein nicht --- belastbar ist erst eine vorwärts
festgelegte, unabhängig prüfbare Vorhersage.

\section{Die verborgene Inkonsistenz der Standard-QM}

Die Instantanität ist kein QM-spezifisches Problem, sondern ein Erbe der
klassischen Physik. In der Newtonschen Mechanik wirkt Gravitation
instantan; das Coulomb-Potential ist instantan; Lagrange- und
Hamilton-Mechanik beschreiben den Zustand $(q(t),p(t))$ zu einem fixen
$t$ --- die Berechnung ist eine zeitlose algebraische Auswertung. Niemand
fragt, wie lange diese Berechnung dauert, weil Berechnungen in dieser
Sprache keine Dauer haben. Pauli hat nur explizit gemacht, was klassisch
stillschweigend war.

Die QM erbt das und fügt eine Spannung hinzu. \textbf{Der Zustand
$\psi(t)$ ist instantan}: eine zeitlose Größe zu einem fixen $t$, keine
Dauer, kein Prozess. \textbf{Die Auswertung $|\psi(t)|^2$ braucht
Zeit}: Um eine Wahrscheinlichkeitsdichte zu verifizieren, braucht man
viele Messungen, also echte physikalische Zeit. Der Formalismus ist
instantan, seine Verifikation ist zeitlich ausgedehnt. Die Zeit, die
für die probabilistische Auswertung gebraucht wird, steht außerhalb des
Formalismus --- sie ist weder Operator noch Variable, sondern das, was
der Experimentator aufwendet, um den Zustand überhaupt zu kennen.

Superposition $c_1|\psi_1\rangle+c_2|\psi_2\rangle$ ist damit doppelt
zeitlos: als Zustand (kein Zeitoperator, kein intrinsisches
\enquote{Wann}) und als Aussage (probabilistisch, braucht Zeit zur
Verifikation, die außerhalb liegt). Was als überlagert existiert, kann
nur durch wiederholte Messungen über echte Zeit belegt werden --- und
jede Messung zerstört die Superposition.

\section{Drei Konsequenzen, die selten gesagt werden}

In den Lehrbüchern wird Zeit als Parameter eingeführt, ohne zu betonen,
was das Fehlen eines Zeitoperators bedeutet. Es bedeutet konkret:

\textbf{Superposition ist zeitlos.} Ein Elektron im Zustand
$c_\uparrow|\uparrow\rangle+c_\downarrow|\downarrow\rangle$ hat keine
intrinsische Zeitskala. Die Frage \enquote{wie lange ist es in dieser
Superposition?} ist im Formalismus nicht stellbar. Was existiert, ist
die Kohärenzzeit $T_2$ --- eine phänomenologische Größe der
Umgebungskopplung, keine intrinsische Eigenschaft des Zustands. In der
FFGFT trägt jede Mode $\tilde T_i = 1/m_i$ als geometrische Größe. Ein
Zustand \emph{ist} seine Zeitskala --- nicht als Dauer einer
Superposition, sondern als strukturelle Eigenschaft der Mode.

\textbf{Verschränkung ohne Gleichzeitigkeit.} Zwei verschränkte
Teilchen teilen einen Zustand in $\mathcal{H}_1\otimes\mathcal{H}_2$
ohne eingebaute Zeitstruktur. Die Frage nach der Gleichzeitigkeit
einer Messung an beiden ist keine QM-Frage, sondern eine klassische
über den externen Parameter --- und damit eine relativistische. Das ist
der tiefste Grund, warum die QM mit Lorentz-Invarianz vereinbar ist:
Gleichzeitigkeit steht außerhalb des Formalismus. In der FFGFT: Wären
zwei Moden mit $m_1\neq m_2$ verschränkt, hätten sie verschiedene
$\tilde T_1\neq\tilde T_2$. Modenübergreifende Verschränkung erzeugt
inkompatible Zeitskalen --- strukturell problematischer als in der QM,
wo der Konflikt nicht auftritt, weil $t$ außen steht.

\textbf{Kollaps ohne Dauer.} Der Übergang $|\psi\rangle\to|\psi_k\rangle$
hat im Formalismus keine Dauer; er ist instantan in $t$. Wie lange der
Kollaps dauert, liegt außerhalb der Standard-QM. Ohne Zeitoperator gibt
es keine Observable \enquote{Kollapsdauer}. In der FFGFT ist der
Typ-I-Schritt (Messakt) ein geometrischer Übergang: Der Massenkreis
$S^1_m$ wird zur Zeitachse $\mathbb{R}_t$, die Windungszahl wird
verworfen. Das ist ein Informationsverlust mit natürlicher Richtung.

\section{Was $\tilde T\cdot m = 1$ strukturell ändert}

Die Grundrelation setzt Zeit nicht als externen Parameter, sondern als
geometrische Größe jeder Mode: Die Periode $R_m = \hbar/(mc^2)$ ist
intrinsisch --- sie gehört zum Teilchen wie seine Masse. Drei
Konsequenzen:

\textbf{Superposition} innerhalb einer Mode --- intermediäre Torsion
--- hat eine natürliche Zeitskala $\tilde T_i$. Superposition
\emph{über} verschiedene Moden würde inkompatible Zeitskalen
überlagern; das ist der strukturelle Grund, warum modenübergreifende
Superpositionen in der FFGFT problematisch sind.\status{B}

\textbf{Verschränkung} innerhalb einer Mode hat Gleichzeitigkeit durch
$\tilde T$ definiert. Über Moden hinweg fehlt ein gemeinsamer
Zeitbegriff --- ähnlich wie in der QM, aber aus anderem Grund: nicht
weil $t$ extern steht, sondern weil verschiedene Moden verschiedene
geometrische Zeitskalen tragen.

\textbf{Kollaps} (Typ I) hat eine natürliche Richtung ---
verlustbehaftet von $S^1_m$ nach $\mathbb{R}_t$ --- und eine natürliche
Asymmetrie: Die Rückrichtung erfordert die verworfene Windungszahl. Der
Zeitpfeil ist keine Zusatzannahme, sondern Konsequenz der Geometrie des
Messens.

Das Pauli-Theorem gilt für einen offenen Zeitparameter. Für eine
kompakte Zeitdimension greift es nicht: Ihr Konjugiertes hat ein
diskretes Spektrum, die Windungszahlen. $\tilde T\cdot m = 1$
realisiert genau diese Situation: Zeit ist kompakt, das konjugierte
Spektrum ist die diskrete Massenleiter. Das Theorem schließt diesen Fall
nicht aus --- es erzwingt ihn als den einzigen konsistenten Weg, Zeit in
den Hilbertraum zu integrieren.

\section{Ausrollen: Typ I und Typ III}

Zwei verschiedene Operationen tragen denselben Namen. \textbf{Typ-I-
Ausrollen} ist der Messakt: verlustbehaftet, die Projektion vom
kompakten Träger auf die reelle Achse ist nicht umkehrbar. Das ist der
Kollaps. \textbf{Typ-III-Ausrollen} ist der Darstellungswechsel: der
bijektive Pullback aus Kapitel~\ref{ch:hilbert}, kein Informationsverlust,
nur ein Sprachwechsel. Der Korrekturfaktor $K_{\mathrm{frak}}$ sitzt in der
Massenformel des Typ-I-Übergangs; am Typ-III-Pullback ist er nicht
beteiligt. Wer in Hilbertraum-Sprache rechnet, sieht $K_{\mathrm{frak}}$
erst bei der SI-Übersetzung.\status{B}

\section{Zusammenfassung}

\begin{center}
\small\renewcommand{\arraystretch}{1.3}
\begin{tabular}{p{2.2cm}p{3.5cm}p{3.7cm}}
\toprule
 & \textbf{Standard-QM} & \textbf{FFGFT} \\
\midrule
Zeit & absent: kein Operator; externer Parameter & geometrisch: $\tilde T_i = 1/m_i$, kompakt, modengebunden \\
Superposition & zeitlos, keine intrinsische Uhr & intermediäre Torsion in einer Mode; $\tilde T_i$ als Skala \\
Modenübergreifend & algebraisch erlaubt, kein Zeitkonflikt & inkompatible $\tilde T_i\neq\tilde T_j$, strukturell problematisch \\
Verschränkung & zeitlos; Gleichzeitigkeit extern (Lorentz) & modenspezifisches $\tilde T$; kein gemeinsamer Zeitbegriff über Moden \\
Kollaps & instantan in $t$, keine Dauer & Typ I: gerichteter Informationsverlust, Zeitpfeil geometrisch \\
Pauli-Theorem & blockiert Zeitoperator für offenes $t$ & greift nicht für kompaktes $\tilde T$ \\
\bottomrule
\end{tabular}
\end{center}

Die Wahl, wo die Zeit sitzt, geht der Mathematik voraus und entscheidet,
ob die Moden geschenkt oder angeheftet sind. Wer die Sektoren als
räumliche Eigenräume liest, hält die vertraute Physik und zahlt mit
der äußeren Anheftung der Zahlen. Wer die Zeit als geschlossene
Dimension in den Raum nimmt, erhält die Moden strukturell und zahlt mit
einer ontologischen Festlegung. Der ontologische Preis der zweiten Wahl
bleibt auf dem Tisch --- aber der methodische Preis der ersten ist nicht
kleiner, er ist nur vertrauter.

\quelle{Dok.~307 (Zeit im Zustandsraum, Juli 2026), Dok.~334 (Superposition ohne Zeit, Aug.\ 2026), Dok.~333 (Typ I/III), Dok.~306 ($T\cdot E = 1$), Dok.~270 (Entkompaktifizierung)}


\appendix
% QMB_n_symbole.tex — Anhang: Formelzeichen und Abkürzungen
\chapter*{Formelzeichen und Abkürzungen}
\addcontentsline{toc}{chapter}{Formelzeichen und Abkürzungen}

\section*{Grundgrößen der FFGFT}

\begin{description}[leftmargin=4.5em,style=sameline,itemsep=0.4em]
  \item[$\xi$] Geometrieparameter der FFGFT, $\xi = 4/30000 =
    1/7500 \approx 1{,}333\times10^{-4}$. Einziger freier Parameter;
    folgt aus Galois-Gruppenordnungen des Trägers.
  \item[$\tilde T$] Kompakte Zeitkoordinate einer Mode; $\tilde T =
    1/m$ in natürlichen Einheiten ($\hbar = c = 1$).
  \item[$T^4/\mathbb{Z}_3$] FFGFT-Träger: vierdimensionaler Torus mit
    $\mathbb{Z}_3$-Symmetrie. Drei räumliche und eine zeitliche
    Wicklungsrichtung; neun Fixpunkte, Trialitätssymmetrie.
  \item[$K_{\mathrm{frak}}$] Fraktaler Korrekturfaktor,
    $K_{\mathrm{frak}} = 1 - 100\xi \approx 0{,}9867$. Entsteht beim
    Typ-I-Übergang vom kompakten Träger zur SI-Skala; kumulativ über
    17 Größenordnungen (RG-Lauf von Planck- zu elektroschwacher Skala).
  \item[$D_f$] Fraktale Raumdimension, $D_f = 3 - \xi \approx
    2{,}9999$. Lokal; nicht verwechseln mit $K_{\mathrm{frak}}$.
  \item[$L_0$] Fundamentale FFGFT-Länge, $L_0 = \xi\,\ell_P \approx
    2{,}15\times10^{-39}$\,m.
  \item[$\ell_P$] Planck-Länge, $\ell_P \approx 1{,}616\times10^{-35}$\,m.
  \item[$t_P$] Planck-Zeit, $t_P \approx 5{,}39\times10^{-44}$\,s.
\end{description}

\section*{Qubit- und Quantenrechner-Symbole}

\begin{description}[leftmargin=4.5em,style=sameline,itemsep=0.4em]
  \item[$(z,r,\theta)$] Zylindrische Phasenraumkoordinaten eines
    FFGFT-Qubits: $z\in[-1,+1]$ Projektion auf Berechnungsachse,
    $r\in[0,1]$ radiale Superpositionsamplitude, $\theta\in[0,2\pi)$
    relative Phase.
  \item[$(\alpha,\beta)$] Standard-QM-Amplituden eines Qubits:
    $|\psi\rangle = \alpha|0\rangle+\beta|1\rangle$,
    $|\alpha|^2+|\beta|^2 = 1$.
  \item[$K_{\mathrm{frak}}$ (Gatter)] In den $\pi$-Gattern
    ($\sigma_x$-Drehung): Drehwinkel $\alpha = \pi K_{\mathrm{frak}}$
    statt $\pi$. Geometrisch aus $D_f$ abgeleitet; prüfbare Vorhersage.
  \item[$\hat F$] Fundamentaler fraktaler Spektraloperator auf
    $\mathcal{H} = L^2(T^4/\mathbb{Z}_3)\otimes\mathbb{C}^3$.
    Selbstadjungiert; $\xi = \lambda_{\min}(\hat F_{D_4})$.
  \item[$D_4$] Wurzelgitter in $\mathbb{R}^4$: alle ganzzahligen
    Vektoren mit gerader Koordinatensumme. Kusszahl 24; Grundlage des
    Modenspektrums.
  \item[$\mathrm{CHSH}$] Clauser-Horne-Shimony-Holt-Parameter;
    lokal-realistisch $\le2$, Tsirelson-Schranke (QM)
    $2\sqrt2\approx2{,}828$.
  \item[$\Delta\mathrm{CHSH}_{\mathrm{elem}}$] Elementare
    FFGFT-Abweichung pro Messung: $\xi/(2\pi)\approx2\times10^{-5}$.
    Hypothese; formale Herleitung ausstehend.
  \item[$T_1,\,T_2$] Relaxationszeit ($T_1$) und Kohärenzzeit ($T_2$)
    eines Qubits; stets $T_2\le2T_1$.
  \item[$Q,\,Q_\varphi$] Register-Dimension: $Q = 2^n$ (Standard-QFT),
    $Q_\varphi = \varphi^n$ ($\phi$-QFT). Beide äquivalent für
    Periodenfindung bis $N < 2^{20}$ auf $\pm\xi$.
\end{description}

\section*{Spektraltheorie und Gitter}

\begin{description}[leftmargin=4.5em,style=sameline,itemsep=0.4em]
  \item[$N(n)$] Entartung der Schale $|k|^2 = n$ im $D_4$-Gitter;
    Koeffizient der Thetareihe
    $\theta_{D_4} = (\theta_3^4+\theta_4^4)/2$.
  \item[$Z_{\mathbb{Z}^4}(s)$] Spektral-Zeta des kubischen Torus:
    $8(1-4^{1-s})\zeta(s)\zeta(s-1)$.
  \item[$\mathrm{Aut}(D_4)$] Automorphismengruppe des $D_4$-Gitters;
    $|\mathrm{Aut}(D_4)| = 1152 = 2^7\cdot3^2$. Teilerfremd zu 5:
    keine Fünffachsymmetrie.
  \item[$\chi_0,\chi_1,\chi_2$] Die drei $\mathbb{Z}_3$-Charaktere des
    Orbifolds; jede $D_4$-Schale zerfällt gleichmäßig in
    $N(n)/3$ Zustände pro Sektor.
  \item[$\lambda_4$] Bogenlänge auf dem Massenkreis; de-Broglie-Relation
    $\lambda_4\cdot m = 2\pi$ exakt pro Mode (Ortsrelation, nicht
    Zeitrelation).
\end{description}

\section*{Zeitstruktur}

\begin{description}[leftmargin=4.5em,style=sameline,itemsep=0.4em]
  \item[Typ I] Messakt (Ausrollen): verlustbehaftet, $S^1_m\to\mathbb{R}_t$.
    Informationsverlust, Zeitpfeil geometrisch.
  \item[Typ III] Darstellungswechsel (Pullback): verlustfrei, bijektiv.
    Kein Informationsverlust; $K_{\mathrm{frak}}$ nicht beteiligt.
  \item[$\xi_{\mathrm{Higgs}}$] Higgs-Kopplungsparameter,
    $\approx 1{,}038\times10^{-5}$. Wirkt im algebraischen
    Zustandsraum (Kopplungsraten, Dekohärenz); \emph{nicht} identisch
    mit $\xi$, der im geometrischen Ortsraum wirkt.
\end{description}

\section*{Statusmarker}

\begin{description}[leftmargin=2.5em,style=sameline,itemsep=0.3em]
  \item[\textbf{[K]}] Konfirmiert — vollständig aus $\xi$ oder einer
    korpusgesicherten Ableitung bewiesen; Prüfskript bestanden.
  \item[\textbf{[B]}] Belegt — numerisch gesichert; analytischer Beweis
    vorhanden oder trivial.
  \item[\textbf{[S]}] Spekulativ — Hypothese oder Hinweis; noch nicht
    vollständig hergeleitet.
  \item[\textbf{[E]}] Etablierte Physik — Standardergebnis, hier nur
    in Erinnerung gerufen.
  \item[\textbf{[X]}] Geprüft negativ — explizit ausgeschlossen mit
    Begründung.
  \item[\textbf{[SETZUNG]}] Axiomatische Setzung — Ausgangspunkt der
    Ableitung, nicht weiter begründet.
\end{description}

% QMB_n_lit.tex — Anhang: Quellenverzeichnis
\chapter*{Quellenverzeichnis}
\addcontentsline{toc}{chapter}{Quellenverzeichnis}

\section*{FFGFT-Korpusdokumente}

Alle Dokumente sind im Repositorium frei zugänglich.

\medskip\noindent\textbf{PDFs (Korpusdokumente):}\\
\url{https://github.com/jpascher/T0-Time-Mass-Duality/tree/main/2/pdf/}

\medskip\noindent\textbf{Dieses Buch (De/En, LaTeX-Quellen):}\\
\url{https://github.com/jpascher/T0-Time-Mass-Duality/tree/main/2/Sources/qmb-ch/}

\medskip\noindent\textbf{Prüfskripte (Python):}\\
\url{https://github.com/jpascher/T0-Time-Mass-Duality/tree/main/2/python/QMB_Skripte/}

\medskip\noindent Die folgende Liste führt die in diesem Buch verwendeten
Korpusdokumente in Nummernreihenfolge auf.

\begin{description}[leftmargin=3.5em,style=sameline,itemsep=0.3em]\small
\item[Dok.~022] T0 QFT-ML-Addendum: kumulative CHSH-Formel (2025)
\item[Dok.~023] Bell-Tests und $\xi$-modifizierte Korrelationen (Okt. 2025)
\item[Dok.~023a/b] Bell-Tests Teil 2; Antwort zum Bell-Video (2025)
\item[Dok.~034] Geometrischer Qubit-Formalismus (Nov. 2025)
\item[Dok.~055] Dynamische Masse, Photonen, Nichtlokalität (2025)
\item[Dok.~073] Vier Algorithmen im FFGFT-Rahmen (2025)
\item[Dok.~074] No-Go-Theoreme und FFGFT (2025)
\item[Dok.~075/076] RSA-Verfahren; RSA-Periodentests (2026)
\item[Dok.~083] Photonenchip-Einordnung (2025)
\item[Dok.~131] Scheinbare Instantanität — Felddynamik statt Fernwirkung (2026)
\item[Dok.~133] Fraktale Korrektur $K_{\mathrm{frak}}$: Herleitung und RG-Lauf (2026)
\item[Dok.~147] Quantencomputer: $\phi$-QFT und IBM Heron~r2 (Mai 2026)
\item[Dok.~161] Coherent Ising Machine und FFGFT (2026)
\item[Dok.~173] Quantenpunkt als Resonator (März 2026)
\item[Dok.~174] Spin und Qubit-Zustände (2026)
\item[Dok.~175] Qubit-Zustandsräume und Photonikhierarchie (März 2026)
\item[Dok.~176] Shor-Algorithmus: Aufbereitung vs.\ Transformation (2026)
\item[Dok.~178] Spin-Stabilität und Permanentmagnet (2026)
\item[Dok.~183] Google Willow: topologische Fehlerkorrektur (April 2026)
\item[Dok.~184] p-Bit: zwischen Bit und Qubit (2026)
\item[Dok.~190] Korrekturregister — Präzisierungen~11 und 12 (Mai 2026)
\item[Dok.~230] Übersetzbarkeit FFGFT $\leftrightarrow$ Hilbertraum (Mai 2026)
\item[Dok.~231] Hilbertraum-Erweiterung (2026)
\item[Dok.~285] Dimensionsbrücke HLV (2026)
\item[Dok.~291] Dynamischer Ort $\theta = 2/9$ (2026)
\item[Dok.~293] Ikosaeder und Phase $\theta = 2/9$ (2026)
\item[Dok.~306] Native Zeit-Energie-Reziprozität $T\cdot E=1$ (2026)
\item[Dok.~307] Zeit im Zustandsraum: eingerollt vs.\ ausgerollt (Juli 2026)
\item[Dok.~311] Vier auf Drei: $24 = 12+12$ (2026)
\item[Dok.~313] Kein Anfang: Schließungsgabelung (2026)
\item[Dok.~314] Das Gitter im Hilbertraum (2026)
\item[Dok.~316] Weyl-Gesetz und Riemann-Hypothese (Aug. 2026)
\item[Dok.~322] Spektraltheorie auf $T^4/\mathbb{Z}_3$ (2026)
\item[Dok.~327] Selbstadjungiertheit von $\hat F$ (Aug. 2026)
\item[Dok.~333] $K_{\mathrm{frak}}$, Ausrollen und $\tilde T\cdot m$-Dualität (Aug. 2026)
\item[Dok.~334] Superposition ohne Zeit (Aug. 2026)
\item[Dok.~338] Masseberechnungen und Galois-Gegenüberstellung (2026)
\item[Dok.~339] Frobenius-Trennung massiv/masselos in $GF(27)^*$ (Aug. 2026)
\item[Dok.~343] Zeta-Funktion, Auflösungsgrenze, Quantenrechner (Sept. 2026)
\item[Dok.~353] Koide-Amplitude $d/c=\sqrt2$ aus $T^4/\mathbb{Z}_3$-Geometrie (Sept. 2026)
\item[Dok.~364] Ikosaeder/D4-Vergleich, Galois-Struktur (Sept. 2026)
\item[Dok.~365] Fundament der FFGFT (Sept. 2026)
\end{description}

\section*{Ausgewählte externe Quellen}

\begin{description}[leftmargin=3.5em,style=sameline,itemsep=0.4em]\small

\item[Bell 1964] J.~S. Bell: \textit{On the Einstein-Podolsky-Rosen
  paradox.} Physics 1 (3), 195–200 (1964). — Das Originalargument
  für die Unmöglichkeit lokal-realistischer verborgener Variablen.

\item[CHSH 1969] J.~F. Clauser, M.~A. Horne, A.~Shimony, R.~A. Holt:
  \textit{Proposed experiment to test local hidden-variable theories.}
  Phys. Rev. Lett. 23, 880 (1969). — Experimentell prüfbare Form der
  Bell-Ungleichung.

\item[Aspect 1982] A.~Aspect, P.~Grangier, G.~Roger: \textit{Experimental
  realization of Einstein-Podolsky-Rosen-Bohm Gedankenexperiment.}
  Phys. Rev. Lett. 49, 91 (1982). — Erster schlupflochfreier
  Bell-Test (Loophole-Schlupflöcher 2015 vollständig geschlossen).

\item[Kochen/Specker 1967] S.~Kochen, E.~P. Specker: \textit{The
  problem of hidden variables in quantum mechanics.} J. Math. Mech.
  17, 59–87 (1967). — Beweis der Kontextabhängigkeit von Observablen.

\item[PBR 2012] M.~F. Pusey, J.~Barrett, T.~Rudolph: \textit{On the
  reality of the quantum state.} Nature Physics 8, 475–478 (2012). —
  Zeigt: Quantenzustand ist ontisch, kein Wissenszustand.

\item[Shor 1994] P.~W. Shor: \textit{Algorithms for quantum computation:
  discrete logarithms and factoring.} Proc. 35th FOCS, 124–134 (1994). —
  Polynomial-Zeit-Faktorisierung auf einem Quantencomputer.

\item[Grover 1996] L.~K. Grover: \textit{A fast quantum mechanical
  algorithm for database search.} Proc. 28th STOC, 212–219 (1996). —
  $O(\sqrt N)$-Suche in einer unsortierten Datenbank.

\item[Loss/DiVincenzo 1998] D.~Loss, D.~P. DiVincenzo: \textit{Quantum
  computation with quantum dots.} Phys. Rev. A 57, 120 (1998). —
  Spin-Qubit-Gatter durch Austauschwechselwirkung.

\item[Willow 2024] Google Quantum AI: \textit{Quantum error correction
  below the surface code threshold.} Nature 634, 2024. — Erstmaliges
  Unterschreiten des Fehlerkorrektur-Schwellenwerts auf Willow
  (105 Qubits, IBM-Heron-Vergleich).

\item[Will 2025] Will et al.: \textit{Floquet topological order in a
  quantum processor.} Nature, Sept. 2025. — Nicht-gleichgewichtliche
  topologische Ordnung (chirale Kantenmoden, Anyonen) auf Willow.

\item[Solovay/Kitaev] A.~Y. Kitaev: \textit{Quantum computations:
  algorithms and error correction.} Russ. Math. Surv. 52 (1997). —
  Universalität von $\{H,T,\mathrm{CNOT}\}$ für die Quantenrechnung.

\item[Tsirelson 1980] B.~S. Tsirelson: \textit{Quantum generalizations
  of Bell's inequality.} Lett. Math. Phys. 4, 93–100 (1980). —
  Obere Schranke der Quantenkorrelationen: $2\sqrt2$.

\item[Kaluza 1921 / Klein 1926] T.~Kaluza (1921), O.~Klein (1926):
  Kompaktifizierte Dimension erzeugt diskreten Modenturm und
  Ladungsquantisierung. Mathematische Blaupause für den
  Zeit-Masse-Mechanismus der FFGFT.

\item[Page/Wootters 1983] D.~N. Page, W.~K. Wootters: \textit{Evolution
  without evolution.} Phys. Rev. D 27, 2885 (1983). — Zeit als
  Verschränkungskorrelation innerhalb des Zustandsraums.

\item[Weyl 1912] H.~Weyl: Eigenwertgesetz kompakter Mannigfaltigkeiten
  $N(\lambda)\sim C\lambda^{d/2}$. Begründet die Unmöglichkeit,
  Riemann-Nullstellen als Laplace-Spektrum zu tragen.

\item[Inagaki 2016 / McMahon 2016] T.~Inagaki et al.; P.~L. McMahon
  et al.: Coherent Ising Machine (CIM) mit DOPO-Netzwerk. Science
  2016. — Experimentelle Realisierung physikalischer Optimierung
  durch Relaxation.

\item[Schiemann 1997] A.~Schiemann: Isospektrale, nicht-isometrische
  Tori in Dimension 4. Zeigt, dass das Spektrum allein das Gitter
  nicht bestimmt. Grenze der Aussagekraft spektraler Methoden.

\item[BSI 2023] Bundesamt für Sicherheit in der Informationstechnik:
  \textit{Status of quantum computer development.} 2023. —
  Aktuelle Einschätzung zur Bedrohung kryptographischer Verfahren.

\end{description}

% QMB_n_ai.tex — Anhang: Hinweis zur KI-Unterstützung
\chapter*{Hinweis zur Entstehung}
\addcontentsline{toc}{chapter}{Hinweis zur Entstehung}

Dieses Buch entstand aus dem FFGFT-Korpus, der seit 2025 über 365
Einzeldokumente umfasst. Die Texte der einzelnen Kapitel wurden von
Johann Pascher konzipiert und inhaltlich verantwortet; Claude
(Anthropic, Sonnet~4.6) hat beim Ausformulieren, beim Zusammenführen
der Quellsektionen und beim Kompilieren der \LaTeX-Vorlage assistiert.

\medskip\noindent\textbf{Was KI-Assistenz in diesem Projekt bedeutete:}
Die Quelldokumente (Dok.~022 bis Dok.~365) enthalten die eigentlichen
Ableitungen, Prüfskripte und epistemischen Marker. Die KI las diese
Dokumente ab, formulierte die Kapitel in eigener Sprache um --- ohne
wörtliche Übernahme, ohne fremde Inhalte --- und prüfte die Konsistenz
der Statusmarker. Alle inhaltlichen Entscheidungen, Grenzziehungen und
Statusmarker sind Johanns Werk; die KI war ein Formulierungswerkzeug,
kein inhaltlicher Autor.

\medskip\noindent\textbf{Was KI-Assistenz nicht bedeutete:}
Keine der physikalischen Aussagen stammt aus dem Trainingskorpus der KI.
Alle Zahlen, Ableitungen und Querverweise gehen auf die genannten
Quelldokumente zurück, die im Repositorium öffentlich einsehbar sind.
Anfragen, ob ein Ergebnis \enquote{neu} oder \enquote{korrekt} sei,
wurden vom Autor selbst beantwortet, nicht von der KI.

\medskip\noindent\textbf{Prüfbarkeit:} Jedes Kapitel trägt am Ende
einen \texttt{\\quelle}-Verweis auf die zugrundeliegenden
Korpusdokumente. Jede inhaltliche Aussage trägt einen Statusmarker
([K], [B], [S], [E], [X]). Wer eine Aussage hinterfragen will,
findet im genannten Dokument den vollständigen Herleitungsweg mit
Prüfskript.


\end{document}
